% Môn: Vật lý
% Lớp: 12
% Chương: Chương III. Từ trường
% Bài: L12C3 Bài 19. Điện từ trường. Mô hình sóng điện từ
% Số câu: 324
% App đọc file này trực tiếp — sửa rồi Commit trên GitHub, bấm Đồng bộ GitHub .tex

% L12C3 Bài 19. Điện từ trường. Mô hình sóng điện từ

% ===== Câu 1 =====
% ID: L12C3B19-01-DS
% Mức: TH
\dangbt{Giải thích truyền thông tin bằng sóng điện từ và các tình huống thực tiễn}
\begin{ex}
% ID: L12C3B19-01-DS
% Mức: TH
Xét Giải thích truyền thông tin bằng sóng điện từ và các tình huống thực tiễn (Tình huống 1). Hãy xác định tính đúng, sai của bốn phát biểu sau.
\choiceTF
{\True Thông tin được ghép lên sóng mang bằng biến điệu, truyền đi và được máy thu giải điều chế.}
{Máy thu nhận thông tin mà không cần chọn sóng hay giải điều chế.}
{\True Khi giải cần xác định đúng dữ kiện, phương chiều, điều kiện áp dụng và đơn vị.}
{Kết quả không phụ thuộc dữ kiện và có thể bỏ qua điều kiện.}
\loigiai{
\begin{itemchoice}
\item (Đúng) Thông tin được ghép lên sóng mang bằng biến điệu, truyền đi và được máy thu giải điều chế. (Vì): Thông tin được ghép lên sóng mang bằng biến điệu, truyền đi và được máy thu giải điều chế. b) Sai: Máy thu nhận thông tin mà không cần chọn sóng hay giải điều chế. c) Đúng: Khi giải cần xác định đúng dữ kiện, phương chiều, điều kiện áp dụng và đơn vị. d) Sai: Kết quả không phụ thuộc dữ kiện và có thể bỏ qua điều kiện. Kiến thức cốt lõi: Thông tin được ghép lên sóng mang bằng biến điệu, truyền đi và được máy thu giải điều chế
\item (Sai) Máy thu nhận thông tin mà không cần chọn sóng hay giải điều chế.
\item (Đúng) Khi giải cần xác định đúng dữ kiện, phương chiều, điều kiện áp dụng và đơn vị.
\item (Sai) Kết quả không phụ thuộc dữ kiện và có thể bỏ qua điều kiện.
\end{itemchoice}
}
\end{ex}

% ===== Câu 2 =====
% ID: L12C3B19-01-TL
% Mức: TH
\begin{bt}
% ID: L12C3B19-01-TL
% Mức: TH
Trình bày kiến thức cốt lõi của Giải thích truyền thông tin bằng sóng điện từ và các tình huống thực tiễn và nêu điều kiện áp dụng.
\loigiai{
Cần nêu được: Thông tin được ghép lên sóng mang bằng biến điệu, truyền đi và được máy thu giải điều chế. Khi vận dụng phải xác định đúng dữ kiện, phương chiều nếu có, điều kiện áp dụng, hệ đơn vị và kiểm tra tính hợp lí. Nhận định “Máy thu nhận thông tin mà không cần chọn sóng hay giải điều chế.” sai.
}
\end{bt}

% ===== Câu 3 =====
% ID: L12C3B19-01-TLN
% Mức: TH
\begin{ex}
% ID: L12C3B19-01-TLN
% Mức: TH
Trong bốn phát biểu về Giải thích truyền thông tin bằng sóng điện từ và các tình huống thực tiễn, có bao nhiêu phát biểu đúng? (Chỉ ghi một số nguyên.) (1) Thông tin được ghép lên sóng mang bằng biến điệu, truyền đi và được máy thu giải điều chế. (2) Máy thu nhận thông tin mà không cần chọn sóng hay giải điều chế. (3) Cần kiểm tra điều kiện áp dụng. (4) Có thể bỏ qua đơn vị.
\shortans{2}
\loigiai{
Đối chiếu với kiến thức: Thông tin được ghép lên sóng mang bằng biến điệu, truyền đi và được máy thu giải điều chế. Có 2 phát biểu đúng.
}
\end{ex}

% ===== Câu 4 =====
% ID: L12C3B19-01-TN
% Mức: NB
\begin{ex}
% ID: L12C3B19-01-TN
% Mức: NB
Phát biểu nào sau đây đúng về Giải thích truyền thông tin bằng sóng điện từ và các tình huống thực tiễn?
\choice
{Máy thu nhận thông tin mà không cần chọn sóng hay giải điều chế.}
{Có thể bỏ qua phương, chiều và đơn vị trong mọi trường hợp.}
{Kết quả không phụ thuộc dữ kiện và điều kiện áp dụng.}
{\True Thông tin được ghép lên sóng mang bằng biến điệu, truyền đi và được máy thu giải điều chế.}
\loigiai{
Thông tin được ghép lên sóng mang bằng biến điệu, truyền đi và được máy thu giải điều chế. Vì vậy phương án $D$ đúng; các phương án còn lại sai.
}
\end{ex}

% ===== Câu 5 =====
% ID: L12C3B19-02-DS
% Mức: TH
\begin{ex}
% ID: L12C3B19-02-DS
% Mức: TH
Xét Giải thích truyền thông tin bằng sóng điện từ và các tình huống thực tiễn (Tình huống 2). Hãy xác định tính đúng, sai của bốn phát biểu sau.
\choiceTF
{Máy thu nhận thông tin mà không cần chọn sóng hay giải điều chế.}
{\True Thông tin được ghép lên sóng mang bằng biến điệu, truyền đi và được máy thu giải điều chế.}
{Kết quả không phụ thuộc dữ kiện và có thể bỏ qua điều kiện.}
{\True Khi giải cần xác định đúng dữ kiện, phương chiều, điều kiện áp dụng và đơn vị.}
\loigiai{
\begin{itemchoice}
\item (Sai) Máy thu nhận thông tin mà không cần chọn sóng hay giải điều chế. (Vì): Máy thu nhận thông tin mà không cần chọn sóng hay giải điều chế. b) Đúng: Thông tin được ghép lên sóng mang bằng biến điệu, truyền đi và được máy thu giải điều chế. c) Sai: Kết quả không phụ thuộc dữ kiện và có thể bỏ qua điều kiện. d) Đúng: Khi giải cần xác định đúng dữ kiện, phương chiều, điều kiện áp dụng và đơn vị. Kiến thức cốt lõi: Thông tin được ghép lên sóng mang bằng biến điệu, truyền đi và được máy thu giải điều chế
\item (Đúng) Thông tin được ghép lên sóng mang bằng biến điệu, truyền đi và được máy thu giải điều chế.
\item (Sai) Kết quả không phụ thuộc dữ kiện và có thể bỏ qua điều kiện.
\item (Đúng) Khi giải cần xác định đúng dữ kiện, phương chiều, điều kiện áp dụng và đơn vị.
\end{itemchoice}
}
\end{ex}

% ===== Câu 6 =====
% ID: L12C3B19-02-TL
% Mức: TH
\begin{bt}
% ID: L12C3B19-02-TL
% Mức: TH
Giải thích vì sao nhận định sau đúng: “Thông tin được ghép lên sóng mang bằng biến điệu, truyền đi và được máy thu giải điều chế.”
\loigiai{
Cần nêu được: Thông tin được ghép lên sóng mang bằng biến điệu, truyền đi và được máy thu giải điều chế. Khi vận dụng phải xác định đúng dữ kiện, phương chiều nếu có, điều kiện áp dụng, hệ đơn vị và kiểm tra tính hợp lí. Nhận định “Máy thu nhận thông tin mà không cần chọn sóng hay giải điều chế.” sai.
}
\end{bt}

% ===== Câu 7 =====
% ID: L12C3B19-02-TLN
% Mức: TH
\begin{ex}
% ID: L12C3B19-02-TLN
% Mức: TH
Trong bốn phát biểu về Giải thích truyền thông tin bằng sóng điện từ và các tình huống thực tiễn, có bao nhiêu phát biểu đúng? (Chỉ ghi một số nguyên.) (1) Máy thu nhận thông tin mà không cần chọn sóng hay giải điều chế. (2) Thông tin được ghép lên sóng mang bằng biến điệu, truyền đi và được máy thu giải điều chế. (3) Kết quả phải phù hợp ý nghĩa vật lí. (4) Mọi đại lượng đều là vô hướng.
\shortans{1}
\loigiai{
Đối chiếu với kiến thức: Thông tin được ghép lên sóng mang bằng biến điệu, truyền đi và được máy thu giải điều chế. Có 1 phát biểu đúng.
}
\end{ex}

% ===== Câu 8 =====
% ID: L12C3B19-02-TN
% Mức: NB
\begin{ex}
% ID: L12C3B19-02-TN
% Mức: NB
Trong nội dung Giải thích truyền thông tin bằng sóng điện từ và các tình huống thực tiễn?
\choice
{\True Thông tin được ghép lên sóng mang bằng biến điệu, truyền đi và được máy thu giải điều chế.}
{Máy thu nhận thông tin mà không cần chọn sóng hay giải điều chế.}
{Có thể bỏ qua phương, chiều và đơn vị trong mọi trường hợp.}
{Kết quả không phụ thuộc dữ kiện và điều kiện áp dụng.}
\loigiai{
Thông tin được ghép lên sóng mang bằng biến điệu, truyền đi và được máy thu giải điều chế. Vì vậy phương án $A$ đúng; các phương án còn lại sai.
}
\end{ex}

% ===== Câu 9 =====
% ID: L12C3B19-03-DS
% Mức: VD
\begin{ex}
% ID: L12C3B19-03-DS
% Mức: VD
Xét Giải thích truyền thông tin bằng sóng điện từ và các tình huống thực tiễn (Tình huống 3). Hãy xác định tính đúng, sai của bốn phát biểu sau.
\choiceTF
{\True Khi giải cần xác định đúng dữ kiện, phương chiều, điều kiện áp dụng và đơn vị.}
{Kết quả không phụ thuộc dữ kiện và có thể bỏ qua điều kiện.}
{\True Thông tin được ghép lên sóng mang bằng biến điệu, truyền đi và được máy thu giải điều chế.}
{Máy thu nhận thông tin mà không cần chọn sóng hay giải điều chế.}
\loigiai{
\begin{itemchoice}
\item (Đúng) Khi giải cần xác định đúng dữ kiện, phương chiều, điều kiện áp dụng và đơn vị. (Vì): Khi giải cần xác định đúng dữ kiện, phương chiều, điều kiện áp dụng và đơn vị. b) Sai: Kết quả không phụ thuộc dữ kiện và có thể bỏ qua điều kiện. c) Đúng: Thông tin được ghép lên sóng mang bằng biến điệu, truyền đi và được máy thu giải điều chế. d) Sai: Máy thu nhận thông tin mà không cần chọn sóng hay giải điều chế. Kiến thức cốt lõi: Thông tin được ghép lên sóng mang bằng biến điệu, truyền đi và được máy thu giải điều chế
\item (Sai) Kết quả không phụ thuộc dữ kiện và có thể bỏ qua điều kiện.
\item (Đúng) Thông tin được ghép lên sóng mang bằng biến điệu, truyền đi và được máy thu giải điều chế.
\item (Sai) Máy thu nhận thông tin mà không cần chọn sóng hay giải điều chế.
\end{itemchoice}
}
\end{ex}

% ===== Câu 10 =====
% ID: L12C3B19-03-TL
% Mức: VD
\begin{bt}
% ID: L12C3B19-03-TL
% Mức: VD
Phân tích sai lầm trong nhận định: “Máy thu nhận thông tin mà không cần chọn sóng hay giải điều chế.”
\loigiai{
Cần nêu được: Thông tin được ghép lên sóng mang bằng biến điệu, truyền đi và được máy thu giải điều chế. Khi vận dụng phải xác định đúng dữ kiện, phương chiều nếu có, điều kiện áp dụng, hệ đơn vị và kiểm tra tính hợp lí. Nhận định “Máy thu nhận thông tin mà không cần chọn sóng hay giải điều chế.” sai.
}
\end{bt}

% ===== Câu 11 =====
% ID: L12C3B19-03-TLN
% Mức: TH
\begin{ex}
% ID: L12C3B19-03-TLN
% Mức: TH
Trong bốn phát biểu về Giải thích truyền thông tin bằng sóng điện từ và các tình huống thực tiễn, có bao nhiêu phát biểu đúng? (Chỉ ghi một số nguyên.) (1) Thông tin được ghép lên sóng mang bằng biến điệu, truyền đi và được máy thu giải điều chế. (2) Máy thu nhận thông tin mà không cần chọn sóng hay giải điều chế. (3) Cần kiểm tra điều kiện áp dụng. (4) Có thể bỏ qua đơn vị.
\shortans{2}
\loigiai{
Đối chiếu với kiến thức: Thông tin được ghép lên sóng mang bằng biến điệu, truyền đi và được máy thu giải điều chế. Có 2 phát biểu đúng.
}
\end{ex}

% ===== Câu 12 =====
% ID: L12C3B19-03-TN
% Mức: NB
\begin{ex}
% ID: L12C3B19-03-TN
% Mức: NB
Tình huống 3: chọn kết luận chính xác về Giải thích truyền thông tin bằng sóng điện từ và các tình huống thực tiễn?
\choice
{Kết quả không phụ thuộc dữ kiện và điều kiện áp dụng.}
{\True Thông tin được ghép lên sóng mang bằng biến điệu, truyền đi và được máy thu giải điều chế.}
{Máy thu nhận thông tin mà không cần chọn sóng hay giải điều chế.}
{Có thể bỏ qua phương, chiều và đơn vị trong mọi trường hợp.}
\loigiai{
Thông tin được ghép lên sóng mang bằng biến điệu, truyền đi và được máy thu giải điều chế. Vì vậy phương án $B$ đúng; các phương án còn lại sai.
}
\end{ex}

% ===== Câu 13 =====
% ID: L12C3B19-04-DS
% Mức: VD
\begin{ex}
% ID: L12C3B19-04-DS
% Mức: VD
Xét Giải thích truyền thông tin bằng sóng điện từ và các tình huống thực tiễn (Tình huống 4). Hãy xác định tính đúng, sai của bốn phát biểu sau.
\choiceTF
{Kết quả không phụ thuộc dữ kiện và có thể bỏ qua điều kiện.}
{\True Khi giải cần xác định đúng dữ kiện, phương chiều, điều kiện áp dụng và đơn vị.}
{Máy thu nhận thông tin mà không cần chọn sóng hay giải điều chế.}
{\True Thông tin được ghép lên sóng mang bằng biến điệu, truyền đi và được máy thu giải điều chế.}
\loigiai{
\begin{itemchoice}
\item (Sai) Kết quả không phụ thuộc dữ kiện và có thể bỏ qua điều kiện. (Vì): Kết quả không phụ thuộc dữ kiện và có thể bỏ qua điều kiện. b) Đúng: Khi giải cần xác định đúng dữ kiện, phương chiều, điều kiện áp dụng và đơn vị. c) Sai: Máy thu nhận thông tin mà không cần chọn sóng hay giải điều chế. d) Đúng: Thông tin được ghép lên sóng mang bằng biến điệu, truyền đi và được máy thu giải điều chế. Kiến thức cốt lõi: Thông tin được ghép lên sóng mang bằng biến điệu, truyền đi và được máy thu giải điều chế
\item (Đúng) Khi giải cần xác định đúng dữ kiện, phương chiều, điều kiện áp dụng và đơn vị.
\item (Sai) Máy thu nhận thông tin mà không cần chọn sóng hay giải điều chế.
\item (Đúng) Thông tin được ghép lên sóng mang bằng biến điệu, truyền đi và được máy thu giải điều chế.
\end{itemchoice}
}
\end{ex}

% ===== Câu 14 =====
% ID: L12C3B19-04-TL
% Mức: VD
\begin{bt}
% ID: L12C3B19-04-TL
% Mức: VD
Đề xuất các bước giải một bài toán thuộc dạng Giải thích truyền thông tin bằng sóng điện từ và các tình huống thực tiễn.
\loigiai{
Cần nêu được: Thông tin được ghép lên sóng mang bằng biến điệu, truyền đi và được máy thu giải điều chế. Khi vận dụng phải xác định đúng dữ kiện, phương chiều nếu có, điều kiện áp dụng, hệ đơn vị và kiểm tra tính hợp lí. Nhận định “Máy thu nhận thông tin mà không cần chọn sóng hay giải điều chế.” sai.
}
\end{bt}

% ===== Câu 15 =====
% ID: L12C3B19-04-TLN
% Mức: VD
\begin{ex}
% ID: L12C3B19-04-TLN
% Mức: VD
Trong bốn phát biểu về Giải thích truyền thông tin bằng sóng điện từ và các tình huống thực tiễn, có bao nhiêu phát biểu đúng? (Chỉ ghi một số nguyên.) (1) Máy thu nhận thông tin mà không cần chọn sóng hay giải điều chế. (2) Thông tin được ghép lên sóng mang bằng biến điệu, truyền đi và được máy thu giải điều chế. (3) Kết quả phải phù hợp ý nghĩa vật lí. (4) Mọi đại lượng đều là vô hướng.
\shortans{2}
\loigiai{
Đối chiếu với kiến thức: Thông tin được ghép lên sóng mang bằng biến điệu, truyền đi và được máy thu giải điều chế. Có 2 phát biểu đúng.
}
\end{ex}

% ===== Câu 16 =====
% ID: L12C3B19-04-TN
% Mức: TH
\begin{ex}
% ID: L12C3B19-04-TN
% Mức: TH
Nội dung nào mô tả đúng nhất Giải thích truyền thông tin bằng sóng điện từ và các tình huống thực tiễn?
\choice
{Có thể bỏ qua phương, chiều và đơn vị trong mọi trường hợp.}
{Kết quả không phụ thuộc dữ kiện và điều kiện áp dụng.}
{\True Thông tin được ghép lên sóng mang bằng biến điệu, truyền đi và được máy thu giải điều chế.}
{Máy thu nhận thông tin mà không cần chọn sóng hay giải điều chế.}
\loigiai{
Thông tin được ghép lên sóng mang bằng biến điệu, truyền đi và được máy thu giải điều chế. Vì vậy phương án $C$ đúng; các phương án còn lại sai.
}
\end{ex}

% ===== Câu 17 =====
% ID: L12C3B19-05-DS
% Mức: VD
\begin{ex}
% ID: L12C3B19-05-DS
% Mức: VD
Xét Giải thích truyền thông tin bằng sóng điện từ và các tình huống thực tiễn (Tình huống 5). Hãy xác định tính đúng, sai của bốn phát biểu sau.
\choiceTF
{\True Thông tin được ghép lên sóng mang bằng biến điệu, truyền đi và được máy thu giải điều chế.}
{Kết quả không phụ thuộc dữ kiện và có thể bỏ qua điều kiện.}
{Máy thu nhận thông tin mà không cần chọn sóng hay giải điều chế.}
{\True Khi giải cần xác định đúng dữ kiện, phương chiều, điều kiện áp dụng và đơn vị.}
\loigiai{
\begin{itemchoice}
\item (Đúng) Thông tin được ghép lên sóng mang bằng biến điệu, truyền đi và được máy thu giải điều chế. (Vì): Thông tin được ghép lên sóng mang bằng biến điệu, truyền đi và được máy thu giải điều chế. b) Sai: Kết quả không phụ thuộc dữ kiện và có thể bỏ qua điều kiện. c) Sai: Máy thu nhận thông tin mà không cần chọn sóng hay giải điều chế. d) Đúng: Khi giải cần xác định đúng dữ kiện, phương chiều, điều kiện áp dụng và đơn vị. Kiến thức cốt lõi: Thông tin được ghép lên sóng mang bằng biến điệu, truyền đi và được máy thu giải điều chế
\item (Sai) Kết quả không phụ thuộc dữ kiện và có thể bỏ qua điều kiện.
\item (Sai) Máy thu nhận thông tin mà không cần chọn sóng hay giải điều chế.
\item (Đúng) Khi giải cần xác định đúng dữ kiện, phương chiều, điều kiện áp dụng và đơn vị.
\end{itemchoice}
}
\end{ex}

% ===== Câu 18 =====
% ID: L12C3B19-05-TL
% Mức: VD
\begin{bt}
% ID: L12C3B19-05-TL
% Mức: VD
Nêu một tình huống thực tiễn liên quan đến Giải thích truyền thông tin bằng sóng điện từ và các tình huống thực tiễn và giải thích bằng kiến thức Vật lí.
\loigiai{
Cần nêu được: Thông tin được ghép lên sóng mang bằng biến điệu, truyền đi và được máy thu giải điều chế. Khi vận dụng phải xác định đúng dữ kiện, phương chiều nếu có, điều kiện áp dụng, hệ đơn vị và kiểm tra tính hợp lí. Nhận định “Máy thu nhận thông tin mà không cần chọn sóng hay giải điều chế.” sai.
}
\end{bt}

% ===== Câu 19 =====
% ID: L12C3B19-05-TLN
% Mức: VD
\begin{ex}
% ID: L12C3B19-05-TLN
% Mức: VD
Trong bốn phát biểu về Giải thích truyền thông tin bằng sóng điện từ và các tình huống thực tiễn, có bao nhiêu phát biểu đúng? (Chỉ ghi một số nguyên.) (1) Thông tin được ghép lên sóng mang bằng biến điệu, truyền đi và được máy thu giải điều chế. (2) Máy thu nhận thông tin mà không cần chọn sóng hay giải điều chế. (3) Cần kiểm tra điều kiện áp dụng. (4) Có thể bỏ qua đơn vị.
\shortans{2}
\loigiai{
Đối chiếu với kiến thức: Thông tin được ghép lên sóng mang bằng biến điệu, truyền đi và được máy thu giải điều chế. Có 2 phát biểu đúng.
}
\end{ex}

% ===== Câu 20 =====
% ID: L12C3B19-05-TN
% Mức: TH
\begin{ex}
% ID: L12C3B19-05-TN
% Mức: TH
Khi phân tích Giải thích truyền thông tin bằng sóng điện từ và các tình huống thực tiễn?
\choice
{Máy thu nhận thông tin mà không cần chọn sóng hay giải điều chế.}
{Có thể bỏ qua phương, chiều và đơn vị trong mọi trường hợp.}
{Kết quả không phụ thuộc dữ kiện và điều kiện áp dụng.}
{\True Thông tin được ghép lên sóng mang bằng biến điệu, truyền đi và được máy thu giải điều chế.}
\loigiai{
Thông tin được ghép lên sóng mang bằng biến điệu, truyền đi và được máy thu giải điều chế. Vì vậy phương án $D$ đúng; các phương án còn lại sai.
}
\end{ex}

% ===== Câu 21 =====
% ID: L12C3B19-06-DS
% Mức: TH
\dangbt{Mô tả sự hình thành và lan truyền của sóng điện từ}
\begin{ex}
% ID: L12C3B19-06-DS
% Mức: TH
Xét Mô tả sự hình thành và lan truyền của sóng điện từ (Tình huống 1). Hãy xác định tính đúng, sai của bốn phát biểu sau.
\choiceTF
{\True Sóng điện từ là điện từ trường lan truyền trong không gian và truyền được trong chân không.}
{Sóng điện từ bắt buộc cần môi trường vật chất để truyền.}
{\True Khi giải cần xác định đúng dữ kiện, phương chiều, điều kiện áp dụng và đơn vị.}
{Kết quả không phụ thuộc dữ kiện và có thể bỏ qua điều kiện.}
\loigiai{
\begin{itemchoice}
\item (Đúng) Sóng điện từ là điện từ trường lan truyền trong không gian và truyền được trong chân không. (Vì): Sóng điện từ là điện từ trường lan truyền trong không gian và truyền được trong chân không. b) Sai: Sóng điện từ bắt buộc cần môi trường vật chất để truyền. c) Đúng: Khi giải cần xác định đúng dữ kiện, phương chiều, điều kiện áp dụng và đơn vị. d) Sai: Kết quả không phụ thuộc dữ kiện và có thể bỏ qua điều kiện. Kiến thức cốt lõi: Sóng điện từ là điện từ trường lan truyền trong không gian và truyền được trong chân không
\item (Sai) Sóng điện từ bắt buộc cần môi trường vật chất để truyền.
\item (Đúng) Khi giải cần xác định đúng dữ kiện, phương chiều, điều kiện áp dụng và đơn vị.
\item (Sai) Kết quả không phụ thuộc dữ kiện và có thể bỏ qua điều kiện.
\end{itemchoice}
}
\end{ex}

% ===== Câu 22 =====
% ID: L12C3B19-06-TL
% Mức: VDC
\dangbt{Giải thích truyền thông tin bằng sóng điện từ và các tình huống thực tiễn}
\begin{bt}
% ID: L12C3B19-06-TL
% Mức: VDC
So sánh nhận định đúng “Thông tin được ghép lên sóng mang bằng biến điệu, truyền đi và được máy thu giải điều chế.” với nhận định sai “Máy thu nhận thông tin mà không cần chọn sóng hay giải điều chế.”; chỉ rõ căn cứ Vật lí.
\loigiai{
Cần nêu được: Thông tin được ghép lên sóng mang bằng biến điệu, truyền đi và được máy thu giải điều chế. Khi vận dụng phải xác định đúng dữ kiện, phương chiều nếu có, điều kiện áp dụng, hệ đơn vị và kiểm tra tính hợp lí. Nhận định “Máy thu nhận thông tin mà không cần chọn sóng hay giải điều chế.” sai.
}
\end{bt}

% ===== Câu 23 =====
% ID: L12C3B19-06-TLN
% Mức: VDC
\begin{ex}
% ID: L12C3B19-06-TLN
% Mức: VDC
Trong bốn phát biểu về Giải thích truyền thông tin bằng sóng điện từ và các tình huống thực tiễn, có bao nhiêu phát biểu đúng? (Chỉ ghi một số nguyên.) (1) Khi giải cần xác định đúng dữ kiện, phương chiều, điều kiện áp dụng và đơn vị. (2) Máy thu nhận thông tin mà không cần chọn sóng hay giải điều chế. (3) Phải dùng đúng định luật cho tình huống. (4) Thông tin được ghép lên sóng mang bằng biến điệu, truyền đi và được máy thu giải điều chế.
\shortans{3}
\loigiai{
Đối chiếu với kiến thức: Thông tin được ghép lên sóng mang bằng biến điệu, truyền đi và được máy thu giải điều chế. Có 3 phát biểu đúng.
}
\end{ex}

% ===== Câu 24 =====
% ID: L12C3B19-06-TN
% Mức: TH
\begin{ex}
% ID: L12C3B19-06-TN
% Mức: TH
Một học sinh ôn tập Giải thích truyền thông tin bằng sóng điện từ và các tình huống thực tiễn?
\choice
{\True Thông tin được ghép lên sóng mang bằng biến điệu, truyền đi và được máy thu giải điều chế.}
{Máy thu nhận thông tin mà không cần chọn sóng hay giải điều chế.}
{Có thể bỏ qua phương, chiều và đơn vị trong mọi trường hợp.}
{Kết quả không phụ thuộc dữ kiện và điều kiện áp dụng.}
\loigiai{
Thông tin được ghép lên sóng mang bằng biến điệu, truyền đi và được máy thu giải điều chế. Vì vậy phương án $A$ đúng; các phương án còn lại sai.
}
\end{ex}

% ===== Câu 25 =====
% ID: L12C3B19-07-DS
% Mức: TH
\dangbt{Mô tả sự hình thành và lan truyền của sóng điện từ}
\begin{ex}
% ID: L12C3B19-07-DS
% Mức: TH
Xét Mô tả sự hình thành và lan truyền của sóng điện từ (Tình huống 2). Hãy xác định tính đúng, sai của bốn phát biểu sau.
\choiceTF
{Sóng điện từ bắt buộc cần môi trường vật chất để truyền.}
{\True Sóng điện từ là điện từ trường lan truyền trong không gian và truyền được trong chân không.}
{Kết quả không phụ thuộc dữ kiện và có thể bỏ qua điều kiện.}
{\True Khi giải cần xác định đúng dữ kiện, phương chiều, điều kiện áp dụng và đơn vị.}
\loigiai{
\begin{itemchoice}
\item (Sai) Sóng điện từ bắt buộc cần môi trường vật chất để truyền. (Vì): óng điện từ bắt buộc cần môi trường vật chất để truyền. b) Đúng: Sóng điện từ là điện từ trường lan truyền trong không gian và truyền được trong chân không. c) Sai: Kết quả không phụ thuộc dữ kiện và có thể bỏ qua điều kiện. d) Đúng: Khi giải cần xác định đúng dữ kiện, phương chiều, điều kiện áp dụng và đơn vị. Kiến thức cốt lõi: Sóng điện từ là điện từ trường lan truyền trong không gian và truyền được trong chân không
\item (Đúng) Sóng điện từ là điện từ trường lan truyền trong không gian và truyền được trong chân không.
\item (Sai) Kết quả không phụ thuộc dữ kiện và có thể bỏ qua điều kiện.
\item (Đúng) Khi giải cần xác định đúng dữ kiện, phương chiều, điều kiện áp dụng và đơn vị.
\end{itemchoice}
}
\end{ex}

% ===== Câu 26 =====
% ID: L12C3B19-07-TL
% Mức: TH
\begin{bt}
% ID: L12C3B19-07-TL
% Mức: TH
Trình bày kiến thức cốt lõi của Mô tả sự hình thành và lan truyền của sóng điện từ và nêu điều kiện áp dụng.
\loigiai{
Cần nêu được: Sóng điện từ là điện từ trường lan truyền trong không gian và truyền được trong chân không. Khi vận dụng phải xác định đúng dữ kiện, phương chiều nếu có, điều kiện áp dụng, hệ đơn vị và kiểm tra tính hợp lí. Nhận định “Sóng điện từ bắt buộc cần môi trường vật chất để truyền.” sai.
}
\end{bt}

% ===== Câu 27 =====
% ID: L12C3B19-07-TLN
% Mức: TH
\begin{ex}
% ID: L12C3B19-07-TLN
% Mức: TH
Trong bốn phát biểu về Mô tả sự hình thành và lan truyền của sóng điện từ, có bao nhiêu phát biểu đúng? (Chỉ ghi một số nguyên.) (1) Sóng điện từ là điện từ trường lan truyền trong không gian và truyền được trong chân không. (2) Sóng điện từ bắt buộc cần môi trường vật chất để truyền. (3) Cần kiểm tra điều kiện áp dụng. (4) Có thể bỏ qua đơn vị.
\shortans{2}
\loigiai{
Đối chiếu với kiến thức: Sóng điện từ là điện từ trường lan truyền trong không gian và truyền được trong chân không. Có 2 phát biểu đúng.
}
\end{ex}

% ===== Câu 28 =====
% ID: L12C3B19-07-TN
% Mức: TH
\dangbt{Giải thích truyền thông tin bằng sóng điện từ và các tình huống thực tiễn}
\begin{ex}
% ID: L12C3B19-07-TN
% Mức: TH
Chọn phát biểu không trái với kiến thức về Giải thích truyền thông tin bằng sóng điện từ và các tình huống thực tiễn?
\choice
{Kết quả không phụ thuộc dữ kiện và điều kiện áp dụng.}
{\True Thông tin được ghép lên sóng mang bằng biến điệu, truyền đi và được máy thu giải điều chế.}
{Máy thu nhận thông tin mà không cần chọn sóng hay giải điều chế.}
{Có thể bỏ qua phương, chiều và đơn vị trong mọi trường hợp.}
\loigiai{
Thông tin được ghép lên sóng mang bằng biến điệu, truyền đi và được máy thu giải điều chế. Vì vậy phương án $B$ đúng; các phương án còn lại sai.
}
\end{ex}

% ===== Câu 29 =====
% ID: L12C3B19-08-DS
% Mức: VD
\dangbt{Mô tả sự hình thành và lan truyền của sóng điện từ}
\begin{ex}
% ID: L12C3B19-08-DS
% Mức: VD
Xét Mô tả sự hình thành và lan truyền của sóng điện từ (Tình huống 3). Hãy xác định tính đúng, sai của bốn phát biểu sau.
\choiceTF
{\True Khi giải cần xác định đúng dữ kiện, phương chiều, điều kiện áp dụng và đơn vị.}
{Kết quả không phụ thuộc dữ kiện và có thể bỏ qua điều kiện.}
{\True Sóng điện từ là điện từ trường lan truyền trong không gian và truyền được trong chân không.}
{Sóng điện từ bắt buộc cần môi trường vật chất để truyền.}
\loigiai{
\begin{itemchoice}
\item (Đúng) Khi giải cần xác định đúng dữ kiện, phương chiều, điều kiện áp dụng và đơn vị. (Vì): Khi giải cần xác định đúng dữ kiện, phương chiều, điều kiện áp dụng và đơn vị. b) Sai: Kết quả không phụ thuộc dữ kiện và có thể bỏ qua điều kiện. c) Đúng: Sóng điện từ là điện từ trường lan truyền trong không gian và truyền được trong chân không. d) Sai: Sóng điện từ bắt buộc cần môi trường vật chất để truyền. Kiến thức cốt lõi: Sóng điện từ là điện từ trường lan truyền trong không gian và truyền được trong chân không
\item (Sai) Kết quả không phụ thuộc dữ kiện và có thể bỏ qua điều kiện.
\item (Đúng) Sóng điện từ là điện từ trường lan truyền trong không gian và truyền được trong chân không.
\item (Sai) Sóng điện từ bắt buộc cần môi trường vật chất để truyền.
\end{itemchoice}
}
\end{ex}

% ===== Câu 30 =====
% ID: L12C3B19-08-TL
% Mức: TH
\begin{bt}
% ID: L12C3B19-08-TL
% Mức: TH
Giải thích vì sao nhận định sau đúng: “Sóng điện từ là điện từ trường lan truyền trong không gian và truyền được trong chân không.”
\loigiai{
Cần nêu được: Sóng điện từ là điện từ trường lan truyền trong không gian và truyền được trong chân không. Khi vận dụng phải xác định đúng dữ kiện, phương chiều nếu có, điều kiện áp dụng, hệ đơn vị và kiểm tra tính hợp lí. Nhận định “Sóng điện từ bắt buộc cần môi trường vật chất để truyền.” sai.
}
\end{bt}

% ===== Câu 31 =====
% ID: L12C3B19-08-TLN
% Mức: TH
\begin{ex}
% ID: L12C3B19-08-TLN
% Mức: TH
Trong bốn phát biểu về Mô tả sự hình thành và lan truyền của sóng điện từ, có bao nhiêu phát biểu đúng? (Chỉ ghi một số nguyên.) (1) Sóng điện từ bắt buộc cần môi trường vật chất để truyền. (2) Sóng điện từ là điện từ trường lan truyền trong không gian và truyền được trong chân không. (3) Kết quả phải phù hợp ý nghĩa vật lí. (4) Mọi đại lượng đều là vô hướng.
\shortans{1}
\loigiai{
Đối chiếu với kiến thức: Sóng điện từ là điện từ trường lan truyền trong không gian và truyền được trong chân không. Có 1 phát biểu đúng.
}
\end{ex}

% ===== Câu 32 =====
% ID: L12C3B19-08-TN
% Mức: VD
\dangbt{Giải thích truyền thông tin bằng sóng điện từ và các tình huống thực tiễn}
\begin{ex}
% ID: L12C3B19-08-TN
% Mức: VD
Cơ sở Vật lí đúng dùng để giải Giải thích truyền thông tin bằng sóng điện từ và các tình huống thực tiễn?
\choice
{Có thể bỏ qua phương, chiều và đơn vị trong mọi trường hợp.}
{Kết quả không phụ thuộc dữ kiện và điều kiện áp dụng.}
{\True Thông tin được ghép lên sóng mang bằng biến điệu, truyền đi và được máy thu giải điều chế.}
{Máy thu nhận thông tin mà không cần chọn sóng hay giải điều chế.}
\loigiai{
Thông tin được ghép lên sóng mang bằng biến điệu, truyền đi và được máy thu giải điều chế. Vì vậy phương án $C$ đúng; các phương án còn lại sai.
}
\end{ex}

% ===== Câu 33 =====
% ID: L12C3B19-09-DS
% Mức: VD
\dangbt{Mô tả sự hình thành và lan truyền của sóng điện từ}
\begin{ex}
% ID: L12C3B19-09-DS
% Mức: VD
Xét Mô tả sự hình thành và lan truyền của sóng điện từ (Tình huống 4). Hãy xác định tính đúng, sai của bốn phát biểu sau.
\choiceTF
{Kết quả không phụ thuộc dữ kiện và có thể bỏ qua điều kiện.}
{\True Khi giải cần xác định đúng dữ kiện, phương chiều, điều kiện áp dụng và đơn vị.}
{Sóng điện từ bắt buộc cần môi trường vật chất để truyền.}
{\True Sóng điện từ là điện từ trường lan truyền trong không gian và truyền được trong chân không.}
\loigiai{
\begin{itemchoice}
\item (Sai) Kết quả không phụ thuộc dữ kiện và có thể bỏ qua điều kiện. (Vì): Kết quả không phụ thuộc dữ kiện và có thể bỏ qua điều kiện. b) Đúng: Khi giải cần xác định đúng dữ kiện, phương chiều, điều kiện áp dụng và đơn vị. c) Sai: Sóng điện từ bắt buộc cần môi trường vật chất để truyền. d) Đúng: Sóng điện từ là điện từ trường lan truyền trong không gian và truyền được trong chân không. Kiến thức cốt lõi: Sóng điện từ là điện từ trường lan truyền trong không gian và truyền được trong chân không
\item (Đúng) Khi giải cần xác định đúng dữ kiện, phương chiều, điều kiện áp dụng và đơn vị.
\item (Sai) Sóng điện từ bắt buộc cần môi trường vật chất để truyền.
\item (Đúng) Sóng điện từ là điện từ trường lan truyền trong không gian và truyền được trong chân không.
\end{itemchoice}
}
\end{ex}

% ===== Câu 34 =====
% ID: L12C3B19-09-TL
% Mức: VD
\begin{bt}
% ID: L12C3B19-09-TL
% Mức: VD
Phân tích sai lầm trong nhận định: “Sóng điện từ bắt buộc cần môi trường vật chất để truyền.”
\loigiai{
Cần nêu được: Sóng điện từ là điện từ trường lan truyền trong không gian và truyền được trong chân không. Khi vận dụng phải xác định đúng dữ kiện, phương chiều nếu có, điều kiện áp dụng, hệ đơn vị và kiểm tra tính hợp lí. Nhận định “Sóng điện từ bắt buộc cần môi trường vật chất để truyền.” sai.
}
\end{bt}

% ===== Câu 35 =====
% ID: L12C3B19-09-TLN
% Mức: TH
\begin{ex}
% ID: L12C3B19-09-TLN
% Mức: TH
Trong bốn phát biểu về Mô tả sự hình thành và lan truyền của sóng điện từ, có bao nhiêu phát biểu đúng? (Chỉ ghi một số nguyên.) (1) Sóng điện từ là điện từ trường lan truyền trong không gian và truyền được trong chân không. (2) Sóng điện từ bắt buộc cần môi trường vật chất để truyền. (3) Cần kiểm tra điều kiện áp dụng. (4) Có thể bỏ qua đơn vị.
\shortans{2}
\loigiai{
Đối chiếu với kiến thức: Sóng điện từ là điện từ trường lan truyền trong không gian và truyền được trong chân không. Có 2 phát biểu đúng.
}
\end{ex}

% ===== Câu 36 =====
% ID: L12C3B19-09-TN
% Mức: VD
\dangbt{Giải thích truyền thông tin bằng sóng điện từ và các tình huống thực tiễn}
\begin{ex}
% ID: L12C3B19-09-TN
% Mức: VD
Trong một bài thực hành về Giải thích truyền thông tin bằng sóng điện từ và các tình huống thực tiễn?
\choice
{Máy thu nhận thông tin mà không cần chọn sóng hay giải điều chế.}
{Có thể bỏ qua phương, chiều và đơn vị trong mọi trường hợp.}
{Kết quả không phụ thuộc dữ kiện và điều kiện áp dụng.}
{\True Thông tin được ghép lên sóng mang bằng biến điệu, truyền đi và được máy thu giải điều chế.}
\loigiai{
Thông tin được ghép lên sóng mang bằng biến điệu, truyền đi và được máy thu giải điều chế. Vì vậy phương án $D$ đúng; các phương án còn lại sai.
}
\end{ex}

% ===== Câu 37 =====
% ID: L12C3B19-10-DS
% Mức: VD
\dangbt{Mô tả sự hình thành và lan truyền của sóng điện từ}
\begin{ex}
% ID: L12C3B19-10-DS
% Mức: VD
Xét Mô tả sự hình thành và lan truyền của sóng điện từ (Tình huống 5). Hãy xác định tính đúng, sai của bốn phát biểu sau.
\choiceTF
{\True Sóng điện từ là điện từ trường lan truyền trong không gian và truyền được trong chân không.}
{Kết quả không phụ thuộc dữ kiện và có thể bỏ qua điều kiện.}
{Sóng điện từ bắt buộc cần môi trường vật chất để truyền.}
{\True Khi giải cần xác định đúng dữ kiện, phương chiều, điều kiện áp dụng và đơn vị.}
\loigiai{
\begin{itemchoice}
\item (Đúng) Sóng điện từ là điện từ trường lan truyền trong không gian và truyền được trong chân không. (Vì): Sóng điện từ là điện từ trường lan truyền trong không gian và truyền được trong chân không. b) Sai: Kết quả không phụ thuộc dữ kiện và có thể bỏ qua điều kiện. c) Sai: Sóng điện từ bắt buộc cần môi trường vật chất để truyền. d) Đúng: Khi giải cần xác định đúng dữ kiện, phương chiều, điều kiện áp dụng và đơn vị. Kiến thức cốt lõi: Sóng điện từ là điện từ trường lan truyền trong không gian và truyền được trong chân không
\item (Sai) Kết quả không phụ thuộc dữ kiện và có thể bỏ qua điều kiện.
\item (Sai) Sóng điện từ bắt buộc cần môi trường vật chất để truyền.
\item (Đúng) Khi giải cần xác định đúng dữ kiện, phương chiều, điều kiện áp dụng và đơn vị.
\end{itemchoice}
}
\end{ex}

% ===== Câu 38 =====
% ID: L12C3B19-10-TL
% Mức: VD
\begin{bt}
% ID: L12C3B19-10-TL
% Mức: VD
Đề xuất các bước giải một bài toán thuộc dạng Mô tả sự hình thành và lan truyền của sóng điện từ.
\loigiai{
Cần nêu được: Sóng điện từ là điện từ trường lan truyền trong không gian và truyền được trong chân không. Khi vận dụng phải xác định đúng dữ kiện, phương chiều nếu có, điều kiện áp dụng, hệ đơn vị và kiểm tra tính hợp lí. Nhận định “Sóng điện từ bắt buộc cần môi trường vật chất để truyền.” sai.
}
\end{bt}

% ===== Câu 39 =====
% ID: L12C3B19-10-TLN
% Mức: VD
\begin{ex}
% ID: L12C3B19-10-TLN
% Mức: VD
Trong bốn phát biểu về Mô tả sự hình thành và lan truyền của sóng điện từ, có bao nhiêu phát biểu đúng? (Chỉ ghi một số nguyên.) (1) Sóng điện từ bắt buộc cần môi trường vật chất để truyền. (2) Sóng điện từ là điện từ trường lan truyền trong không gian và truyền được trong chân không. (3) Kết quả phải phù hợp ý nghĩa vật lí. (4) Mọi đại lượng đều là vô hướng.
\shortans{2}
\loigiai{
Đối chiếu với kiến thức: Sóng điện từ là điện từ trường lan truyền trong không gian và truyền được trong chân không. Có 2 phát biểu đúng.
}
\end{ex}

% ===== Câu 40 =====
% ID: L12C3B19-10-TN
% Mức: VD
\dangbt{Giải thích truyền thông tin bằng sóng điện từ và các tình huống thực tiễn}
\begin{ex}
% ID: L12C3B19-10-TN
% Mức: VD
Kết luận đúng khi vận dụng Giải thích truyền thông tin bằng sóng điện từ và các tình huống thực tiễn?
\choice
{\True Thông tin được ghép lên sóng mang bằng biến điệu, truyền đi và được máy thu giải điều chế.}
{Máy thu nhận thông tin mà không cần chọn sóng hay giải điều chế.}
{Có thể bỏ qua phương, chiều và đơn vị trong mọi trường hợp.}
{Kết quả không phụ thuộc dữ kiện và điều kiện áp dụng.}
\loigiai{
Thông tin được ghép lên sóng mang bằng biến điệu, truyền đi và được máy thu giải điều chế. Vì vậy phương án $A$ đúng; các phương án còn lại sai.
}
\end{ex}

% ===== Câu 41 =====
% ID: L12C3B19-100-TN
% Mức: VD
\dangbt{Tính bước sóng, tần số và tốc độ truyền sóng điện từ}
\begin{ex}
% ID: L12C3B19-100-TN
% Mức: VD
Kết luận đúng khi vận dụng Tính bước sóng, tần số và tốc độ truyền sóng điện từ?
\choice
{\True Trong chân không, bước sóng lambda = c/f với c xấp xỉ 3.10^ $8\,\text{m}$ /s.}
{Bước sóng trong chân không bằng f/c.}
{Có thể bỏ qua phương, chiều và đơn vị trong mọi trường hợp.}
{Kết quả không phụ thuộc dữ kiện và điều kiện áp dụng.}
\loigiai{
Trong chân không, bước sóng lambda = c/f với c xấp xỉ 3.10^ $8\,\text{m}$ /s. Vì vậy phương án $A$ đúng; các phương án còn lại sai.
}
\end{ex}

% ===== Câu 42 =====
% ID: L12C3B19-101-TN
% Mức: NB
\dangbt{Phân loại phổ sóng điện từ và nêu ứng dụng}
\begin{ex}
% ID: L12C3B19-101-TN
% Mức: NB
Phát biểu nào sau đây đúng về Phân loại phổ sóng điện từ và nêu ứng dụng?
\choice
{Âm thanh là một miền của phổ sóng điện từ.}
{Có thể bỏ qua phương, chiều và đơn vị trong mọi trường hợp.}
{Kết quả không phụ thuộc dữ kiện và điều kiện áp dụng.}
{\True Phổ điện từ gồm sóng vô tuyến, vi ba, hồng ngoại, ánh sáng nhìn thấy, tử ngoại, tia $X$ và gamma.}
\loigiai{
Phổ điện từ gồm sóng vô tuyến, vi ba, hồng ngoại, ánh sáng nhìn thấy, tử ngoại, tia $X$ và gamma. Vì vậy phương án $D$ đúng; các phương án còn lại sai.
}
\end{ex}

% ===== Câu 43 =====
% ID: L12C3B19-102-TN
% Mức: NB
\begin{ex}
% ID: L12C3B19-102-TN
% Mức: NB
Trong nội dung Phân loại phổ sóng điện từ và nêu ứng dụng?
\choice
{\True Phổ điện từ gồm sóng vô tuyến, vi ba, hồng ngoại, ánh sáng nhìn thấy, tử ngoại, tia $X$ và gamma.}
{Âm thanh là một miền của phổ sóng điện từ.}
{Có thể bỏ qua phương, chiều và đơn vị trong mọi trường hợp.}
{Kết quả không phụ thuộc dữ kiện và điều kiện áp dụng.}
\loigiai{
Phổ điện từ gồm sóng vô tuyến, vi ba, hồng ngoại, ánh sáng nhìn thấy, tử ngoại, tia $X$ và gamma. Vì vậy phương án $A$ đúng; các phương án còn lại sai.
}
\end{ex}

% ===== Câu 44 =====
% ID: L12C3B19-103-TN
% Mức: NB
\begin{ex}
% ID: L12C3B19-103-TN
% Mức: NB
Tình huống 3: chọn kết luận chính xác về Phân loại phổ sóng điện từ và nêu ứng dụng?
\choice
{Kết quả không phụ thuộc dữ kiện và điều kiện áp dụng.}
{\True Phổ điện từ gồm sóng vô tuyến, vi ba, hồng ngoại, ánh sáng nhìn thấy, tử ngoại, tia $X$ và gamma.}
{Âm thanh là một miền của phổ sóng điện từ.}
{Có thể bỏ qua phương, chiều và đơn vị trong mọi trường hợp.}
\loigiai{
Phổ điện từ gồm sóng vô tuyến, vi ba, hồng ngoại, ánh sáng nhìn thấy, tử ngoại, tia $X$ và gamma. Vì vậy phương án $B$ đúng; các phương án còn lại sai.
}
\end{ex}

% ===== Câu 45 =====
% ID: L12C3B19-104-TN
% Mức: TH
\begin{ex}
% ID: L12C3B19-104-TN
% Mức: TH
Nội dung nào mô tả đúng nhất Phân loại phổ sóng điện từ và nêu ứng dụng?
\choice
{Có thể bỏ qua phương, chiều và đơn vị trong mọi trường hợp.}
{Kết quả không phụ thuộc dữ kiện và điều kiện áp dụng.}
{\True Phổ điện từ gồm sóng vô tuyến, vi ba, hồng ngoại, ánh sáng nhìn thấy, tử ngoại, tia $X$ và gamma.}
{Âm thanh là một miền của phổ sóng điện từ.}
\loigiai{
Phổ điện từ gồm sóng vô tuyến, vi ba, hồng ngoại, ánh sáng nhìn thấy, tử ngoại, tia $X$ và gamma. Vì vậy phương án $C$ đúng; các phương án còn lại sai.
}
\end{ex}

% ===== Câu 46 =====
% ID: L12C3B19-105-TN
% Mức: TH
\begin{ex}
% ID: L12C3B19-105-TN
% Mức: TH
Khi phân tích Phân loại phổ sóng điện từ và nêu ứng dụng?
\choice
{Âm thanh là một miền của phổ sóng điện từ.}
{Có thể bỏ qua phương, chiều và đơn vị trong mọi trường hợp.}
{Kết quả không phụ thuộc dữ kiện và điều kiện áp dụng.}
{\True Phổ điện từ gồm sóng vô tuyến, vi ba, hồng ngoại, ánh sáng nhìn thấy, tử ngoại, tia $X$ và gamma.}
\loigiai{
Phổ điện từ gồm sóng vô tuyến, vi ba, hồng ngoại, ánh sáng nhìn thấy, tử ngoại, tia $X$ và gamma. Vì vậy phương án $D$ đúng; các phương án còn lại sai.
}
\end{ex}

% ===== Câu 47 =====
% ID: L12C3B19-106-TN
% Mức: TH
\begin{ex}
% ID: L12C3B19-106-TN
% Mức: TH
Một học sinh ôn tập Phân loại phổ sóng điện từ và nêu ứng dụng?
\choice
{\True Phổ điện từ gồm sóng vô tuyến, vi ba, hồng ngoại, ánh sáng nhìn thấy, tử ngoại, tia $X$ và gamma.}
{Âm thanh là một miền của phổ sóng điện từ.}
{Có thể bỏ qua phương, chiều và đơn vị trong mọi trường hợp.}
{Kết quả không phụ thuộc dữ kiện và điều kiện áp dụng.}
\loigiai{
Phổ điện từ gồm sóng vô tuyến, vi ba, hồng ngoại, ánh sáng nhìn thấy, tử ngoại, tia $X$ và gamma. Vì vậy phương án $A$ đúng; các phương án còn lại sai.
}
\end{ex}

% ===== Câu 48 =====
% ID: L12C3B19-107-TN
% Mức: TH
\begin{ex}
% ID: L12C3B19-107-TN
% Mức: TH
Chọn phát biểu không trái với kiến thức về Phân loại phổ sóng điện từ và nêu ứng dụng?
\choice
{Kết quả không phụ thuộc dữ kiện và điều kiện áp dụng.}
{\True Phổ điện từ gồm sóng vô tuyến, vi ba, hồng ngoại, ánh sáng nhìn thấy, tử ngoại, tia $X$ và gamma.}
{Âm thanh là một miền của phổ sóng điện từ.}
{Có thể bỏ qua phương, chiều và đơn vị trong mọi trường hợp.}
\loigiai{
Phổ điện từ gồm sóng vô tuyến, vi ba, hồng ngoại, ánh sáng nhìn thấy, tử ngoại, tia $X$ và gamma. Vì vậy phương án $B$ đúng; các phương án còn lại sai.
}
\end{ex}

% ===== Câu 49 =====
% ID: L12C3B19-108-TN
% Mức: VD
\begin{ex}
% ID: L12C3B19-108-TN
% Mức: VD
Cơ sở Vật lí đúng dùng để giải Phân loại phổ sóng điện từ và nêu ứng dụng?
\choice
{Có thể bỏ qua phương, chiều và đơn vị trong mọi trường hợp.}
{Kết quả không phụ thuộc dữ kiện và điều kiện áp dụng.}
{\True Phổ điện từ gồm sóng vô tuyến, vi ba, hồng ngoại, ánh sáng nhìn thấy, tử ngoại, tia $X$ và gamma.}
{Âm thanh là một miền của phổ sóng điện từ.}
\loigiai{
Phổ điện từ gồm sóng vô tuyến, vi ba, hồng ngoại, ánh sáng nhìn thấy, tử ngoại, tia $X$ và gamma. Vì vậy phương án $C$ đúng; các phương án còn lại sai.
}
\end{ex}

% ===== Câu 50 =====
% ID: L12C3B19-109-TN
% Mức: VD
\begin{ex}
% ID: L12C3B19-109-TN
% Mức: VD
Trong một bài thực hành về Phân loại phổ sóng điện từ và nêu ứng dụng?
\choice
{Âm thanh là một miền của phổ sóng điện từ.}
{Có thể bỏ qua phương, chiều và đơn vị trong mọi trường hợp.}
{Kết quả không phụ thuộc dữ kiện và điều kiện áp dụng.}
{\True Phổ điện từ gồm sóng vô tuyến, vi ba, hồng ngoại, ánh sáng nhìn thấy, tử ngoại, tia $X$ và gamma.}
\loigiai{
Phổ điện từ gồm sóng vô tuyến, vi ba, hồng ngoại, ánh sáng nhìn thấy, tử ngoại, tia $X$ và gamma. Vì vậy phương án $D$ đúng; các phương án còn lại sai.
}
\end{ex}

% ===== Câu 51 =====
% ID: L12C3B19-11-DS
% Mức: TH
\dangbt{Nhận biết các tính chất cơ bản của sóng điện từ}
\begin{ex}
% ID: L12C3B19-11-DS
% Mức: TH
Xét Nhận biết các tính chất cơ bản của sóng điện từ (Tình huống 1). Hãy xác định tính đúng, sai của bốn phát biểu sau.
\choiceTF
{\True Trong chân không, vectơ $E$ và $B$ vuông góc nhau và cùng vuông góc phương truyền sóng.}
{Trong sóng điện từ, $E$ và $B$ song song với phương truyền.}
{\True Khi giải cần xác định đúng dữ kiện, phương chiều, điều kiện áp dụng và đơn vị.}
{Kết quả không phụ thuộc dữ kiện và có thể bỏ qua điều kiện.}
\loigiai{
\begin{itemchoice}
\item (Đúng) Trong chân không, vectơ $E$ và $B$ vuông góc nhau và cùng vuông góc phương truyền sóng. (Vì): Trong chân không, vectơ $E$ và $B$ vuông góc nhau và cùng vuông góc phương truyền sóng. b) Sai: Trong sóng điện từ, $E$ và $B$ song song với phương truyền. c) Đúng: Khi giải cần xác định đúng dữ kiện, phương chiều, điều kiện áp dụng và đơn vị. d) Sai: Kết quả không phụ thuộc dữ kiện và có thể bỏ qua điều kiện. Kiến thức cốt lõi: Trong chân không, vectơ $E$ và $B$ vuông góc nhau và cùng vuông góc phương truyền sóng
\item (Sai) Trong sóng điện từ, $E$ và $B$ song song với phương truyền.
\item (Đúng) Khi giải cần xác định đúng dữ kiện, phương chiều, điều kiện áp dụng và đơn vị.
\item (Sai) Kết quả không phụ thuộc dữ kiện và có thể bỏ qua điều kiện.
\end{itemchoice}
}
\end{ex}

% ===== Câu 52 =====
% ID: L12C3B19-11-TL
% Mức: VD
\dangbt{Mô tả sự hình thành và lan truyền của sóng điện từ}
\begin{bt}
% ID: L12C3B19-11-TL
% Mức: VD
Nêu một tình huống thực tiễn liên quan đến Mô tả sự hình thành và lan truyền của sóng điện từ và giải thích bằng kiến thức Vật lí.
\loigiai{
Cần nêu được: Sóng điện từ là điện từ trường lan truyền trong không gian và truyền được trong chân không. Khi vận dụng phải xác định đúng dữ kiện, phương chiều nếu có, điều kiện áp dụng, hệ đơn vị và kiểm tra tính hợp lí. Nhận định “Sóng điện từ bắt buộc cần môi trường vật chất để truyền.” sai.
}
\end{bt}

% ===== Câu 53 =====
% ID: L12C3B19-11-TLN
% Mức: VD
\begin{ex}
% ID: L12C3B19-11-TLN
% Mức: VD
Trong bốn phát biểu về Mô tả sự hình thành và lan truyền của sóng điện từ, có bao nhiêu phát biểu đúng? (Chỉ ghi một số nguyên.) (1) Sóng điện từ là điện từ trường lan truyền trong không gian và truyền được trong chân không. (2) Sóng điện từ bắt buộc cần môi trường vật chất để truyền. (3) Cần kiểm tra điều kiện áp dụng. (4) Có thể bỏ qua đơn vị.
\shortans{2}
\loigiai{
Đối chiếu với kiến thức: Sóng điện từ là điện từ trường lan truyền trong không gian và truyền được trong chân không. Có 2 phát biểu đúng.
}
\end{ex}

% ===== Câu 54 =====
% ID: L12C3B19-11-TN
% Mức: NB
\begin{ex}
% ID: L12C3B19-11-TN
% Mức: NB
Phát biểu nào sau đây đúng về Mô tả sự hình thành và lan truyền của sóng điện từ?
\choice
{Sóng điện từ bắt buộc cần môi trường vật chất để truyền.}
{Có thể bỏ qua phương, chiều và đơn vị trong mọi trường hợp.}
{Kết quả không phụ thuộc dữ kiện và điều kiện áp dụng.}
{\True Sóng điện từ là điện từ trường lan truyền trong không gian và truyền được trong chân không.}
\loigiai{
Sóng điện từ là điện từ trường lan truyền trong không gian và truyền được trong chân không. Vì vậy phương án $D$ đúng; các phương án còn lại sai.
}
\end{ex}

% ===== Câu 55 =====
% ID: L12C3B19-110-TN
% Mức: VD
\dangbt{Phân loại phổ sóng điện từ và nêu ứng dụng}
\begin{ex}
% ID: L12C3B19-110-TN
% Mức: VD
Kết luận đúng khi vận dụng Phân loại phổ sóng điện từ và nêu ứng dụng?
\choice
{\True Phổ điện từ gồm sóng vô tuyến, vi ba, hồng ngoại, ánh sáng nhìn thấy, tử ngoại, tia $X$ và gamma.}
{Âm thanh là một miền của phổ sóng điện từ.}
{Có thể bỏ qua phương, chiều và đơn vị trong mọi trường hợp.}
{Kết quả không phụ thuộc dữ kiện và điều kiện áp dụng.}
\loigiai{
Phổ điện từ gồm sóng vô tuyến, vi ba, hồng ngoại, ánh sáng nhìn thấy, tử ngoại, tia $X$ và gamma. Vì vậy phương án $A$ đúng; các phương án còn lại sai.
}
\end{ex}

% ===== Câu 56 =====
% ID: L12C3B19-111-TN
% Mức: NB
\dangbt{Giải thích truyền thông tin bằng sóng điện từ và các tình huống thực tiễn}
\begin{ex}
% ID: L12C3B19-111-TN
% Mức: NB
Phát biểu nào sau đây đúng về Giải thích truyền thông tin bằng sóng điện từ và các tình huống thực tiễn?
\choice
{Máy thu nhận thông tin mà không cần chọn sóng hay giải điều chế.}
{Có thể bỏ qua phương, chiều và đơn vị trong mọi trường hợp.}
{Kết quả không phụ thuộc dữ kiện và điều kiện áp dụng.}
{\True Thông tin được ghép lên sóng mang bằng biến điệu, truyền đi và được máy thu giải điều chế.}
\loigiai{
Thông tin được ghép lên sóng mang bằng biến điệu, truyền đi và được máy thu giải điều chế. Vì vậy phương án $D$ đúng; các phương án còn lại sai.
}
\end{ex}

% ===== Câu 57 =====
% ID: L12C3B19-112-TN
% Mức: NB
\begin{ex}
% ID: L12C3B19-112-TN
% Mức: NB
Trong nội dung Giải thích truyền thông tin bằng sóng điện từ và các tình huống thực tiễn?
\choice
{\True Thông tin được ghép lên sóng mang bằng biến điệu, truyền đi và được máy thu giải điều chế.}
{Máy thu nhận thông tin mà không cần chọn sóng hay giải điều chế.}
{Có thể bỏ qua phương, chiều và đơn vị trong mọi trường hợp.}
{Kết quả không phụ thuộc dữ kiện và điều kiện áp dụng.}
\loigiai{
Thông tin được ghép lên sóng mang bằng biến điệu, truyền đi và được máy thu giải điều chế. Vì vậy phương án $A$ đúng; các phương án còn lại sai.
}
\end{ex}

% ===== Câu 58 =====
% ID: L12C3B19-113-TN
% Mức: NB
\begin{ex}
% ID: L12C3B19-113-TN
% Mức: NB
Tình huống 3: chọn kết luận chính xác về Giải thích truyền thông tin bằng sóng điện từ và các tình huống thực tiễn?
\choice
{Kết quả không phụ thuộc dữ kiện và điều kiện áp dụng.}
{\True Thông tin được ghép lên sóng mang bằng biến điệu, truyền đi và được máy thu giải điều chế.}
{Máy thu nhận thông tin mà không cần chọn sóng hay giải điều chế.}
{Có thể bỏ qua phương, chiều và đơn vị trong mọi trường hợp.}
\loigiai{
Thông tin được ghép lên sóng mang bằng biến điệu, truyền đi và được máy thu giải điều chế. Vì vậy phương án $B$ đúng; các phương án còn lại sai.
}
\end{ex}

% ===== Câu 59 =====
% ID: L12C3B19-114-TN
% Mức: TH
\begin{ex}
% ID: L12C3B19-114-TN
% Mức: TH
Nội dung nào mô tả đúng nhất Giải thích truyền thông tin bằng sóng điện từ và các tình huống thực tiễn?
\choice
{Có thể bỏ qua phương, chiều và đơn vị trong mọi trường hợp.}
{Kết quả không phụ thuộc dữ kiện và điều kiện áp dụng.}
{\True Thông tin được ghép lên sóng mang bằng biến điệu, truyền đi và được máy thu giải điều chế.}
{Máy thu nhận thông tin mà không cần chọn sóng hay giải điều chế.}
\loigiai{
Thông tin được ghép lên sóng mang bằng biến điệu, truyền đi và được máy thu giải điều chế. Vì vậy phương án $C$ đúng; các phương án còn lại sai.
}
\end{ex}

% ===== Câu 60 =====
% ID: L12C3B19-115-TN
% Mức: TH
\begin{ex}
% ID: L12C3B19-115-TN
% Mức: TH
Khi phân tích Giải thích truyền thông tin bằng sóng điện từ và các tình huống thực tiễn?
\choice
{Máy thu nhận thông tin mà không cần chọn sóng hay giải điều chế.}
{Có thể bỏ qua phương, chiều và đơn vị trong mọi trường hợp.}
{Kết quả không phụ thuộc dữ kiện và điều kiện áp dụng.}
{\True Thông tin được ghép lên sóng mang bằng biến điệu, truyền đi và được máy thu giải điều chế.}
\loigiai{
Thông tin được ghép lên sóng mang bằng biến điệu, truyền đi và được máy thu giải điều chế. Vì vậy phương án $D$ đúng; các phương án còn lại sai.
}
\end{ex}

% ===== Câu 61 =====
% ID: L12C3B19-116-TN
% Mức: TH
\begin{ex}
% ID: L12C3B19-116-TN
% Mức: TH
Một học sinh ôn tập Giải thích truyền thông tin bằng sóng điện từ và các tình huống thực tiễn?
\choice
{\True Thông tin được ghép lên sóng mang bằng biến điệu, truyền đi và được máy thu giải điều chế.}
{Máy thu nhận thông tin mà không cần chọn sóng hay giải điều chế.}
{Có thể bỏ qua phương, chiều và đơn vị trong mọi trường hợp.}
{Kết quả không phụ thuộc dữ kiện và điều kiện áp dụng.}
\loigiai{
Thông tin được ghép lên sóng mang bằng biến điệu, truyền đi và được máy thu giải điều chế. Vì vậy phương án $A$ đúng; các phương án còn lại sai.
}
\end{ex}

% ===== Câu 62 =====
% ID: L12C3B19-117-TN
% Mức: TH
\begin{ex}
% ID: L12C3B19-117-TN
% Mức: TH
Chọn phát biểu không trái với kiến thức về Giải thích truyền thông tin bằng sóng điện từ và các tình huống thực tiễn?
\choice
{Kết quả không phụ thuộc dữ kiện và điều kiện áp dụng.}
{\True Thông tin được ghép lên sóng mang bằng biến điệu, truyền đi và được máy thu giải điều chế.}
{Máy thu nhận thông tin mà không cần chọn sóng hay giải điều chế.}
{Có thể bỏ qua phương, chiều và đơn vị trong mọi trường hợp.}
\loigiai{
Thông tin được ghép lên sóng mang bằng biến điệu, truyền đi và được máy thu giải điều chế. Vì vậy phương án $B$ đúng; các phương án còn lại sai.
}
\end{ex}

% ===== Câu 63 =====
% ID: L12C3B19-118-TN
% Mức: VD
\begin{ex}
% ID: L12C3B19-118-TN
% Mức: VD
Cơ sở Vật lí đúng dùng để giải Giải thích truyền thông tin bằng sóng điện từ và các tình huống thực tiễn?
\choice
{Có thể bỏ qua phương, chiều và đơn vị trong mọi trường hợp.}
{Kết quả không phụ thuộc dữ kiện và điều kiện áp dụng.}
{\True Thông tin được ghép lên sóng mang bằng biến điệu, truyền đi và được máy thu giải điều chế.}
{Máy thu nhận thông tin mà không cần chọn sóng hay giải điều chế.}
\loigiai{
Thông tin được ghép lên sóng mang bằng biến điệu, truyền đi và được máy thu giải điều chế. Vì vậy phương án $C$ đúng; các phương án còn lại sai.
}
\end{ex}

% ===== Câu 64 =====
% ID: L12C3B19-119-TN
% Mức: VD
\begin{ex}
% ID: L12C3B19-119-TN
% Mức: VD
Trong một bài thực hành về Giải thích truyền thông tin bằng sóng điện từ và các tình huống thực tiễn?
\choice
{Máy thu nhận thông tin mà không cần chọn sóng hay giải điều chế.}
{Có thể bỏ qua phương, chiều và đơn vị trong mọi trường hợp.}
{Kết quả không phụ thuộc dữ kiện và điều kiện áp dụng.}
{\True Thông tin được ghép lên sóng mang bằng biến điệu, truyền đi và được máy thu giải điều chế.}
\loigiai{
Thông tin được ghép lên sóng mang bằng biến điệu, truyền đi và được máy thu giải điều chế. Vì vậy phương án $D$ đúng; các phương án còn lại sai.
}
\end{ex}

% ===== Câu 65 =====
% ID: L12C3B19-12-DS
% Mức: TH
\dangbt{Nhận biết các tính chất cơ bản của sóng điện từ}
\begin{ex}
% ID: L12C3B19-12-DS
% Mức: TH
Xét Nhận biết các tính chất cơ bản của sóng điện từ (Tình huống 2). Hãy xác định tính đúng, sai của bốn phát biểu sau.
\choiceTF
{Trong sóng điện từ, $E$ và $B$ song song với phương truyền.}
{\True Trong chân không, vectơ $E$ và $B$ vuông góc nhau và cùng vuông góc phương truyền sóng.}
{Kết quả không phụ thuộc dữ kiện và có thể bỏ qua điều kiện.}
{\True Khi giải cần xác định đúng dữ kiện, phương chiều, điều kiện áp dụng và đơn vị.}
\loigiai{
\begin{itemchoice}
\item (Sai) Trong sóng điện từ, $E$ và $B$ song song với phương truyền. (Vì): Trong sóng điện từ, $E$ và $B$ song song với phương truyền. b) Đúng: Trong chân không, vectơ $E$ và $B$ vuông góc nhau và cùng vuông góc phương truyền sóng. c) Sai: Kết quả không phụ thuộc dữ kiện và có thể bỏ qua điều kiện. d) Đúng: Khi giải cần xác định đúng dữ kiện, phương chiều, điều kiện áp dụng và đơn vị. Kiến thức cốt lõi: Trong chân không, vectơ $E$ và $B$ vuông góc nhau và cùng vuông góc phương truyền sóng
\item (Đúng) Trong chân không, vectơ $E$ và $B$ vuông góc nhau và cùng vuông góc phương truyền sóng.
\item (Sai) Kết quả không phụ thuộc dữ kiện và có thể bỏ qua điều kiện.
\item (Đúng) Khi giải cần xác định đúng dữ kiện, phương chiều, điều kiện áp dụng và đơn vị.
\end{itemchoice}
}
\end{ex}

% ===== Câu 66 =====
% ID: L12C3B19-12-TL
% Mức: VDC
\dangbt{Mô tả sự hình thành và lan truyền của sóng điện từ}
\begin{bt}
% ID: L12C3B19-12-TL
% Mức: VDC
So sánh nhận định đúng “Sóng điện từ là điện từ trường lan truyền trong không gian và truyền được trong chân không.” với nhận định sai “Sóng điện từ bắt buộc cần môi trường vật chất để truyền.”; chỉ rõ căn cứ Vật lí.
\loigiai{
Cần nêu được: Sóng điện từ là điện từ trường lan truyền trong không gian và truyền được trong chân không. Khi vận dụng phải xác định đúng dữ kiện, phương chiều nếu có, điều kiện áp dụng, hệ đơn vị và kiểm tra tính hợp lí. Nhận định “Sóng điện từ bắt buộc cần môi trường vật chất để truyền.” sai.
}
\end{bt}

% ===== Câu 67 =====
% ID: L12C3B19-12-TLN
% Mức: VDC
\begin{ex}
% ID: L12C3B19-12-TLN
% Mức: VDC
Trong bốn phát biểu về Mô tả sự hình thành và lan truyền của sóng điện từ, có bao nhiêu phát biểu đúng? (Chỉ ghi một số nguyên.) (1) Khi giải cần xác định đúng dữ kiện, phương chiều, điều kiện áp dụng và đơn vị. (2) Sóng điện từ bắt buộc cần môi trường vật chất để truyền. (3) Phải dùng đúng định luật cho tình huống. (4) Sóng điện từ là điện từ trường lan truyền trong không gian và truyền được trong chân không.
\shortans{3}
\loigiai{
Đối chiếu với kiến thức: Sóng điện từ là điện từ trường lan truyền trong không gian và truyền được trong chân không. Có 3 phát biểu đúng.
}
\end{ex}

% ===== Câu 68 =====
% ID: L12C3B19-12-TN
% Mức: NB
\begin{ex}
% ID: L12C3B19-12-TN
% Mức: NB
Trong nội dung Mô tả sự hình thành và lan truyền của sóng điện từ?
\choice
{\True Sóng điện từ là điện từ trường lan truyền trong không gian và truyền được trong chân không.}
{Sóng điện từ bắt buộc cần môi trường vật chất để truyền.}
{Có thể bỏ qua phương, chiều và đơn vị trong mọi trường hợp.}
{Kết quả không phụ thuộc dữ kiện và điều kiện áp dụng.}
\loigiai{
Sóng điện từ là điện từ trường lan truyền trong không gian và truyền được trong chân không. Vì vậy phương án $A$ đúng; các phương án còn lại sai.
}
\end{ex}

% ===== Câu 69 =====
% ID: L12C3B19-120-TN
% Mức: VD
\dangbt{Giải thích truyền thông tin bằng sóng điện từ và các tình huống thực tiễn}
\begin{ex}
% ID: L12C3B19-120-TN
% Mức: VD
Kết luận đúng khi vận dụng Giải thích truyền thông tin bằng sóng điện từ và các tình huống thực tiễn?
\choice
{\True Thông tin được ghép lên sóng mang bằng biến điệu, truyền đi và được máy thu giải điều chế.}
{Máy thu nhận thông tin mà không cần chọn sóng hay giải điều chế.}
{Có thể bỏ qua phương, chiều và đơn vị trong mọi trường hợp.}
{Kết quả không phụ thuộc dữ kiện và điều kiện áp dụng.}
\loigiai{
Thông tin được ghép lên sóng mang bằng biến điệu, truyền đi và được máy thu giải điều chế. Vì vậy phương án $A$ đúng; các phương án còn lại sai.
}
\end{ex}

% ===== Câu 70 =====
% ID: L12C3B19-13-DS
% Mức: VD
\dangbt{Nhận biết các tính chất cơ bản của sóng điện từ}
\begin{ex}
% ID: L12C3B19-13-DS
% Mức: VD
Xét Nhận biết các tính chất cơ bản của sóng điện từ (Tình huống 3). Hãy xác định tính đúng, sai của bốn phát biểu sau.
\choiceTF
{\True Khi giải cần xác định đúng dữ kiện, phương chiều, điều kiện áp dụng và đơn vị.}
{Kết quả không phụ thuộc dữ kiện và có thể bỏ qua điều kiện.}
{\True Trong chân không, vectơ $E$ và $B$ vuông góc nhau và cùng vuông góc phương truyền sóng.}
{Trong sóng điện từ, $E$ và $B$ song song với phương truyền.}
\loigiai{
\begin{itemchoice}
\item (Đúng) Khi giải cần xác định đúng dữ kiện, phương chiều, điều kiện áp dụng và đơn vị. (Vì): Khi giải cần xác định đúng dữ kiện, phương chiều, điều kiện áp dụng và đơn vị. b) Sai: Kết quả không phụ thuộc dữ kiện và có thể bỏ qua điều kiện. c) Đúng: Trong chân không, vectơ $E$ và $B$ vuông góc nhau và cùng vuông góc phương truyền sóng. d) Sai: Trong sóng điện từ, $E$ và $B$ song song với phương truyền. Kiến thức cốt lõi: Trong chân không, vectơ $E$ và $B$ vuông góc nhau và cùng vuông góc phương truyền sóng
\item (Sai) Kết quả không phụ thuộc dữ kiện và có thể bỏ qua điều kiện.
\item (Đúng) Trong chân không, vectơ $E$ và $B$ vuông góc nhau và cùng vuông góc phương truyền sóng.
\item (Sai) Trong sóng điện từ, $E$ và $B$ song song với phương truyền.
\end{itemchoice}
}
\end{ex}

% ===== Câu 71 =====
% ID: L12C3B19-13-TL
% Mức: TH
\begin{bt}
% ID: L12C3B19-13-TL
% Mức: TH
Trình bày kiến thức cốt lõi của Nhận biết các tính chất cơ bản của sóng điện từ và nêu điều kiện áp dụng.
\loigiai{
Cần nêu được: Trong chân không, vectơ $E$ và $B$ vuông góc nhau và cùng vuông góc phương truyền sóng. Khi vận dụng phải xác định đúng dữ kiện, phương chiều nếu có, điều kiện áp dụng, hệ đơn vị và kiểm tra tính hợp lí. Nhận định “Trong sóng điện từ, $E$ và $B$ song song với phương truyền.” sai.
}
\end{bt}

% ===== Câu 72 =====
% ID: L12C3B19-13-TLN
% Mức: TH
\begin{ex}
% ID: L12C3B19-13-TLN
% Mức: TH
Trong bốn phát biểu về Nhận biết các tính chất cơ bản của sóng điện từ, có bao nhiêu phát biểu đúng? (Chỉ ghi một số nguyên.) (1) Trong chân không, vectơ $E$ và $B$ vuông góc nhau và cùng vuông góc phương truyền sóng. (2) Trong sóng điện từ, $E$ và $B$ song song với phương truyền. (3) Cần kiểm tra điều kiện áp dụng. (4) Có thể bỏ qua đơn vị.
\shortans{2}
\loigiai{
Đối chiếu với kiến thức: Trong chân không, vectơ $E$ và $B$ vuông góc nhau và cùng vuông góc phương truyền sóng. Có 2 phát biểu đúng.
}
\end{ex}

% ===== Câu 73 =====
% ID: L12C3B19-13-TN
% Mức: NB
\dangbt{Mô tả sự hình thành và lan truyền của sóng điện từ}
\begin{ex}
% ID: L12C3B19-13-TN
% Mức: NB
Tình huống 3: chọn kết luận chính xác về Mô tả sự hình thành và lan truyền của sóng điện từ?
\choice
{Kết quả không phụ thuộc dữ kiện và điều kiện áp dụng.}
{\True Sóng điện từ là điện từ trường lan truyền trong không gian và truyền được trong chân không.}
{Sóng điện từ bắt buộc cần môi trường vật chất để truyền.}
{Có thể bỏ qua phương, chiều và đơn vị trong mọi trường hợp.}
\loigiai{
Sóng điện từ là điện từ trường lan truyền trong không gian và truyền được trong chân không. Vì vậy phương án $B$ đúng; các phương án còn lại sai.
}
\end{ex}

% ===== Câu 74 =====
% ID: L12C3B19-14-DS
% Mức: VD
\dangbt{Nhận biết các tính chất cơ bản của sóng điện từ}
\begin{ex}
% ID: L12C3B19-14-DS
% Mức: VD
Xét Nhận biết các tính chất cơ bản của sóng điện từ (Tình huống 4). Hãy xác định tính đúng, sai của bốn phát biểu sau.
\choiceTF
{Kết quả không phụ thuộc dữ kiện và có thể bỏ qua điều kiện.}
{\True Khi giải cần xác định đúng dữ kiện, phương chiều, điều kiện áp dụng và đơn vị.}
{Trong sóng điện từ, $E$ và $B$ song song với phương truyền.}
{\True Trong chân không, vectơ $E$ và $B$ vuông góc nhau và cùng vuông góc phương truyền sóng.}
\loigiai{
\begin{itemchoice}
\item (Sai) Kết quả không phụ thuộc dữ kiện và có thể bỏ qua điều kiện. (Vì): Kết quả không phụ thuộc dữ kiện và có thể bỏ qua điều kiện. b) Đúng: Khi giải cần xác định đúng dữ kiện, phương chiều, điều kiện áp dụng và đơn vị. c) Sai: Trong sóng điện từ, $E$ và $B$ song song với phương truyền. d) Đúng: Trong chân không, vectơ $E$ và $B$ vuông góc nhau và cùng vuông góc phương truyền sóng. Kiến thức cốt lõi: Trong chân không, vectơ $E$ và $B$ vuông góc nhau và cùng vuông góc phương truyền sóng
\item (Đúng) Khi giải cần xác định đúng dữ kiện, phương chiều, điều kiện áp dụng và đơn vị.
\item (Sai) Trong sóng điện từ, $E$ và $B$ song song với phương truyền.
\item (Đúng) Trong chân không, vectơ $E$ và $B$ vuông góc nhau và cùng vuông góc phương truyền sóng.
\end{itemchoice}
}
\end{ex}

% ===== Câu 75 =====
% ID: L12C3B19-14-TL
% Mức: TH
\begin{bt}
% ID: L12C3B19-14-TL
% Mức: TH
Giải thích vì sao nhận định sau đúng: “Trong chân không, vectơ $E$ và $B$ vuông góc nhau và cùng vuông góc phương truyền sóng.”
\loigiai{
Cần nêu được: Trong chân không, vectơ $E$ và $B$ vuông góc nhau và cùng vuông góc phương truyền sóng. Khi vận dụng phải xác định đúng dữ kiện, phương chiều nếu có, điều kiện áp dụng, hệ đơn vị và kiểm tra tính hợp lí. Nhận định “Trong sóng điện từ, $E$ và $B$ song song với phương truyền.” sai.
}
\end{bt}

% ===== Câu 76 =====
% ID: L12C3B19-14-TLN
% Mức: TH
\begin{ex}
% ID: L12C3B19-14-TLN
% Mức: TH
Trong bốn phát biểu về Nhận biết các tính chất cơ bản của sóng điện từ, có bao nhiêu phát biểu đúng? (Chỉ ghi một số nguyên.) (1) Trong sóng điện từ, $E$ và $B$ song song với phương truyền. (2) Trong chân không, vectơ $E$ và $B$ vuông góc nhau và cùng vuông góc phương truyền sóng. (3) Kết quả phải phù hợp ý nghĩa vật lí. (4) Mọi đại lượng đều là vô hướng.
\shortans{1}
\loigiai{
Đối chiếu với kiến thức: Trong chân không, vectơ $E$ và $B$ vuông góc nhau và cùng vuông góc phương truyền sóng. Có 1 phát biểu đúng.
}
\end{ex}

% ===== Câu 77 =====
% ID: L12C3B19-14-TN
% Mức: TH
\dangbt{Mô tả sự hình thành và lan truyền của sóng điện từ}
\begin{ex}
% ID: L12C3B19-14-TN
% Mức: TH
Nội dung nào mô tả đúng nhất Mô tả sự hình thành và lan truyền của sóng điện từ?
\choice
{Có thể bỏ qua phương, chiều và đơn vị trong mọi trường hợp.}
{Kết quả không phụ thuộc dữ kiện và điều kiện áp dụng.}
{\True Sóng điện từ là điện từ trường lan truyền trong không gian và truyền được trong chân không.}
{Sóng điện từ bắt buộc cần môi trường vật chất để truyền.}
\loigiai{
Sóng điện từ là điện từ trường lan truyền trong không gian và truyền được trong chân không. Vì vậy phương án $C$ đúng; các phương án còn lại sai.
}
\end{ex}

% ===== Câu 78 =====
% ID: L12C3B19-15-DS
% Mức: VD
\dangbt{Nhận biết các tính chất cơ bản của sóng điện từ}
\begin{ex}
% ID: L12C3B19-15-DS
% Mức: VD
Xét Nhận biết các tính chất cơ bản của sóng điện từ (Tình huống 5). Hãy xác định tính đúng, sai của bốn phát biểu sau.
\choiceTF
{\True Trong chân không, vectơ $E$ và $B$ vuông góc nhau và cùng vuông góc phương truyền sóng.}
{Kết quả không phụ thuộc dữ kiện và có thể bỏ qua điều kiện.}
{Trong sóng điện từ, $E$ và $B$ song song với phương truyền.}
{\True Khi giải cần xác định đúng dữ kiện, phương chiều, điều kiện áp dụng và đơn vị.}
\loigiai{
\begin{itemchoice}
\item (Đúng) Trong chân không, vectơ $E$ và $B$ vuông góc nhau và cùng vuông góc phương truyền sóng. (Vì): Trong chân không, vectơ $E$ và $B$ vuông góc nhau và cùng vuông góc phương truyền sóng. b) Sai: Kết quả không phụ thuộc dữ kiện và có thể bỏ qua điều kiện. c) Sai: Trong sóng điện từ, $E$ và $B$ song song với phương truyền. d) Đúng: Khi giải cần xác định đúng dữ kiện, phương chiều, điều kiện áp dụng và đơn vị. Kiến thức cốt lõi: Trong chân không, vectơ $E$ và $B$ vuông góc nhau và cùng vuông góc phương truyền sóng
\item (Sai) Kết quả không phụ thuộc dữ kiện và có thể bỏ qua điều kiện.
\item (Sai) Trong sóng điện từ, $E$ và $B$ song song với phương truyền.
\item (Đúng) Khi giải cần xác định đúng dữ kiện, phương chiều, điều kiện áp dụng và đơn vị.
\end{itemchoice}
}
\end{ex}

% ===== Câu 79 =====
% ID: L12C3B19-15-TL
% Mức: VD
\begin{bt}
% ID: L12C3B19-15-TL
% Mức: VD
Phân tích sai lầm trong nhận định: “Trong sóng điện từ, $E$ và $B$ song song với phương truyền.”
\loigiai{
Cần nêu được: Trong chân không, vectơ $E$ và $B$ vuông góc nhau và cùng vuông góc phương truyền sóng. Khi vận dụng phải xác định đúng dữ kiện, phương chiều nếu có, điều kiện áp dụng, hệ đơn vị và kiểm tra tính hợp lí. Nhận định “Trong sóng điện từ, $E$ và $B$ song song với phương truyền.” sai.
}
\end{bt}

% ===== Câu 80 =====
% ID: L12C3B19-15-TLN
% Mức: TH
\begin{ex}
% ID: L12C3B19-15-TLN
% Mức: TH
Trong bốn phát biểu về Nhận biết các tính chất cơ bản của sóng điện từ, có bao nhiêu phát biểu đúng? (Chỉ ghi một số nguyên.) (1) Trong chân không, vectơ $E$ và $B$ vuông góc nhau và cùng vuông góc phương truyền sóng. (2) Trong sóng điện từ, $E$ và $B$ song song với phương truyền. (3) Cần kiểm tra điều kiện áp dụng. (4) Có thể bỏ qua đơn vị.
\shortans{2}
\loigiai{
Đối chiếu với kiến thức: Trong chân không, vectơ $E$ và $B$ vuông góc nhau và cùng vuông góc phương truyền sóng. Có 2 phát biểu đúng.
}
\end{ex}

% ===== Câu 81 =====
% ID: L12C3B19-15-TN
% Mức: TH
\dangbt{Mô tả sự hình thành và lan truyền của sóng điện từ}
\begin{ex}
% ID: L12C3B19-15-TN
% Mức: TH
Khi phân tích Mô tả sự hình thành và lan truyền của sóng điện từ?
\choice
{Sóng điện từ bắt buộc cần môi trường vật chất để truyền.}
{Có thể bỏ qua phương, chiều và đơn vị trong mọi trường hợp.}
{Kết quả không phụ thuộc dữ kiện và điều kiện áp dụng.}
{\True Sóng điện từ là điện từ trường lan truyền trong không gian và truyền được trong chân không.}
\loigiai{
Sóng điện từ là điện từ trường lan truyền trong không gian và truyền được trong chân không. Vì vậy phương án $D$ đúng; các phương án còn lại sai.
}
\end{ex}

% ===== Câu 82 =====
% ID: L12C3B19-16-DS
% Mức: TH
\dangbt{Nhận biết điện trường biến thiên, từ trường biến thiên và điện từ trường}
\begin{ex}
% ID: L12C3B19-16-DS
% Mức: TH
Xét Nhận biết điện trường biến thiên, từ trường biến thiên và điện từ trường (Tình huống 1). Hãy xác định tính đúng, sai của bốn phát biểu sau.
\choiceTF
{\True Điện trường biến thiên sinh ra từ trường biến thiên và từ trường biến thiên sinh ra điện trường xoáy.}
{Điện trường và từ trường biến thiên tồn tại hoàn toàn độc lập.}
{\True Khi giải cần xác định đúng dữ kiện, phương chiều, điều kiện áp dụng và đơn vị.}
{Kết quả không phụ thuộc dữ kiện và có thể bỏ qua điều kiện.}
\loigiai{
\begin{itemchoice}
\item (Đúng) Điện trường biến thiên sinh ra từ trường biến thiên và từ trường biến thiên sinh ra điện trường xoáy. (Vì): iện trường biến thiên sinh ra từ trường biến thiên và từ trường biến thiên sinh ra điện trường xoáy. b) Sai: Điện trường và từ trường biến thiên tồn tại hoàn toàn độc lập. c) Đúng: Khi giải cần xác định đúng dữ kiện, phương chiều, điều kiện áp dụng và đơn vị. d) Sai: Kết quả không phụ thuộc dữ kiện và có thể bỏ qua điều kiện. Kiến thức cốt lõi: Điện trường biến thiên sinh ra từ trường biến thiên và từ trường biến thiên sinh ra điện trường xoáy
\item (Sai) Điện trường và từ trường biến thiên tồn tại hoàn toàn độc lập.
\item (Đúng) Khi giải cần xác định đúng dữ kiện, phương chiều, điều kiện áp dụng và đơn vị.
\item (Sai) Kết quả không phụ thuộc dữ kiện và có thể bỏ qua điều kiện.
\end{itemchoice}
}
\end{ex}

% ===== Câu 83 =====
% ID: L12C3B19-16-TL
% Mức: VD
\dangbt{Nhận biết các tính chất cơ bản của sóng điện từ}
\begin{bt}
% ID: L12C3B19-16-TL
% Mức: VD
Đề xuất các bước giải một bài toán thuộc dạng Nhận biết các tính chất cơ bản của sóng điện từ.
\loigiai{
Cần nêu được: Trong chân không, vectơ $E$ và $B$ vuông góc nhau và cùng vuông góc phương truyền sóng. Khi vận dụng phải xác định đúng dữ kiện, phương chiều nếu có, điều kiện áp dụng, hệ đơn vị và kiểm tra tính hợp lí. Nhận định “Trong sóng điện từ, $E$ và $B$ song song với phương truyền.” sai.
}
\end{bt}

% ===== Câu 84 =====
% ID: L12C3B19-16-TLN
% Mức: VD
\begin{ex}
% ID: L12C3B19-16-TLN
% Mức: VD
Trong bốn phát biểu về Nhận biết các tính chất cơ bản của sóng điện từ, có bao nhiêu phát biểu đúng? (Chỉ ghi một số nguyên.) (1) Trong sóng điện từ, $E$ và $B$ song song với phương truyền. (2) Trong chân không, vectơ $E$ và $B$ vuông góc nhau và cùng vuông góc phương truyền sóng. (3) Kết quả phải phù hợp ý nghĩa vật lí. (4) Mọi đại lượng đều là vô hướng.
\shortans{2}
\loigiai{
Đối chiếu với kiến thức: Trong chân không, vectơ $E$ và $B$ vuông góc nhau và cùng vuông góc phương truyền sóng. Có 2 phát biểu đúng.
}
\end{ex}

% ===== Câu 85 =====
% ID: L12C3B19-16-TN
% Mức: TH
\dangbt{Mô tả sự hình thành và lan truyền của sóng điện từ}
\begin{ex}
% ID: L12C3B19-16-TN
% Mức: TH
Một học sinh ôn tập Mô tả sự hình thành và lan truyền của sóng điện từ?
\choice
{\True Sóng điện từ là điện từ trường lan truyền trong không gian và truyền được trong chân không.}
{Sóng điện từ bắt buộc cần môi trường vật chất để truyền.}
{Có thể bỏ qua phương, chiều và đơn vị trong mọi trường hợp.}
{Kết quả không phụ thuộc dữ kiện và điều kiện áp dụng.}
\loigiai{
Sóng điện từ là điện từ trường lan truyền trong không gian và truyền được trong chân không. Vì vậy phương án $A$ đúng; các phương án còn lại sai.
}
\end{ex}

% ===== Câu 86 =====
% ID: L12C3B19-17-DS
% Mức: TH
\dangbt{Nhận biết điện trường biến thiên, từ trường biến thiên và điện từ trường}
\begin{ex}
% ID: L12C3B19-17-DS
% Mức: TH
Xét Nhận biết điện trường biến thiên, từ trường biến thiên và điện từ trường (Tình huống 2). Hãy xác định tính đúng, sai của bốn phát biểu sau.
\choiceTF
{Điện trường và từ trường biến thiên tồn tại hoàn toàn độc lập.}
{\True Điện trường biến thiên sinh ra từ trường biến thiên và từ trường biến thiên sinh ra điện trường xoáy.}
{Kết quả không phụ thuộc dữ kiện và có thể bỏ qua điều kiện.}
{\True Khi giải cần xác định đúng dữ kiện, phương chiều, điều kiện áp dụng và đơn vị.}
\loigiai{
\begin{itemchoice}
\item (Sai) Điện trường và từ trường biến thiên tồn tại hoàn toàn độc lập. (Vì): Điện trường và từ trường biến thiên tồn tại hoàn toàn độc lập. b) Đúng: Điện trường biến thiên sinh ra từ trường biến thiên và từ trường biến thiên sinh ra điện trường xoáy. c) Sai: Kết quả không phụ thuộc dữ kiện và có thể bỏ qua điều kiện. d) Đúng: Khi giải cần xác định đúng dữ kiện, phương chiều, điều kiện áp dụng và đơn vị. Kiến thức cốt lõi: Điện trường biến thiên sinh ra từ trường biến thiên và từ trường biến thiên sinh ra điện trường xoáy
\item (Đúng) Điện trường biến thiên sinh ra từ trường biến thiên và từ trường biến thiên sinh ra điện trường xoáy.
\item (Sai) Kết quả không phụ thuộc dữ kiện và có thể bỏ qua điều kiện.
\item (Đúng) Khi giải cần xác định đúng dữ kiện, phương chiều, điều kiện áp dụng và đơn vị.
\end{itemchoice}
}
\end{ex}

% ===== Câu 87 =====
% ID: L12C3B19-17-TL
% Mức: VD
\dangbt{Nhận biết các tính chất cơ bản của sóng điện từ}
\begin{bt}
% ID: L12C3B19-17-TL
% Mức: VD
Nêu một tình huống thực tiễn liên quan đến Nhận biết các tính chất cơ bản của sóng điện từ và giải thích bằng kiến thức Vật lí.
\loigiai{
Cần nêu được: Trong chân không, vectơ $E$ và $B$ vuông góc nhau và cùng vuông góc phương truyền sóng. Khi vận dụng phải xác định đúng dữ kiện, phương chiều nếu có, điều kiện áp dụng, hệ đơn vị và kiểm tra tính hợp lí. Nhận định “Trong sóng điện từ, $E$ và $B$ song song với phương truyền.” sai.
}
\end{bt}

% ===== Câu 88 =====
% ID: L12C3B19-17-TLN
% Mức: VD
\begin{ex}
% ID: L12C3B19-17-TLN
% Mức: VD
Trong bốn phát biểu về Nhận biết các tính chất cơ bản của sóng điện từ, có bao nhiêu phát biểu đúng? (Chỉ ghi một số nguyên.) (1) Trong chân không, vectơ $E$ và $B$ vuông góc nhau và cùng vuông góc phương truyền sóng. (2) Trong sóng điện từ, $E$ và $B$ song song với phương truyền. (3) Cần kiểm tra điều kiện áp dụng. (4) Có thể bỏ qua đơn vị.
\shortans{2}
\loigiai{
Đối chiếu với kiến thức: Trong chân không, vectơ $E$ và $B$ vuông góc nhau và cùng vuông góc phương truyền sóng. Có 2 phát biểu đúng.
}
\end{ex}

% ===== Câu 89 =====
% ID: L12C3B19-17-TN
% Mức: TH
\dangbt{Mô tả sự hình thành và lan truyền của sóng điện từ}
\begin{ex}
% ID: L12C3B19-17-TN
% Mức: TH
Chọn phát biểu không trái với kiến thức về Mô tả sự hình thành và lan truyền của sóng điện từ?
\choice
{Kết quả không phụ thuộc dữ kiện và điều kiện áp dụng.}
{\True Sóng điện từ là điện từ trường lan truyền trong không gian và truyền được trong chân không.}
{Sóng điện từ bắt buộc cần môi trường vật chất để truyền.}
{Có thể bỏ qua phương, chiều và đơn vị trong mọi trường hợp.}
\loigiai{
Sóng điện từ là điện từ trường lan truyền trong không gian và truyền được trong chân không. Vì vậy phương án $B$ đúng; các phương án còn lại sai.
}
\end{ex}

% ===== Câu 90 =====
% ID: L12C3B19-18-DS
% Mức: VD
\dangbt{Nhận biết điện trường biến thiên, từ trường biến thiên và điện từ trường}
\begin{ex}
% ID: L12C3B19-18-DS
% Mức: VD
Xét Nhận biết điện trường biến thiên, từ trường biến thiên và điện từ trường (Tình huống 3). Hãy xác định tính đúng, sai của bốn phát biểu sau.
\choiceTF
{\True Khi giải cần xác định đúng dữ kiện, phương chiều, điều kiện áp dụng và đơn vị.}
{Kết quả không phụ thuộc dữ kiện và có thể bỏ qua điều kiện.}
{\True Điện trường biến thiên sinh ra từ trường biến thiên và từ trường biến thiên sinh ra điện trường xoáy.}
{Điện trường và từ trường biến thiên tồn tại hoàn toàn độc lập.}
\loigiai{
\begin{itemchoice}
\item (Đúng) Khi giải cần xác định đúng dữ kiện, phương chiều, điều kiện áp dụng và đơn vị. (Vì): Khi giải cần xác định đúng dữ kiện, phương chiều, điều kiện áp dụng và đơn vị. b) Sai: Kết quả không phụ thuộc dữ kiện và có thể bỏ qua điều kiện. c) Đúng: Điện trường biến thiên sinh ra từ trường biến thiên và từ trường biến thiên sinh ra điện trường xoáy. d) Sai: Điện trường và từ trường biến thiên tồn tại hoàn toàn độc lập. Kiến thức cốt lõi: Điện trường biến thiên sinh ra từ trường biến thiên và từ trường biến thiên sinh ra điện trường xoáy
\item (Sai) Kết quả không phụ thuộc dữ kiện và có thể bỏ qua điều kiện.
\item (Đúng) Điện trường biến thiên sinh ra từ trường biến thiên và từ trường biến thiên sinh ra điện trường xoáy.
\item (Sai) Điện trường và từ trường biến thiên tồn tại hoàn toàn độc lập.
\end{itemchoice}
}
\end{ex}

% ===== Câu 91 =====
% ID: L12C3B19-18-TL
% Mức: VDC
\dangbt{Nhận biết các tính chất cơ bản của sóng điện từ}
\begin{bt}
% ID: L12C3B19-18-TL
% Mức: VDC
So sánh nhận định đúng “Trong chân không, vectơ $E$ và $B$ vuông góc nhau và cùng vuông góc phương truyền sóng.” với nhận định sai “Trong sóng điện từ, $E$ và $B$ song song với phương truyền.”; chỉ rõ căn cứ Vật lí.
\loigiai{
Cần nêu được: Trong chân không, vectơ $E$ và $B$ vuông góc nhau và cùng vuông góc phương truyền sóng. Khi vận dụng phải xác định đúng dữ kiện, phương chiều nếu có, điều kiện áp dụng, hệ đơn vị và kiểm tra tính hợp lí. Nhận định “Trong sóng điện từ, $E$ và $B$ song song với phương truyền.” sai.
}
\end{bt}

% ===== Câu 92 =====
% ID: L12C3B19-18-TLN
% Mức: VDC
\begin{ex}
% ID: L12C3B19-18-TLN
% Mức: VDC
Trong bốn phát biểu về Nhận biết các tính chất cơ bản của sóng điện từ, có bao nhiêu phát biểu đúng? (Chỉ ghi một số nguyên.) (1) Khi giải cần xác định đúng dữ kiện, phương chiều, điều kiện áp dụng và đơn vị. (2) Trong sóng điện từ, $E$ và $B$ song song với phương truyền. (3) Phải dùng đúng định luật cho tình huống. (4) Trong chân không, vectơ $E$ và $B$ vuông góc nhau và cùng vuông góc phương truyền sóng.
\shortans{3}
\loigiai{
Đối chiếu với kiến thức: Trong chân không, vectơ $E$ và $B$ vuông góc nhau và cùng vuông góc phương truyền sóng. Có 3 phát biểu đúng.
}
\end{ex}

% ===== Câu 93 =====
% ID: L12C3B19-18-TN
% Mức: VD
\dangbt{Mô tả sự hình thành và lan truyền của sóng điện từ}
\begin{ex}
% ID: L12C3B19-18-TN
% Mức: VD
Cơ sở Vật lí đúng dùng để giải Mô tả sự hình thành và lan truyền của sóng điện từ?
\choice
{Có thể bỏ qua phương, chiều và đơn vị trong mọi trường hợp.}
{Kết quả không phụ thuộc dữ kiện và điều kiện áp dụng.}
{\True Sóng điện từ là điện từ trường lan truyền trong không gian và truyền được trong chân không.}
{Sóng điện từ bắt buộc cần môi trường vật chất để truyền.}
\loigiai{
Sóng điện từ là điện từ trường lan truyền trong không gian và truyền được trong chân không. Vì vậy phương án $C$ đúng; các phương án còn lại sai.
}
\end{ex}

% ===== Câu 94 =====
% ID: L12C3B19-19-DS
% Mức: VD
\dangbt{Nhận biết điện trường biến thiên, từ trường biến thiên và điện từ trường}
\begin{ex}
% ID: L12C3B19-19-DS
% Mức: VD
Xét Nhận biết điện trường biến thiên, từ trường biến thiên và điện từ trường (Tình huống 4). Hãy xác định tính đúng, sai của bốn phát biểu sau.
\choiceTF
{Kết quả không phụ thuộc dữ kiện và có thể bỏ qua điều kiện.}
{\True Khi giải cần xác định đúng dữ kiện, phương chiều, điều kiện áp dụng và đơn vị.}
{Điện trường và từ trường biến thiên tồn tại hoàn toàn độc lập.}
{\True Điện trường biến thiên sinh ra từ trường biến thiên và từ trường biến thiên sinh ra điện trường xoáy.}
\loigiai{
\begin{itemchoice}
\item (Sai) Kết quả không phụ thuộc dữ kiện và có thể bỏ qua điều kiện. (Vì): Kết quả không phụ thuộc dữ kiện và có thể bỏ qua điều kiện. b) Đúng: Khi giải cần xác định đúng dữ kiện, phương chiều, điều kiện áp dụng và đơn vị. c) Sai: Điện trường và từ trường biến thiên tồn tại hoàn toàn độc lập. d) Đúng: Điện trường biến thiên sinh ra từ trường biến thiên và từ trường biến thiên sinh ra điện trường xoáy. Kiến thức cốt lõi: Điện trường biến thiên sinh ra từ trường biến thiên và từ trường biến thiên sinh ra điện trường xoáy
\item (Đúng) Khi giải cần xác định đúng dữ kiện, phương chiều, điều kiện áp dụng và đơn vị.
\item (Sai) Điện trường và từ trường biến thiên tồn tại hoàn toàn độc lập.
\item (Đúng) Điện trường biến thiên sinh ra từ trường biến thiên và từ trường biến thiên sinh ra điện trường xoáy.
\end{itemchoice}
}
\end{ex}

% ===== Câu 95 =====
% ID: L12C3B19-19-TL
% Mức: TH
\begin{bt}
% ID: L12C3B19-19-TL
% Mức: TH
Trình bày kiến thức cốt lõi của Nhận biết điện trường biến thiên, từ trường biến thiên và điện từ trường và nêu điều kiện áp dụng.
\loigiai{
Cần nêu được: Điện trường biến thiên sinh ra từ trường biến thiên và từ trường biến thiên sinh ra điện trường xoáy. Khi vận dụng phải xác định đúng dữ kiện, phương chiều nếu có, điều kiện áp dụng, hệ đơn vị và kiểm tra tính hợp lí. Nhận định “Điện trường và từ trường biến thiên tồn tại hoàn toàn độc lập.” sai.
}
\end{bt}

% ===== Câu 96 =====
% ID: L12C3B19-19-TLN
% Mức: TH
\begin{ex}
% ID: L12C3B19-19-TLN
% Mức: TH
Trong bốn phát biểu về Nhận biết điện trường biến thiên, từ trường biến thiên và điện từ trường, có bao nhiêu phát biểu đúng? (Chỉ ghi một số nguyên.) (1) Điện trường biến thiên sinh ra từ trường biến thiên và từ trường biến thiên sinh ra điện trường xoáy. (2) Điện trường và từ trường biến thiên tồn tại hoàn toàn độc lập. (3) Cần kiểm tra điều kiện áp dụng. (4) Có thể bỏ qua đơn vị.
\shortans{2}
\loigiai{
Đối chiếu với kiến thức: Điện trường biến thiên sinh ra từ trường biến thiên và từ trường biến thiên sinh ra điện trường xoáy. Có 2 phát biểu đúng.
}
\end{ex}

% ===== Câu 97 =====
% ID: L12C3B19-19-TN
% Mức: VD
\dangbt{Mô tả sự hình thành và lan truyền của sóng điện từ}
\begin{ex}
% ID: L12C3B19-19-TN
% Mức: VD
Trong một bài thực hành về Mô tả sự hình thành và lan truyền của sóng điện từ?
\choice
{Sóng điện từ bắt buộc cần môi trường vật chất để truyền.}
{Có thể bỏ qua phương, chiều và đơn vị trong mọi trường hợp.}
{Kết quả không phụ thuộc dữ kiện và điều kiện áp dụng.}
{\True Sóng điện từ là điện từ trường lan truyền trong không gian và truyền được trong chân không.}
\loigiai{
Sóng điện từ là điện từ trường lan truyền trong không gian và truyền được trong chân không. Vì vậy phương án $D$ đúng; các phương án còn lại sai.
}
\end{ex}

% ===== Câu 98 =====
% ID: L12C3B19-20-DS
% Mức: VD
\dangbt{Nhận biết điện trường biến thiên, từ trường biến thiên và điện từ trường}
\begin{ex}
% ID: L12C3B19-20-DS
% Mức: VD
Xét Nhận biết điện trường biến thiên, từ trường biến thiên và điện từ trường (Tình huống 5). Hãy xác định tính đúng, sai của bốn phát biểu sau.
\choiceTF
{\True Điện trường biến thiên sinh ra từ trường biến thiên và từ trường biến thiên sinh ra điện trường xoáy.}
{Kết quả không phụ thuộc dữ kiện và có thể bỏ qua điều kiện.}
{Điện trường và từ trường biến thiên tồn tại hoàn toàn độc lập.}
{\True Khi giải cần xác định đúng dữ kiện, phương chiều, điều kiện áp dụng và đơn vị.}
\loigiai{
\begin{itemchoice}
\item (Đúng) Điện trường biến thiên sinh ra từ trường biến thiên và từ trường biến thiên sinh ra điện trường xoáy. (Vì): iện trường biến thiên sinh ra từ trường biến thiên và từ trường biến thiên sinh ra điện trường xoáy. b) Sai: Kết quả không phụ thuộc dữ kiện và có thể bỏ qua điều kiện. c) Sai: Điện trường và từ trường biến thiên tồn tại hoàn toàn độc lập. d) Đúng: Khi giải cần xác định đúng dữ kiện, phương chiều, điều kiện áp dụng và đơn vị. Kiến thức cốt lõi: Điện trường biến thiên sinh ra từ trường biến thiên và từ trường biến thiên sinh ra điện trường xoáy
\item (Sai) Kết quả không phụ thuộc dữ kiện và có thể bỏ qua điều kiện.
\item (Sai) Điện trường và từ trường biến thiên tồn tại hoàn toàn độc lập.
\item (Đúng) Khi giải cần xác định đúng dữ kiện, phương chiều, điều kiện áp dụng và đơn vị.
\end{itemchoice}
}
\end{ex}

% ===== Câu 99 =====
% ID: L12C3B19-20-TL
% Mức: TH
\begin{bt}
% ID: L12C3B19-20-TL
% Mức: TH
Giải thích vì sao nhận định sau đúng: “Điện trường biến thiên sinh ra từ trường biến thiên và từ trường biến thiên sinh ra điện trường xoáy.”
\loigiai{
Cần nêu được: Điện trường biến thiên sinh ra từ trường biến thiên và từ trường biến thiên sinh ra điện trường xoáy. Khi vận dụng phải xác định đúng dữ kiện, phương chiều nếu có, điều kiện áp dụng, hệ đơn vị và kiểm tra tính hợp lí. Nhận định “Điện trường và từ trường biến thiên tồn tại hoàn toàn độc lập.” sai.
}
\end{bt}

% ===== Câu 100 =====
% ID: L12C3B19-20-TLN
% Mức: TH
\begin{ex}
% ID: L12C3B19-20-TLN
% Mức: TH
Trong bốn phát biểu về Nhận biết điện trường biến thiên, từ trường biến thiên và điện từ trường, có bao nhiêu phát biểu đúng? (Chỉ ghi một số nguyên.) (1) Điện trường và từ trường biến thiên tồn tại hoàn toàn độc lập. (2) Điện trường biến thiên sinh ra từ trường biến thiên và từ trường biến thiên sinh ra điện trường xoáy. (3) Kết quả phải phù hợp ý nghĩa vật lí. (4) Mọi đại lượng đều là vô hướng.
\shortans{1}
\loigiai{
Đối chiếu với kiến thức: Điện trường biến thiên sinh ra từ trường biến thiên và từ trường biến thiên sinh ra điện trường xoáy. Có 1 phát biểu đúng.
}
\end{ex}

% ===== Câu 101 =====
% ID: L12C3B19-20-TN
% Mức: VD
\dangbt{Mô tả sự hình thành và lan truyền của sóng điện từ}
\begin{ex}
% ID: L12C3B19-20-TN
% Mức: VD
Kết luận đúng khi vận dụng Mô tả sự hình thành và lan truyền của sóng điện từ?
\choice
{\True Sóng điện từ là điện từ trường lan truyền trong không gian và truyền được trong chân không.}
{Sóng điện từ bắt buộc cần môi trường vật chất để truyền.}
{Có thể bỏ qua phương, chiều và đơn vị trong mọi trường hợp.}
{Kết quả không phụ thuộc dữ kiện và điều kiện áp dụng.}
\loigiai{
Sóng điện từ là điện từ trường lan truyền trong không gian và truyền được trong chân không. Vì vậy phương án $A$ đúng; các phương án còn lại sai.
}
\end{ex}

% ===== Câu 102 =====
% ID: L12C3B19-21-DS
% Mức: TH
\dangbt{Phân loại phổ sóng điện từ và nêu ứng dụng}
\begin{ex}
% ID: L12C3B19-21-DS
% Mức: TH
Xét Phân loại phổ sóng điện từ và nêu ứng dụng (Tình huống 1). Hãy xác định tính đúng, sai của bốn phát biểu sau.
\choiceTF
{\True Phổ điện từ gồm sóng vô tuyến, vi ba, hồng ngoại, ánh sáng nhìn thấy, tử ngoại, tia $X$ và gamma.}
{Âm thanh là một miền của phổ sóng điện từ.}
{\True Khi giải cần xác định đúng dữ kiện, phương chiều, điều kiện áp dụng và đơn vị.}
{Kết quả không phụ thuộc dữ kiện và có thể bỏ qua điều kiện.}
\loigiai{
\begin{itemchoice}
\item (Đúng) Phổ điện từ gồm sóng vô tuyến, vi ba, hồng ngoại, ánh sáng nhìn thấy, tử ngoại, tia $X$ và gamma. (Vì): Phổ điện từ gồm sóng vô tuyến, vi ba, hồng ngoại, ánh sáng nhìn thấy, tử ngoại, tia $X$ và gamma. b) Sai: Âm thanh là một miền của phổ sóng điện từ. c) Đúng: Khi giải cần xác định đúng dữ kiện, phương chiều, điều kiện áp dụng và đơn vị. d) Sai: Kết quả không phụ thuộc dữ kiện và có thể bỏ qua điều kiện. Kiến thức cốt lõi: Phổ điện từ gồm sóng vô tuyến, vi ba, hồng ngoại, ánh sáng nhìn thấy, tử ngoại, tia $X$ và gamma
\item (Sai) Âm thanh là một miền của phổ sóng điện từ.
\item (Đúng) Khi giải cần xác định đúng dữ kiện, phương chiều, điều kiện áp dụng và đơn vị.
\item (Sai) Kết quả không phụ thuộc dữ kiện và có thể bỏ qua điều kiện.
\end{itemchoice}
}
\end{ex}

% ===== Câu 103 =====
% ID: L12C3B19-21-TL
% Mức: VD
\dangbt{Nhận biết điện trường biến thiên, từ trường biến thiên và điện từ trường}
\begin{bt}
% ID: L12C3B19-21-TL
% Mức: VD
Phân tích sai lầm trong nhận định: “Điện trường và từ trường biến thiên tồn tại hoàn toàn độc lập.”
\loigiai{
Cần nêu được: Điện trường biến thiên sinh ra từ trường biến thiên và từ trường biến thiên sinh ra điện trường xoáy. Khi vận dụng phải xác định đúng dữ kiện, phương chiều nếu có, điều kiện áp dụng, hệ đơn vị và kiểm tra tính hợp lí. Nhận định “Điện trường và từ trường biến thiên tồn tại hoàn toàn độc lập.” sai.
}
\end{bt}

% ===== Câu 104 =====
% ID: L12C3B19-21-TLN
% Mức: TH
\begin{ex}
% ID: L12C3B19-21-TLN
% Mức: TH
Trong bốn phát biểu về Nhận biết điện trường biến thiên, từ trường biến thiên và điện từ trường, có bao nhiêu phát biểu đúng? (Chỉ ghi một số nguyên.) (1) Điện trường biến thiên sinh ra từ trường biến thiên và từ trường biến thiên sinh ra điện trường xoáy. (2) Điện trường và từ trường biến thiên tồn tại hoàn toàn độc lập. (3) Cần kiểm tra điều kiện áp dụng. (4) Có thể bỏ qua đơn vị.
\shortans{2}
\loigiai{
Đối chiếu với kiến thức: Điện trường biến thiên sinh ra từ trường biến thiên và từ trường biến thiên sinh ra điện trường xoáy. Có 2 phát biểu đúng.
}
\end{ex}

% ===== Câu 105 =====
% ID: L12C3B19-21-TN
% Mức: NB
\dangbt{Nhận biết các tính chất cơ bản của sóng điện từ}
\begin{ex}
% ID: L12C3B19-21-TN
% Mức: NB
Phát biểu nào sau đây đúng về Nhận biết các tính chất cơ bản của sóng điện từ?
\choice
{Trong sóng điện từ, $E$ và $B$ song song với phương truyền.}
{Có thể bỏ qua phương, chiều và đơn vị trong mọi trường hợp.}
{Kết quả không phụ thuộc dữ kiện và điều kiện áp dụng.}
{\True Trong chân không, vectơ $E$ và $B$ vuông góc nhau và cùng vuông góc phương truyền sóng.}
\loigiai{
Trong chân không, vectơ $E$ và $B$ vuông góc nhau và cùng vuông góc phương truyền sóng. Vì vậy phương án $D$ đúng; các phương án còn lại sai.
}
\end{ex}

% ===== Câu 106 =====
% ID: L12C3B19-22-DS
% Mức: TH
\dangbt{Phân loại phổ sóng điện từ và nêu ứng dụng}
\begin{ex}
% ID: L12C3B19-22-DS
% Mức: TH
Xét Phân loại phổ sóng điện từ và nêu ứng dụng (Tình huống 2). Hãy xác định tính đúng, sai của bốn phát biểu sau.
\choiceTF
{Âm thanh là một miền của phổ sóng điện từ.}
{\True Phổ điện từ gồm sóng vô tuyến, vi ba, hồng ngoại, ánh sáng nhìn thấy, tử ngoại, tia $X$ và gamma.}
{Kết quả không phụ thuộc dữ kiện và có thể bỏ qua điều kiện.}
{\True Khi giải cần xác định đúng dữ kiện, phương chiều, điều kiện áp dụng và đơn vị.}
\loigiai{
\begin{itemchoice}
\item (Sai) Âm thanh là một miền của phổ sóng điện từ. (Vì): Âm thanh là một miền của phổ sóng điện từ. b) Đúng: Phổ điện từ gồm sóng vô tuyến, vi ba, hồng ngoại, ánh sáng nhìn thấy, tử ngoại, tia $X$ và gamma. c) Sai: Kết quả không phụ thuộc dữ kiện và có thể bỏ qua điều kiện. d) Đúng: Khi giải cần xác định đúng dữ kiện, phương chiều, điều kiện áp dụng và đơn vị. Kiến thức cốt lõi: Phổ điện từ gồm sóng vô tuyến, vi ba, hồng ngoại, ánh sáng nhìn thấy, tử ngoại, tia $X$ và gamma
\item (Đúng) Phổ điện từ gồm sóng vô tuyến, vi ba, hồng ngoại, ánh sáng nhìn thấy, tử ngoại, tia $X$ và gamma.
\item (Sai) Kết quả không phụ thuộc dữ kiện và có thể bỏ qua điều kiện.
\item (Đúng) Khi giải cần xác định đúng dữ kiện, phương chiều, điều kiện áp dụng và đơn vị.
\end{itemchoice}
}
\end{ex}

% ===== Câu 107 =====
% ID: L12C3B19-22-TL
% Mức: VD
\dangbt{Nhận biết điện trường biến thiên, từ trường biến thiên và điện từ trường}
\begin{bt}
% ID: L12C3B19-22-TL
% Mức: VD
Đề xuất các bước giải một bài toán thuộc dạng Nhận biết điện trường biến thiên, từ trường biến thiên và điện từ trường.
\loigiai{
Cần nêu được: Điện trường biến thiên sinh ra từ trường biến thiên và từ trường biến thiên sinh ra điện trường xoáy. Khi vận dụng phải xác định đúng dữ kiện, phương chiều nếu có, điều kiện áp dụng, hệ đơn vị và kiểm tra tính hợp lí. Nhận định “Điện trường và từ trường biến thiên tồn tại hoàn toàn độc lập.” sai.
}
\end{bt}

% ===== Câu 108 =====
% ID: L12C3B19-22-TLN
% Mức: VD
\begin{ex}
% ID: L12C3B19-22-TLN
% Mức: VD
Trong bốn phát biểu về Nhận biết điện trường biến thiên, từ trường biến thiên và điện từ trường, có bao nhiêu phát biểu đúng? (Chỉ ghi một số nguyên.) (1) Điện trường và từ trường biến thiên tồn tại hoàn toàn độc lập. (2) Điện trường biến thiên sinh ra từ trường biến thiên và từ trường biến thiên sinh ra điện trường xoáy. (3) Kết quả phải phù hợp ý nghĩa vật lí. (4) Mọi đại lượng đều là vô hướng.
\shortans{2}
\loigiai{
Đối chiếu với kiến thức: Điện trường biến thiên sinh ra từ trường biến thiên và từ trường biến thiên sinh ra điện trường xoáy. Có 2 phát biểu đúng.
}
\end{ex}

% ===== Câu 109 =====
% ID: L12C3B19-22-TN
% Mức: NB
\dangbt{Nhận biết các tính chất cơ bản của sóng điện từ}
\begin{ex}
% ID: L12C3B19-22-TN
% Mức: NB
Trong nội dung Nhận biết các tính chất cơ bản của sóng điện từ?
\choice
{\True Trong chân không, vectơ $E$ và $B$ vuông góc nhau và cùng vuông góc phương truyền sóng.}
{Trong sóng điện từ, $E$ và $B$ song song với phương truyền.}
{Có thể bỏ qua phương, chiều và đơn vị trong mọi trường hợp.}
{Kết quả không phụ thuộc dữ kiện và điều kiện áp dụng.}
\loigiai{
Trong chân không, vectơ $E$ và $B$ vuông góc nhau và cùng vuông góc phương truyền sóng. Vì vậy phương án $A$ đúng; các phương án còn lại sai.
}
\end{ex}

% ===== Câu 110 =====
% ID: L12C3B19-23-DS
% Mức: VD
\dangbt{Phân loại phổ sóng điện từ và nêu ứng dụng}
\begin{ex}
% ID: L12C3B19-23-DS
% Mức: VD
Xét Phân loại phổ sóng điện từ và nêu ứng dụng (Tình huống 3). Hãy xác định tính đúng, sai của bốn phát biểu sau.
\choiceTF
{\True Khi giải cần xác định đúng dữ kiện, phương chiều, điều kiện áp dụng và đơn vị.}
{Kết quả không phụ thuộc dữ kiện và có thể bỏ qua điều kiện.}
{\True Phổ điện từ gồm sóng vô tuyến, vi ba, hồng ngoại, ánh sáng nhìn thấy, tử ngoại, tia $X$ và gamma.}
{Âm thanh là một miền của phổ sóng điện từ.}
\loigiai{
\begin{itemchoice}
\item (Đúng) Khi giải cần xác định đúng dữ kiện, phương chiều, điều kiện áp dụng và đơn vị. (Vì): Khi giải cần xác định đúng dữ kiện, phương chiều, điều kiện áp dụng và đơn vị. b) Sai: Kết quả không phụ thuộc dữ kiện và có thể bỏ qua điều kiện. c) Đúng: Phổ điện từ gồm sóng vô tuyến, vi ba, hồng ngoại, ánh sáng nhìn thấy, tử ngoại, tia $X$ và gamma. d) Sai: Âm thanh là một miền của phổ sóng điện từ. Kiến thức cốt lõi: Phổ điện từ gồm sóng vô tuyến, vi ba, hồng ngoại, ánh sáng nhìn thấy, tử ngoại, tia $X$ và gamma
\item (Sai) Kết quả không phụ thuộc dữ kiện và có thể bỏ qua điều kiện.
\item (Đúng) Phổ điện từ gồm sóng vô tuyến, vi ba, hồng ngoại, ánh sáng nhìn thấy, tử ngoại, tia $X$ và gamma.
\item (Sai) Âm thanh là một miền của phổ sóng điện từ.
\end{itemchoice}
}
\end{ex}

% ===== Câu 111 =====
% ID: L12C3B19-23-TL
% Mức: VD
\dangbt{Nhận biết điện trường biến thiên, từ trường biến thiên và điện từ trường}
\begin{bt}
% ID: L12C3B19-23-TL
% Mức: VD
Nêu một tình huống thực tiễn liên quan đến Nhận biết điện trường biến thiên, từ trường biến thiên và điện từ trường và giải thích bằng kiến thức Vật lí.
\loigiai{
Cần nêu được: Điện trường biến thiên sinh ra từ trường biến thiên và từ trường biến thiên sinh ra điện trường xoáy. Khi vận dụng phải xác định đúng dữ kiện, phương chiều nếu có, điều kiện áp dụng, hệ đơn vị và kiểm tra tính hợp lí. Nhận định “Điện trường và từ trường biến thiên tồn tại hoàn toàn độc lập.” sai.
}
\end{bt}

% ===== Câu 112 =====
% ID: L12C3B19-23-TLN
% Mức: VD
\begin{ex}
% ID: L12C3B19-23-TLN
% Mức: VD
Trong bốn phát biểu về Nhận biết điện trường biến thiên, từ trường biến thiên và điện từ trường, có bao nhiêu phát biểu đúng? (Chỉ ghi một số nguyên.) (1) Điện trường biến thiên sinh ra từ trường biến thiên và từ trường biến thiên sinh ra điện trường xoáy. (2) Điện trường và từ trường biến thiên tồn tại hoàn toàn độc lập. (3) Cần kiểm tra điều kiện áp dụng. (4) Có thể bỏ qua đơn vị.
\shortans{2}
\loigiai{
Đối chiếu với kiến thức: Điện trường biến thiên sinh ra từ trường biến thiên và từ trường biến thiên sinh ra điện trường xoáy. Có 2 phát biểu đúng.
}
\end{ex}

% ===== Câu 113 =====
% ID: L12C3B19-23-TN
% Mức: NB
\dangbt{Nhận biết các tính chất cơ bản của sóng điện từ}
\begin{ex}
% ID: L12C3B19-23-TN
% Mức: NB
Tình huống 3: chọn kết luận chính xác về Nhận biết các tính chất cơ bản của sóng điện từ?
\choice
{Kết quả không phụ thuộc dữ kiện và điều kiện áp dụng.}
{\True Trong chân không, vectơ $E$ và $B$ vuông góc nhau và cùng vuông góc phương truyền sóng.}
{Trong sóng điện từ, $E$ và $B$ song song với phương truyền.}
{Có thể bỏ qua phương, chiều và đơn vị trong mọi trường hợp.}
\loigiai{
Trong chân không, vectơ $E$ và $B$ vuông góc nhau và cùng vuông góc phương truyền sóng. Vì vậy phương án $B$ đúng; các phương án còn lại sai.
}
\end{ex}

% ===== Câu 114 =====
% ID: L12C3B19-24-DS
% Mức: VD
\dangbt{Phân loại phổ sóng điện từ và nêu ứng dụng}
\begin{ex}
% ID: L12C3B19-24-DS
% Mức: VD
Xét Phân loại phổ sóng điện từ và nêu ứng dụng (Tình huống 4). Hãy xác định tính đúng, sai của bốn phát biểu sau.
\choiceTF
{Kết quả không phụ thuộc dữ kiện và có thể bỏ qua điều kiện.}
{\True Khi giải cần xác định đúng dữ kiện, phương chiều, điều kiện áp dụng và đơn vị.}
{Âm thanh là một miền của phổ sóng điện từ.}
{\True Phổ điện từ gồm sóng vô tuyến, vi ba, hồng ngoại, ánh sáng nhìn thấy, tử ngoại, tia $X$ và gamma.}
\loigiai{
\begin{itemchoice}
\item (Sai) Kết quả không phụ thuộc dữ kiện và có thể bỏ qua điều kiện. (Vì): Kết quả không phụ thuộc dữ kiện và có thể bỏ qua điều kiện. b) Đúng: Khi giải cần xác định đúng dữ kiện, phương chiều, điều kiện áp dụng và đơn vị. c) Sai: Âm thanh là một miền của phổ sóng điện từ. d) Đúng: Phổ điện từ gồm sóng vô tuyến, vi ba, hồng ngoại, ánh sáng nhìn thấy, tử ngoại, tia $X$ và gamma. Kiến thức cốt lõi: Phổ điện từ gồm sóng vô tuyến, vi ba, hồng ngoại, ánh sáng nhìn thấy, tử ngoại, tia $X$ và gamma
\item (Đúng) Khi giải cần xác định đúng dữ kiện, phương chiều, điều kiện áp dụng và đơn vị.
\item (Sai) Âm thanh là một miền của phổ sóng điện từ.
\item (Đúng) Phổ điện từ gồm sóng vô tuyến, vi ba, hồng ngoại, ánh sáng nhìn thấy, tử ngoại, tia $X$ và gamma.
\end{itemchoice}
}
\end{ex}

% ===== Câu 115 =====
% ID: L12C3B19-24-TL
% Mức: VDC
\dangbt{Nhận biết điện trường biến thiên, từ trường biến thiên và điện từ trường}
\begin{bt}
% ID: L12C3B19-24-TL
% Mức: VDC
So sánh nhận định đúng “Điện trường biến thiên sinh ra từ trường biến thiên và từ trường biến thiên sinh ra điện trường xoáy.” với nhận định sai “Điện trường và từ trường biến thiên tồn tại hoàn toàn độc lập.”; chỉ rõ căn cứ Vật lí.
\loigiai{
Cần nêu được: Điện trường biến thiên sinh ra từ trường biến thiên và từ trường biến thiên sinh ra điện trường xoáy. Khi vận dụng phải xác định đúng dữ kiện, phương chiều nếu có, điều kiện áp dụng, hệ đơn vị và kiểm tra tính hợp lí. Nhận định “Điện trường và từ trường biến thiên tồn tại hoàn toàn độc lập.” sai.
}
\end{bt}

% ===== Câu 116 =====
% ID: L12C3B19-24-TLN
% Mức: VDC
\begin{ex}
% ID: L12C3B19-24-TLN
% Mức: VDC
Trong bốn phát biểu về Nhận biết điện trường biến thiên, từ trường biến thiên và điện từ trường, có bao nhiêu phát biểu đúng? (Chỉ ghi một số nguyên.) (1) Khi giải cần xác định đúng dữ kiện, phương chiều, điều kiện áp dụng và đơn vị. (2) Điện trường và từ trường biến thiên tồn tại hoàn toàn độc lập. (3) Phải dùng đúng định luật cho tình huống. (4) Điện trường biến thiên sinh ra từ trường biến thiên và từ trường biến thiên sinh ra điện trường xoáy.
\shortans{3}
\loigiai{
Đối chiếu với kiến thức: Điện trường biến thiên sinh ra từ trường biến thiên và từ trường biến thiên sinh ra điện trường xoáy. Có 3 phát biểu đúng.
}
\end{ex}

% ===== Câu 117 =====
% ID: L12C3B19-24-TN
% Mức: TH
\dangbt{Nhận biết các tính chất cơ bản của sóng điện từ}
\begin{ex}
% ID: L12C3B19-24-TN
% Mức: TH
Nội dung nào mô tả đúng nhất Nhận biết các tính chất cơ bản của sóng điện từ?
\choice
{Có thể bỏ qua phương, chiều và đơn vị trong mọi trường hợp.}
{Kết quả không phụ thuộc dữ kiện và điều kiện áp dụng.}
{\True Trong chân không, vectơ $E$ và $B$ vuông góc nhau và cùng vuông góc phương truyền sóng.}
{Trong sóng điện từ, $E$ và $B$ song song với phương truyền.}
\loigiai{
Trong chân không, vectơ $E$ và $B$ vuông góc nhau và cùng vuông góc phương truyền sóng. Vì vậy phương án $C$ đúng; các phương án còn lại sai.
}
\end{ex}

% ===== Câu 118 =====
% ID: L12C3B19-25-DS
% Mức: VD
\dangbt{Phân loại phổ sóng điện từ và nêu ứng dụng}
\begin{ex}
% ID: L12C3B19-25-DS
% Mức: VD
Xét Phân loại phổ sóng điện từ và nêu ứng dụng (Tình huống 5). Hãy xác định tính đúng, sai của bốn phát biểu sau.
\choiceTF
{\True Phổ điện từ gồm sóng vô tuyến, vi ba, hồng ngoại, ánh sáng nhìn thấy, tử ngoại, tia $X$ và gamma.}
{Kết quả không phụ thuộc dữ kiện và có thể bỏ qua điều kiện.}
{Âm thanh là một miền của phổ sóng điện từ.}
{\True Khi giải cần xác định đúng dữ kiện, phương chiều, điều kiện áp dụng và đơn vị.}
\loigiai{
\begin{itemchoice}
\item (Đúng) Phổ điện từ gồm sóng vô tuyến, vi ba, hồng ngoại, ánh sáng nhìn thấy, tử ngoại, tia $X$ và gamma. (Vì): Phổ điện từ gồm sóng vô tuyến, vi ba, hồng ngoại, ánh sáng nhìn thấy, tử ngoại, tia $X$ và gamma. b) Sai: Kết quả không phụ thuộc dữ kiện và có thể bỏ qua điều kiện. c) Sai: Âm thanh là một miền của phổ sóng điện từ. d) Đúng: Khi giải cần xác định đúng dữ kiện, phương chiều, điều kiện áp dụng và đơn vị. Kiến thức cốt lõi: Phổ điện từ gồm sóng vô tuyến, vi ba, hồng ngoại, ánh sáng nhìn thấy, tử ngoại, tia $X$ và gamma
\item (Sai) Kết quả không phụ thuộc dữ kiện và có thể bỏ qua điều kiện.
\item (Sai) Âm thanh là một miền của phổ sóng điện từ.
\item (Đúng) Khi giải cần xác định đúng dữ kiện, phương chiều, điều kiện áp dụng và đơn vị.
\end{itemchoice}
}
\end{ex}

% ===== Câu 119 =====
% ID: L12C3B19-25-TL
% Mức: TH
\begin{bt}
% ID: L12C3B19-25-TL
% Mức: TH
Trình bày kiến thức cốt lõi của Phân loại phổ sóng điện từ và nêu ứng dụng và nêu điều kiện áp dụng.
\loigiai{
Cần nêu được: Phổ điện từ gồm sóng vô tuyến, vi ba, hồng ngoại, ánh sáng nhìn thấy, tử ngoại, tia $X$ và gamma. Khi vận dụng phải xác định đúng dữ kiện, phương chiều nếu có, điều kiện áp dụng, hệ đơn vị và kiểm tra tính hợp lí. Nhận định “Âm thanh là một miền của phổ sóng điện từ.” sai.
}
\end{bt}

% ===== Câu 120 =====
% ID: L12C3B19-25-TLN
% Mức: TH
\begin{ex}
% ID: L12C3B19-25-TLN
% Mức: TH
Trong bốn phát biểu về Phân loại phổ sóng điện từ và nêu ứng dụng, có bao nhiêu phát biểu đúng? (Chỉ ghi một số nguyên.) (1) Phổ điện từ gồm sóng vô tuyến, vi ba, hồng ngoại, ánh sáng nhìn thấy, tử ngoại, tia $X$ và gamma. (2) Âm thanh là một miền của phổ sóng điện từ. (3) Cần kiểm tra điều kiện áp dụng. (4) Có thể bỏ qua đơn vị.
\shortans{2}
\loigiai{
Đối chiếu với kiến thức: Phổ điện từ gồm sóng vô tuyến, vi ba, hồng ngoại, ánh sáng nhìn thấy, tử ngoại, tia $X$ và gamma. Có 2 phát biểu đúng.
}
\end{ex}

% ===== Câu 121 =====
% ID: L12C3B19-25-TN
% Mức: TH
\dangbt{Nhận biết các tính chất cơ bản của sóng điện từ}
\begin{ex}
% ID: L12C3B19-25-TN
% Mức: TH
Khi phân tích Nhận biết các tính chất cơ bản của sóng điện từ?
\choice
{Trong sóng điện từ, $E$ và $B$ song song với phương truyền.}
{Có thể bỏ qua phương, chiều và đơn vị trong mọi trường hợp.}
{Kết quả không phụ thuộc dữ kiện và điều kiện áp dụng.}
{\True Trong chân không, vectơ $E$ và $B$ vuông góc nhau và cùng vuông góc phương truyền sóng.}
\loigiai{
Trong chân không, vectơ $E$ và $B$ vuông góc nhau và cùng vuông góc phương truyền sóng. Vì vậy phương án $D$ đúng; các phương án còn lại sai.
}
\end{ex}

% ===== Câu 122 =====
% ID: L12C3B19-26-DS
% Mức: TH
\dangbt{Tính bước sóng, tần số và tốc độ truyền sóng điện từ}
\begin{ex}
% ID: L12C3B19-26-DS
% Mức: TH
Xét Tính bước sóng, tần số và tốc độ truyền sóng điện từ (Tình huống 1). Hãy xác định tính đúng, sai của bốn phát biểu sau.
\choiceTF
{\True Trong chân không, bước sóng lambda = c/f với c xấp xỉ 3.10^ $8\,\text{m}$ /s.}
{Bước sóng trong chân không bằng f/c.}
{\True Khi giải cần xác định đúng dữ kiện, phương chiều, điều kiện áp dụng và đơn vị.}
{Kết quả không phụ thuộc dữ kiện và có thể bỏ qua điều kiện.}
\loigiai{
\begin{itemchoice}
\item (Đúng) Trong chân không, bước sóng lambda = c/f với c xấp xỉ 3.10^ $8\,\text{m}$ /s. (Vì): Trong chân không, bước sóng lambda = c/f với c xấp xỉ 3.10^ $8\,\text{m}$ /s. b) Sai: Bước sóng trong chân không bằng f/c. c) Đúng: Khi giải cần xác định đúng dữ kiện, phương chiều, điều kiện áp dụng và đơn vị. d) Sai: Kết quả không phụ thuộc dữ kiện và có thể bỏ qua điều kiện. Kiến thức cốt lõi: Trong chân không, bước sóng lambda = c/f với c xấp xỉ 3.10^ $8\,\text{m}$ /s
\item (Sai) Bước sóng trong chân không bằng f/c.
\item (Đúng) Khi giải cần xác định đúng dữ kiện, phương chiều, điều kiện áp dụng và đơn vị.
\item (Sai) Kết quả không phụ thuộc dữ kiện và có thể bỏ qua điều kiện.
\end{itemchoice}
}
\end{ex}

% ===== Câu 123 =====
% ID: L12C3B19-26-TL
% Mức: TH
\dangbt{Phân loại phổ sóng điện từ và nêu ứng dụng}
\begin{bt}
% ID: L12C3B19-26-TL
% Mức: TH
Giải thích vì sao nhận định sau đúng: “Phổ điện từ gồm sóng vô tuyến, vi ba, hồng ngoại, ánh sáng nhìn thấy, tử ngoại, tia $X$ và gamma.”
\loigiai{
Cần nêu được: Phổ điện từ gồm sóng vô tuyến, vi ba, hồng ngoại, ánh sáng nhìn thấy, tử ngoại, tia $X$ và gamma. Khi vận dụng phải xác định đúng dữ kiện, phương chiều nếu có, điều kiện áp dụng, hệ đơn vị và kiểm tra tính hợp lí. Nhận định “Âm thanh là một miền của phổ sóng điện từ.” sai.
}
\end{bt}

% ===== Câu 124 =====
% ID: L12C3B19-26-TLN
% Mức: TH
\begin{ex}
% ID: L12C3B19-26-TLN
% Mức: TH
Trong bốn phát biểu về Phân loại phổ sóng điện từ và nêu ứng dụng, có bao nhiêu phát biểu đúng? (Chỉ ghi một số nguyên.) (1) Âm thanh là một miền của phổ sóng điện từ. (2) Phổ điện từ gồm sóng vô tuyến, vi ba, hồng ngoại, ánh sáng nhìn thấy, tử ngoại, tia $X$ và gamma. (3) Kết quả phải phù hợp ý nghĩa vật lí. (4) Mọi đại lượng đều là vô hướng.
\shortans{1}
\loigiai{
Đối chiếu với kiến thức: Phổ điện từ gồm sóng vô tuyến, vi ba, hồng ngoại, ánh sáng nhìn thấy, tử ngoại, tia $X$ và gamma. Có 1 phát biểu đúng.
}
\end{ex}

% ===== Câu 125 =====
% ID: L12C3B19-26-TN
% Mức: TH
\dangbt{Nhận biết các tính chất cơ bản của sóng điện từ}
\begin{ex}
% ID: L12C3B19-26-TN
% Mức: TH
Một học sinh ôn tập Nhận biết các tính chất cơ bản của sóng điện từ?
\choice
{\True Trong chân không, vectơ $E$ và $B$ vuông góc nhau và cùng vuông góc phương truyền sóng.}
{Trong sóng điện từ, $E$ và $B$ song song với phương truyền.}
{Có thể bỏ qua phương, chiều và đơn vị trong mọi trường hợp.}
{Kết quả không phụ thuộc dữ kiện và điều kiện áp dụng.}
\loigiai{
Trong chân không, vectơ $E$ và $B$ vuông góc nhau và cùng vuông góc phương truyền sóng. Vì vậy phương án $A$ đúng; các phương án còn lại sai.
}
\end{ex}

% ===== Câu 126 =====
% ID: L12C3B19-27-DS
% Mức: TH
\dangbt{Tính bước sóng, tần số và tốc độ truyền sóng điện từ}
\begin{ex}
% ID: L12C3B19-27-DS
% Mức: TH
Xét Tính bước sóng, tần số và tốc độ truyền sóng điện từ (Tình huống 2). Hãy xác định tính đúng, sai của bốn phát biểu sau.
\choiceTF
{Bước sóng trong chân không bằng f/c.}
{\True Trong chân không, bước sóng lambda = c/f với c xấp xỉ 3.10^ $8\,\text{m}$ /s.}
{Kết quả không phụ thuộc dữ kiện và có thể bỏ qua điều kiện.}
{\True Khi giải cần xác định đúng dữ kiện, phương chiều, điều kiện áp dụng và đơn vị.}
\loigiai{
\begin{itemchoice}
\item (Sai) Bước sóng trong chân không bằng f/c. (Vì): Bước sóng trong chân không bằng f/c. b) Đúng: Trong chân không, bước sóng lambda = c/f với c xấp xỉ 3.10^ $8\,\text{m}$ /s. c) Sai: Kết quả không phụ thuộc dữ kiện và có thể bỏ qua điều kiện. d) Đúng: Khi giải cần xác định đúng dữ kiện, phương chiều, điều kiện áp dụng và đơn vị. Kiến thức cốt lõi: Trong chân không, bước sóng lambda = c/f với c xấp xỉ 3.10^ $8\,\text{m}$ /s
\item (Đúng) Trong chân không, bước sóng lambda = c/f với c xấp xỉ 3.10^ $8\,\text{m}$ /s.
\item (Sai) Kết quả không phụ thuộc dữ kiện và có thể bỏ qua điều kiện.
\item (Đúng) Khi giải cần xác định đúng dữ kiện, phương chiều, điều kiện áp dụng và đơn vị.
\end{itemchoice}
}
\end{ex}

% ===== Câu 127 =====
% ID: L12C3B19-27-TL
% Mức: VD
\dangbt{Phân loại phổ sóng điện từ và nêu ứng dụng}
\begin{bt}
% ID: L12C3B19-27-TL
% Mức: VD
Phân tích sai lầm trong nhận định: “Âm thanh là một miền của phổ sóng điện từ.”
\loigiai{
Cần nêu được: Phổ điện từ gồm sóng vô tuyến, vi ba, hồng ngoại, ánh sáng nhìn thấy, tử ngoại, tia $X$ và gamma. Khi vận dụng phải xác định đúng dữ kiện, phương chiều nếu có, điều kiện áp dụng, hệ đơn vị và kiểm tra tính hợp lí. Nhận định “Âm thanh là một miền của phổ sóng điện từ.” sai.
}
\end{bt}

% ===== Câu 128 =====
% ID: L12C3B19-27-TLN
% Mức: TH
\begin{ex}
% ID: L12C3B19-27-TLN
% Mức: TH
Trong bốn phát biểu về Phân loại phổ sóng điện từ và nêu ứng dụng, có bao nhiêu phát biểu đúng? (Chỉ ghi một số nguyên.) (1) Phổ điện từ gồm sóng vô tuyến, vi ba, hồng ngoại, ánh sáng nhìn thấy, tử ngoại, tia $X$ và gamma. (2) Âm thanh là một miền của phổ sóng điện từ. (3) Cần kiểm tra điều kiện áp dụng. (4) Có thể bỏ qua đơn vị.
\shortans{2}
\loigiai{
Đối chiếu với kiến thức: Phổ điện từ gồm sóng vô tuyến, vi ba, hồng ngoại, ánh sáng nhìn thấy, tử ngoại, tia $X$ và gamma. Có 2 phát biểu đúng.
}
\end{ex}

% ===== Câu 129 =====
% ID: L12C3B19-27-TN
% Mức: TH
\dangbt{Nhận biết các tính chất cơ bản của sóng điện từ}
\begin{ex}
% ID: L12C3B19-27-TN
% Mức: TH
Chọn phát biểu không trái với kiến thức về Nhận biết các tính chất cơ bản của sóng điện từ?
\choice
{Kết quả không phụ thuộc dữ kiện và điều kiện áp dụng.}
{\True Trong chân không, vectơ $E$ và $B$ vuông góc nhau và cùng vuông góc phương truyền sóng.}
{Trong sóng điện từ, $E$ và $B$ song song với phương truyền.}
{Có thể bỏ qua phương, chiều và đơn vị trong mọi trường hợp.}
\loigiai{
Trong chân không, vectơ $E$ và $B$ vuông góc nhau và cùng vuông góc phương truyền sóng. Vì vậy phương án $B$ đúng; các phương án còn lại sai.
}
\end{ex}

% ===== Câu 130 =====
% ID: L12C3B19-28-DS
% Mức: VD
\dangbt{Tính bước sóng, tần số và tốc độ truyền sóng điện từ}
\begin{ex}
% ID: L12C3B19-28-DS
% Mức: VD
Xét Tính bước sóng, tần số và tốc độ truyền sóng điện từ (Tình huống 3). Hãy xác định tính đúng, sai của bốn phát biểu sau.
\choiceTF
{\True Khi giải cần xác định đúng dữ kiện, phương chiều, điều kiện áp dụng và đơn vị.}
{Kết quả không phụ thuộc dữ kiện và có thể bỏ qua điều kiện.}
{\True Trong chân không, bước sóng lambda = c/f với c xấp xỉ 3.10^ $8\,\text{m}$ /s.}
{Bước sóng trong chân không bằng f/c.}
\loigiai{
\begin{itemchoice}
\item (Đúng) Khi giải cần xác định đúng dữ kiện, phương chiều, điều kiện áp dụng và đơn vị. (Vì): Khi giải cần xác định đúng dữ kiện, phương chiều, điều kiện áp dụng và đơn vị. b) Sai: Kết quả không phụ thuộc dữ kiện và có thể bỏ qua điều kiện. c) Đúng: Trong chân không, bước sóng lambda = c/f với c xấp xỉ 3.10^ $8\,\text{m}$ /s. d) Sai: Bước sóng trong chân không bằng f/c. Kiến thức cốt lõi: Trong chân không, bước sóng lambda = c/f với c xấp xỉ 3.10^ $8\,\text{m}$ /s
\item (Sai) Kết quả không phụ thuộc dữ kiện và có thể bỏ qua điều kiện.
\item (Đúng) Trong chân không, bước sóng lambda = c/f với c xấp xỉ 3.10^ $8\,\text{m}$ /s.
\item (Sai) Bước sóng trong chân không bằng f/c.
\end{itemchoice}
}
\end{ex}

% ===== Câu 131 =====
% ID: L12C3B19-28-TL
% Mức: VD
\dangbt{Phân loại phổ sóng điện từ và nêu ứng dụng}
\begin{bt}
% ID: L12C3B19-28-TL
% Mức: VD
Đề xuất các bước giải một bài toán thuộc dạng Phân loại phổ sóng điện từ và nêu ứng dụng.
\loigiai{
Cần nêu được: Phổ điện từ gồm sóng vô tuyến, vi ba, hồng ngoại, ánh sáng nhìn thấy, tử ngoại, tia $X$ và gamma. Khi vận dụng phải xác định đúng dữ kiện, phương chiều nếu có, điều kiện áp dụng, hệ đơn vị và kiểm tra tính hợp lí. Nhận định “Âm thanh là một miền của phổ sóng điện từ.” sai.
}
\end{bt}

% ===== Câu 132 =====
% ID: L12C3B19-28-TLN
% Mức: VD
\begin{ex}
% ID: L12C3B19-28-TLN
% Mức: VD
Trong bốn phát biểu về Phân loại phổ sóng điện từ và nêu ứng dụng, có bao nhiêu phát biểu đúng? (Chỉ ghi một số nguyên.) (1) Âm thanh là một miền của phổ sóng điện từ. (2) Phổ điện từ gồm sóng vô tuyến, vi ba, hồng ngoại, ánh sáng nhìn thấy, tử ngoại, tia $X$ và gamma. (3) Kết quả phải phù hợp ý nghĩa vật lí. (4) Mọi đại lượng đều là vô hướng.
\shortans{2}
\loigiai{
Đối chiếu với kiến thức: Phổ điện từ gồm sóng vô tuyến, vi ba, hồng ngoại, ánh sáng nhìn thấy, tử ngoại, tia $X$ và gamma. Có 2 phát biểu đúng.
}
\end{ex}

% ===== Câu 133 =====
% ID: L12C3B19-28-TN
% Mức: VD
\dangbt{Nhận biết các tính chất cơ bản của sóng điện từ}
\begin{ex}
% ID: L12C3B19-28-TN
% Mức: VD
Cơ sở Vật lí đúng dùng để giải Nhận biết các tính chất cơ bản của sóng điện từ?
\choice
{Có thể bỏ qua phương, chiều và đơn vị trong mọi trường hợp.}
{Kết quả không phụ thuộc dữ kiện và điều kiện áp dụng.}
{\True Trong chân không, vectơ $E$ và $B$ vuông góc nhau và cùng vuông góc phương truyền sóng.}
{Trong sóng điện từ, $E$ và $B$ song song với phương truyền.}
\loigiai{
Trong chân không, vectơ $E$ và $B$ vuông góc nhau và cùng vuông góc phương truyền sóng. Vì vậy phương án $C$ đúng; các phương án còn lại sai.
}
\end{ex}

% ===== Câu 134 =====
% ID: L12C3B19-29-DS
% Mức: VD
\dangbt{Tính bước sóng, tần số và tốc độ truyền sóng điện từ}
\begin{ex}
% ID: L12C3B19-29-DS
% Mức: VD
Xét Tính bước sóng, tần số và tốc độ truyền sóng điện từ (Tình huống 4). Hãy xác định tính đúng, sai của bốn phát biểu sau.
\choiceTF
{Kết quả không phụ thuộc dữ kiện và có thể bỏ qua điều kiện.}
{\True Khi giải cần xác định đúng dữ kiện, phương chiều, điều kiện áp dụng và đơn vị.}
{Bước sóng trong chân không bằng f/c.}
{\True Trong chân không, bước sóng lambda = c/f với c xấp xỉ 3.10^ $8\,\text{m}$ /s.}
\loigiai{
\begin{itemchoice}
\item (Sai) Kết quả không phụ thuộc dữ kiện và có thể bỏ qua điều kiện. (Vì): Kết quả không phụ thuộc dữ kiện và có thể bỏ qua điều kiện. b) Đúng: Khi giải cần xác định đúng dữ kiện, phương chiều, điều kiện áp dụng và đơn vị. c) Sai: Bước sóng trong chân không bằng f/c. d) Đúng: Trong chân không, bước sóng lambda = c/f với c xấp xỉ 3.10^ $8\,\text{m}$ /s. Kiến thức cốt lõi: Trong chân không, bước sóng lambda = c/f với c xấp xỉ 3.10^ $8\,\text{m}$ /s
\item (Đúng) Khi giải cần xác định đúng dữ kiện, phương chiều, điều kiện áp dụng và đơn vị.
\item (Sai) Bước sóng trong chân không bằng f/c.
\item (Đúng) Trong chân không, bước sóng lambda = c/f với c xấp xỉ 3.10^ $8\,\text{m}$ /s.
\end{itemchoice}
}
\end{ex}

% ===== Câu 135 =====
% ID: L12C3B19-29-TL
% Mức: VD
\dangbt{Phân loại phổ sóng điện từ và nêu ứng dụng}
\begin{bt}
% ID: L12C3B19-29-TL
% Mức: VD
Nêu một tình huống thực tiễn liên quan đến Phân loại phổ sóng điện từ và nêu ứng dụng và giải thích bằng kiến thức Vật lí.
\loigiai{
Cần nêu được: Phổ điện từ gồm sóng vô tuyến, vi ba, hồng ngoại, ánh sáng nhìn thấy, tử ngoại, tia $X$ và gamma. Khi vận dụng phải xác định đúng dữ kiện, phương chiều nếu có, điều kiện áp dụng, hệ đơn vị và kiểm tra tính hợp lí. Nhận định “Âm thanh là một miền của phổ sóng điện từ.” sai.
}
\end{bt}

% ===== Câu 136 =====
% ID: L12C3B19-29-TLN
% Mức: VD
\begin{ex}
% ID: L12C3B19-29-TLN
% Mức: VD
Trong bốn phát biểu về Phân loại phổ sóng điện từ và nêu ứng dụng, có bao nhiêu phát biểu đúng? (Chỉ ghi một số nguyên.) (1) Phổ điện từ gồm sóng vô tuyến, vi ba, hồng ngoại, ánh sáng nhìn thấy, tử ngoại, tia $X$ và gamma. (2) Âm thanh là một miền của phổ sóng điện từ. (3) Cần kiểm tra điều kiện áp dụng. (4) Có thể bỏ qua đơn vị.
\shortans{2}
\loigiai{
Đối chiếu với kiến thức: Phổ điện từ gồm sóng vô tuyến, vi ba, hồng ngoại, ánh sáng nhìn thấy, tử ngoại, tia $X$ và gamma. Có 2 phát biểu đúng.
}
\end{ex}

% ===== Câu 137 =====
% ID: L12C3B19-29-TN
% Mức: VD
\dangbt{Nhận biết các tính chất cơ bản của sóng điện từ}
\begin{ex}
% ID: L12C3B19-29-TN
% Mức: VD
Trong một bài thực hành về Nhận biết các tính chất cơ bản của sóng điện từ?
\choice
{Trong sóng điện từ, $E$ và $B$ song song với phương truyền.}
{Có thể bỏ qua phương, chiều và đơn vị trong mọi trường hợp.}
{Kết quả không phụ thuộc dữ kiện và điều kiện áp dụng.}
{\True Trong chân không, vectơ $E$ và $B$ vuông góc nhau và cùng vuông góc phương truyền sóng.}
\loigiai{
Trong chân không, vectơ $E$ và $B$ vuông góc nhau và cùng vuông góc phương truyền sóng. Vì vậy phương án $D$ đúng; các phương án còn lại sai.
}
\end{ex}

% ===== Câu 138 =====
% ID: L12C3B19-30-DS
% Mức: VD
\dangbt{Tính bước sóng, tần số và tốc độ truyền sóng điện từ}
\begin{ex}
% ID: L12C3B19-30-DS
% Mức: VD
Xét Tính bước sóng, tần số và tốc độ truyền sóng điện từ (Tình huống 5). Hãy xác định tính đúng, sai của bốn phát biểu sau.
\choiceTF
{\True Trong chân không, bước sóng lambda = c/f với c xấp xỉ 3.10^ $8\,\text{m}$ /s.}
{Kết quả không phụ thuộc dữ kiện và có thể bỏ qua điều kiện.}
{Bước sóng trong chân không bằng f/c.}
{\True Khi giải cần xác định đúng dữ kiện, phương chiều, điều kiện áp dụng và đơn vị.}
\loigiai{
\begin{itemchoice}
\item (Đúng) Trong chân không, bước sóng lambda = c/f với c xấp xỉ 3.10^ $8\,\text{m}$ /s. (Vì): Trong chân không, bước sóng lambda = c/f với c xấp xỉ 3.10^ $8\,\text{m}$ /s. b) Sai: Kết quả không phụ thuộc dữ kiện và có thể bỏ qua điều kiện. c) Sai: Bước sóng trong chân không bằng f/c. d) Đúng: Khi giải cần xác định đúng dữ kiện, phương chiều, điều kiện áp dụng và đơn vị. Kiến thức cốt lõi: Trong chân không, bước sóng lambda = c/f với c xấp xỉ 3.10^ $8\,\text{m}$ /s
\item (Sai) Kết quả không phụ thuộc dữ kiện và có thể bỏ qua điều kiện.
\item (Sai) Bước sóng trong chân không bằng f/c.
\item (Đúng) Khi giải cần xác định đúng dữ kiện, phương chiều, điều kiện áp dụng và đơn vị.
\end{itemchoice}
}
\end{ex}

% ===== Câu 139 =====
% ID: L12C3B19-30-TL
% Mức: VDC
\dangbt{Phân loại phổ sóng điện từ và nêu ứng dụng}
\begin{bt}
% ID: L12C3B19-30-TL
% Mức: VDC
So sánh nhận định đúng “Phổ điện từ gồm sóng vô tuyến, vi ba, hồng ngoại, ánh sáng nhìn thấy, tử ngoại, tia $X$ và gamma.” với nhận định sai “Âm thanh là một miền của phổ sóng điện từ.”; chỉ rõ căn cứ Vật lí.
\loigiai{
Cần nêu được: Phổ điện từ gồm sóng vô tuyến, vi ba, hồng ngoại, ánh sáng nhìn thấy, tử ngoại, tia $X$ và gamma. Khi vận dụng phải xác định đúng dữ kiện, phương chiều nếu có, điều kiện áp dụng, hệ đơn vị và kiểm tra tính hợp lí. Nhận định “Âm thanh là một miền của phổ sóng điện từ.” sai.
}
\end{bt}

% ===== Câu 140 =====
% ID: L12C3B19-30-TLN
% Mức: VDC
\begin{ex}
% ID: L12C3B19-30-TLN
% Mức: VDC
Trong bốn phát biểu về Phân loại phổ sóng điện từ và nêu ứng dụng, có bao nhiêu phát biểu đúng? (Chỉ ghi một số nguyên.) (1) Khi giải cần xác định đúng dữ kiện, phương chiều, điều kiện áp dụng và đơn vị. (2) Âm thanh là một miền của phổ sóng điện từ. (3) Phải dùng đúng định luật cho tình huống. (4) Phổ điện từ gồm sóng vô tuyến, vi ba, hồng ngoại, ánh sáng nhìn thấy, tử ngoại, tia $X$ và gamma.
\shortans{3}
\loigiai{
Đối chiếu với kiến thức: Phổ điện từ gồm sóng vô tuyến, vi ba, hồng ngoại, ánh sáng nhìn thấy, tử ngoại, tia $X$ và gamma. Có 3 phát biểu đúng.
}
\end{ex}

% ===== Câu 141 =====
% ID: L12C3B19-30-TN
% Mức: VD
\dangbt{Nhận biết các tính chất cơ bản của sóng điện từ}
\begin{ex}
% ID: L12C3B19-30-TN
% Mức: VD
Kết luận đúng khi vận dụng Nhận biết các tính chất cơ bản của sóng điện từ?
\choice
{\True Trong chân không, vectơ $E$ và $B$ vuông góc nhau và cùng vuông góc phương truyền sóng.}
{Trong sóng điện từ, $E$ và $B$ song song với phương truyền.}
{Có thể bỏ qua phương, chiều và đơn vị trong mọi trường hợp.}
{Kết quả không phụ thuộc dữ kiện và điều kiện áp dụng.}
\loigiai{
Trong chân không, vectơ $E$ và $B$ vuông góc nhau và cùng vuông góc phương truyền sóng. Vì vậy phương án $A$ đúng; các phương án còn lại sai.
}
\end{ex}

% ===== Câu 142 =====
% ID: L12C3B19-31-DS
% Mức: TH
\dangbt{Nhận biết điện trường biến thiên, từ trường biến thiên và điện từ trường}
\begin{ex}
% ID: L12C3B19-31-DS
% Mức: TH
Xét Nhận biết điện trường biến thiên, từ trường biến thiên và điện từ trường (Tình huống 1). Hãy xác định tính đúng, sai của bốn phát biểu sau.
\choiceTF
{\True Điện trường biến thiên sinh ra từ trường biến thiên và từ trường biến thiên sinh ra điện trường xoáy.}
{Điện trường và từ trường biến thiên tồn tại hoàn toàn độc lập.}
{\True Khi giải cần xác định đúng dữ kiện, phương chiều, điều kiện áp dụng và đơn vị.}
{Kết quả không phụ thuộc dữ kiện và có thể bỏ qua điều kiện.}
\loigiai{
\begin{itemchoice}
\item (Đúng) Điện trường biến thiên sinh ra từ trường biến thiên và từ trường biến thiên sinh ra điện trường xoáy. (Vì): iện trường biến thiên sinh ra từ trường biến thiên và từ trường biến thiên sinh ra điện trường xoáy. b) Sai: Điện trường và từ trường biến thiên tồn tại hoàn toàn độc lập. c) Đúng: Khi giải cần xác định đúng dữ kiện, phương chiều, điều kiện áp dụng và đơn vị. d) Sai: Kết quả không phụ thuộc dữ kiện và có thể bỏ qua điều kiện. Kiến thức cốt lõi: Điện trường biến thiên sinh ra từ trường biến thiên và từ trường biến thiên sinh ra điện trường xoáy
\item (Sai) Điện trường và từ trường biến thiên tồn tại hoàn toàn độc lập.
\item (Đúng) Khi giải cần xác định đúng dữ kiện, phương chiều, điều kiện áp dụng và đơn vị.
\item (Sai) Kết quả không phụ thuộc dữ kiện và có thể bỏ qua điều kiện.
\end{itemchoice}
}
\end{ex}

% ===== Câu 143 =====
% ID: L12C3B19-31-TL
% Mức: TH
\dangbt{Tính bước sóng, tần số và tốc độ truyền sóng điện từ}
\begin{bt}
% ID: L12C3B19-31-TL
% Mức: TH
Trình bày kiến thức cốt lõi của Tính bước sóng, tần số và tốc độ truyền sóng điện từ và nêu điều kiện áp dụng.
\loigiai{
Cần nêu được: Trong chân không, bước sóng lambda = c/f với c xấp xỉ 3.10^ $8\,\text{m}$ /s. Khi vận dụng phải xác định đúng dữ kiện, phương chiều nếu có, điều kiện áp dụng, hệ đơn vị và kiểm tra tính hợp lí. Nhận định “Bước sóng trong chân không bằng f/c.” sai.
}
\end{bt}

% ===== Câu 144 =====
% ID: L12C3B19-31-TLN
% Mức: TH
\begin{ex}
% ID: L12C3B19-31-TLN
% Mức: TH
Trong bốn phát biểu về Tính bước sóng, tần số và tốc độ truyền sóng điện từ, có bao nhiêu phát biểu đúng? (Chỉ ghi một số nguyên.) (1) Trong chân không, bước sóng lambda = c/f với c xấp xỉ 3.10^ $8\,\text{m}$ /s. (2) Bước sóng trong chân không bằng f/c. (3) Cần kiểm tra điều kiện áp dụng. (4) Có thể bỏ qua đơn vị.
\shortans{2}
\loigiai{
Đối chiếu với kiến thức: Trong chân không, bước sóng lambda = c/f với c xấp xỉ 3.10^ $8\,\text{m}$ /s. Có 2 phát biểu đúng.
}
\end{ex}

% ===== Câu 145 =====
% ID: L12C3B19-31-TN
% Mức: NB
\dangbt{Nhận biết điện trường biến thiên, từ trường biến thiên và điện từ trường}
\begin{ex}
% ID: L12C3B19-31-TN
% Mức: NB
Phát biểu nào sau đây đúng về Nhận biết điện trường biến thiên, từ trường biến thiên và điện từ trường?
\choice
{Điện trường và từ trường biến thiên tồn tại hoàn toàn độc lập.}
{Có thể bỏ qua phương, chiều và đơn vị trong mọi trường hợp.}
{Kết quả không phụ thuộc dữ kiện và điều kiện áp dụng.}
{\True Điện trường biến thiên sinh ra từ trường biến thiên và từ trường biến thiên sinh ra điện trường xoáy.}
\loigiai{
Điện trường biến thiên sinh ra từ trường biến thiên và từ trường biến thiên sinh ra điện trường xoáy. Vì vậy phương án $D$ đúng; các phương án còn lại sai.
}
\end{ex}

% ===== Câu 146 =====
% ID: L12C3B19-32-DS
% Mức: TH
\begin{ex}
% ID: L12C3B19-32-DS
% Mức: TH
Xét Nhận biết điện trường biến thiên, từ trường biến thiên và điện từ trường (Tình huống 2). Hãy xác định tính đúng, sai của bốn phát biểu sau.
\choiceTF
{Điện trường và từ trường biến thiên tồn tại hoàn toàn độc lập.}
{\True Điện trường biến thiên sinh ra từ trường biến thiên và từ trường biến thiên sinh ra điện trường xoáy.}
{Kết quả không phụ thuộc dữ kiện và có thể bỏ qua điều kiện.}
{\True Khi giải cần xác định đúng dữ kiện, phương chiều, điều kiện áp dụng và đơn vị.}
\loigiai{
\begin{itemchoice}
\item (Sai) Điện trường và từ trường biến thiên tồn tại hoàn toàn độc lập. (Vì): Điện trường và từ trường biến thiên tồn tại hoàn toàn độc lập. b) Đúng: Điện trường biến thiên sinh ra từ trường biến thiên và từ trường biến thiên sinh ra điện trường xoáy. c) Sai: Kết quả không phụ thuộc dữ kiện và có thể bỏ qua điều kiện. d) Đúng: Khi giải cần xác định đúng dữ kiện, phương chiều, điều kiện áp dụng và đơn vị. Kiến thức cốt lõi: Điện trường biến thiên sinh ra từ trường biến thiên và từ trường biến thiên sinh ra điện trường xoáy
\item (Đúng) Điện trường biến thiên sinh ra từ trường biến thiên và từ trường biến thiên sinh ra điện trường xoáy.
\item (Sai) Kết quả không phụ thuộc dữ kiện và có thể bỏ qua điều kiện.
\item (Đúng) Khi giải cần xác định đúng dữ kiện, phương chiều, điều kiện áp dụng và đơn vị.
\end{itemchoice}
}
\end{ex}

% ===== Câu 147 =====
% ID: L12C3B19-32-TL
% Mức: TH
\dangbt{Tính bước sóng, tần số và tốc độ truyền sóng điện từ}
\begin{bt}
% ID: L12C3B19-32-TL
% Mức: TH
Giải thích vì sao nhận định sau đúng: “Trong chân không, bước sóng lambda = c/f với c xấp xỉ 3.10^ $8\,\text{m}$ /s.”
\loigiai{
Cần nêu được: Trong chân không, bước sóng lambda = c/f với c xấp xỉ 3.10^ $8\,\text{m}$ /s. Khi vận dụng phải xác định đúng dữ kiện, phương chiều nếu có, điều kiện áp dụng, hệ đơn vị và kiểm tra tính hợp lí. Nhận định “Bước sóng trong chân không bằng f/c.” sai.
}
\end{bt}

% ===== Câu 148 =====
% ID: L12C3B19-32-TLN
% Mức: TH
\begin{ex}
% ID: L12C3B19-32-TLN
% Mức: TH
Trong bốn phát biểu về Tính bước sóng, tần số và tốc độ truyền sóng điện từ, có bao nhiêu phát biểu đúng? (Chỉ ghi một số nguyên.) (1) Bước sóng trong chân không bằng f/c. (2) Trong chân không, bước sóng lambda = c/f với c xấp xỉ 3.10^ $8\,\text{m}$ /s. (3) Kết quả phải phù hợp ý nghĩa vật lí. (4) Mọi đại lượng đều là vô hướng.
\shortans{1}
\loigiai{
Đối chiếu với kiến thức: Trong chân không, bước sóng lambda = c/f với c xấp xỉ 3.10^ $8\,\text{m}$ /s. Có 1 phát biểu đúng.
}
\end{ex}

% ===== Câu 149 =====
% ID: L12C3B19-32-TN
% Mức: NB
\dangbt{Nhận biết điện trường biến thiên, từ trường biến thiên và điện từ trường}
\begin{ex}
% ID: L12C3B19-32-TN
% Mức: NB
Trong nội dung Nhận biết điện trường biến thiên, từ trường biến thiên và điện từ trường?
\choice
{\True Điện trường biến thiên sinh ra từ trường biến thiên và từ trường biến thiên sinh ra điện trường xoáy.}
{Điện trường và từ trường biến thiên tồn tại hoàn toàn độc lập.}
{Có thể bỏ qua phương, chiều và đơn vị trong mọi trường hợp.}
{Kết quả không phụ thuộc dữ kiện và điều kiện áp dụng.}
\loigiai{
Điện trường biến thiên sinh ra từ trường biến thiên và từ trường biến thiên sinh ra điện trường xoáy. Vì vậy phương án $A$ đúng; các phương án còn lại sai.
}
\end{ex}

% ===== Câu 150 =====
% ID: L12C3B19-33-DS
% Mức: VD
\begin{ex}
% ID: L12C3B19-33-DS
% Mức: VD
Xét Nhận biết điện trường biến thiên, từ trường biến thiên và điện từ trường (Tình huống 3). Hãy xác định tính đúng, sai của bốn phát biểu sau.
\choiceTF
{\True Khi giải cần xác định đúng dữ kiện, phương chiều, điều kiện áp dụng và đơn vị.}
{Kết quả không phụ thuộc dữ kiện và có thể bỏ qua điều kiện.}
{\True Điện trường biến thiên sinh ra từ trường biến thiên và từ trường biến thiên sinh ra điện trường xoáy.}
{Điện trường và từ trường biến thiên tồn tại hoàn toàn độc lập.}
\loigiai{
\begin{itemchoice}
\item (Đúng) Khi giải cần xác định đúng dữ kiện, phương chiều, điều kiện áp dụng và đơn vị. (Vì): Khi giải cần xác định đúng dữ kiện, phương chiều, điều kiện áp dụng và đơn vị. b) Sai: Kết quả không phụ thuộc dữ kiện và có thể bỏ qua điều kiện. c) Đúng: Điện trường biến thiên sinh ra từ trường biến thiên và từ trường biến thiên sinh ra điện trường xoáy. d) Sai: Điện trường và từ trường biến thiên tồn tại hoàn toàn độc lập. Kiến thức cốt lõi: Điện trường biến thiên sinh ra từ trường biến thiên và từ trường biến thiên sinh ra điện trường xoáy
\item (Sai) Kết quả không phụ thuộc dữ kiện và có thể bỏ qua điều kiện.
\item (Đúng) Điện trường biến thiên sinh ra từ trường biến thiên và từ trường biến thiên sinh ra điện trường xoáy.
\item (Sai) Điện trường và từ trường biến thiên tồn tại hoàn toàn độc lập.
\end{itemchoice}
}
\end{ex}

% ===== Câu 151 =====
% ID: L12C3B19-33-TL
% Mức: VD
\dangbt{Tính bước sóng, tần số và tốc độ truyền sóng điện từ}
\begin{bt}
% ID: L12C3B19-33-TL
% Mức: VD
Phân tích sai lầm trong nhận định: “Bước sóng trong chân không bằng f/c.”
\loigiai{
Cần nêu được: Trong chân không, bước sóng lambda = c/f với c xấp xỉ 3.10^ $8\,\text{m}$ /s. Khi vận dụng phải xác định đúng dữ kiện, phương chiều nếu có, điều kiện áp dụng, hệ đơn vị và kiểm tra tính hợp lí. Nhận định “Bước sóng trong chân không bằng f/c.” sai.
}
\end{bt}

% ===== Câu 152 =====
% ID: L12C3B19-33-TLN
% Mức: TH
\begin{ex}
% ID: L12C3B19-33-TLN
% Mức: TH
Trong bốn phát biểu về Tính bước sóng, tần số và tốc độ truyền sóng điện từ, có bao nhiêu phát biểu đúng? (Chỉ ghi một số nguyên.) (1) Trong chân không, bước sóng lambda = c/f với c xấp xỉ 3.10^ $8\,\text{m}$ /s. (2) Bước sóng trong chân không bằng f/c. (3) Cần kiểm tra điều kiện áp dụng. (4) Có thể bỏ qua đơn vị.
\shortans{2}
\loigiai{
Đối chiếu với kiến thức: Trong chân không, bước sóng lambda = c/f với c xấp xỉ 3.10^ $8\,\text{m}$ /s. Có 2 phát biểu đúng.
}
\end{ex}

% ===== Câu 153 =====
% ID: L12C3B19-33-TN
% Mức: NB
\dangbt{Nhận biết điện trường biến thiên, từ trường biến thiên và điện từ trường}
\begin{ex}
% ID: L12C3B19-33-TN
% Mức: NB
Tình huống 3: chọn kết luận chính xác về Nhận biết điện trường biến thiên, từ trường biến thiên và điện từ trường?
\choice
{Kết quả không phụ thuộc dữ kiện và điều kiện áp dụng.}
{\True Điện trường biến thiên sinh ra từ trường biến thiên và từ trường biến thiên sinh ra điện trường xoáy.}
{Điện trường và từ trường biến thiên tồn tại hoàn toàn độc lập.}
{Có thể bỏ qua phương, chiều và đơn vị trong mọi trường hợp.}
\loigiai{
Điện trường biến thiên sinh ra từ trường biến thiên và từ trường biến thiên sinh ra điện trường xoáy. Vì vậy phương án $B$ đúng; các phương án còn lại sai.
}
\end{ex}

% ===== Câu 154 =====
% ID: L12C3B19-34-DS
% Mức: VD
\begin{ex}
% ID: L12C3B19-34-DS
% Mức: VD
Xét Nhận biết điện trường biến thiên, từ trường biến thiên và điện từ trường (Tình huống 4). Hãy xác định tính đúng, sai của bốn phát biểu sau.
\choiceTF
{Kết quả không phụ thuộc dữ kiện và có thể bỏ qua điều kiện.}
{\True Khi giải cần xác định đúng dữ kiện, phương chiều, điều kiện áp dụng và đơn vị.}
{Điện trường và từ trường biến thiên tồn tại hoàn toàn độc lập.}
{\True Điện trường biến thiên sinh ra từ trường biến thiên và từ trường biến thiên sinh ra điện trường xoáy.}
\loigiai{
\begin{itemchoice}
\item (Sai) Kết quả không phụ thuộc dữ kiện và có thể bỏ qua điều kiện. (Vì): Kết quả không phụ thuộc dữ kiện và có thể bỏ qua điều kiện. b) Đúng: Khi giải cần xác định đúng dữ kiện, phương chiều, điều kiện áp dụng và đơn vị. c) Sai: Điện trường và từ trường biến thiên tồn tại hoàn toàn độc lập. d) Đúng: Điện trường biến thiên sinh ra từ trường biến thiên và từ trường biến thiên sinh ra điện trường xoáy. Kiến thức cốt lõi: Điện trường biến thiên sinh ra từ trường biến thiên và từ trường biến thiên sinh ra điện trường xoáy
\item (Đúng) Khi giải cần xác định đúng dữ kiện, phương chiều, điều kiện áp dụng và đơn vị.
\item (Sai) Điện trường và từ trường biến thiên tồn tại hoàn toàn độc lập.
\item (Đúng) Điện trường biến thiên sinh ra từ trường biến thiên và từ trường biến thiên sinh ra điện trường xoáy.
\end{itemchoice}
}
\end{ex}

% ===== Câu 155 =====
% ID: L12C3B19-34-TL
% Mức: VD
\dangbt{Tính bước sóng, tần số và tốc độ truyền sóng điện từ}
\begin{bt}
% ID: L12C3B19-34-TL
% Mức: VD
Đề xuất các bước giải một bài toán thuộc dạng Tính bước sóng, tần số và tốc độ truyền sóng điện từ.
\loigiai{
Cần nêu được: Trong chân không, bước sóng lambda = c/f với c xấp xỉ 3.10^ $8\,\text{m}$ /s. Khi vận dụng phải xác định đúng dữ kiện, phương chiều nếu có, điều kiện áp dụng, hệ đơn vị và kiểm tra tính hợp lí. Nhận định “Bước sóng trong chân không bằng f/c.” sai.
}
\end{bt}

% ===== Câu 156 =====
% ID: L12C3B19-34-TLN
% Mức: VD
\begin{ex}
% ID: L12C3B19-34-TLN
% Mức: VD
Trong bốn phát biểu về Tính bước sóng, tần số và tốc độ truyền sóng điện từ, có bao nhiêu phát biểu đúng? (Chỉ ghi một số nguyên.) (1) Bước sóng trong chân không bằng f/c. (2) Trong chân không, bước sóng lambda = c/f với c xấp xỉ 3.10^ $8\,\text{m}$ /s. (3) Kết quả phải phù hợp ý nghĩa vật lí. (4) Mọi đại lượng đều là vô hướng.
\shortans{2}
\loigiai{
Đối chiếu với kiến thức: Trong chân không, bước sóng lambda = c/f với c xấp xỉ 3.10^ $8\,\text{m}$ /s. Có 2 phát biểu đúng.
}
\end{ex}

% ===== Câu 157 =====
% ID: L12C3B19-34-TN
% Mức: TH
\dangbt{Nhận biết điện trường biến thiên, từ trường biến thiên và điện từ trường}
\begin{ex}
% ID: L12C3B19-34-TN
% Mức: TH
Nội dung nào mô tả đúng nhất Nhận biết điện trường biến thiên, từ trường biến thiên và điện từ trường?
\choice
{Có thể bỏ qua phương, chiều và đơn vị trong mọi trường hợp.}
{Kết quả không phụ thuộc dữ kiện và điều kiện áp dụng.}
{\True Điện trường biến thiên sinh ra từ trường biến thiên và từ trường biến thiên sinh ra điện trường xoáy.}
{Điện trường và từ trường biến thiên tồn tại hoàn toàn độc lập.}
\loigiai{
Điện trường biến thiên sinh ra từ trường biến thiên và từ trường biến thiên sinh ra điện trường xoáy. Vì vậy phương án $C$ đúng; các phương án còn lại sai.
}
\end{ex}

% ===== Câu 158 =====
% ID: L12C3B19-35-DS
% Mức: VD
\begin{ex}
% ID: L12C3B19-35-DS
% Mức: VD
Xét Nhận biết điện trường biến thiên, từ trường biến thiên và điện từ trường (Tình huống 5). Hãy xác định tính đúng, sai của bốn phát biểu sau.
\choiceTF
{\True Điện trường biến thiên sinh ra từ trường biến thiên và từ trường biến thiên sinh ra điện trường xoáy.}
{Kết quả không phụ thuộc dữ kiện và có thể bỏ qua điều kiện.}
{Điện trường và từ trường biến thiên tồn tại hoàn toàn độc lập.}
{\True Khi giải cần xác định đúng dữ kiện, phương chiều, điều kiện áp dụng và đơn vị.}
\loigiai{
\begin{itemchoice}
\item (Đúng) Điện trường biến thiên sinh ra từ trường biến thiên và từ trường biến thiên sinh ra điện trường xoáy. (Vì): iện trường biến thiên sinh ra từ trường biến thiên và từ trường biến thiên sinh ra điện trường xoáy. b) Sai: Kết quả không phụ thuộc dữ kiện và có thể bỏ qua điều kiện. c) Sai: Điện trường và từ trường biến thiên tồn tại hoàn toàn độc lập. d) Đúng: Khi giải cần xác định đúng dữ kiện, phương chiều, điều kiện áp dụng và đơn vị. Kiến thức cốt lõi: Điện trường biến thiên sinh ra từ trường biến thiên và từ trường biến thiên sinh ra điện trường xoáy
\item (Sai) Kết quả không phụ thuộc dữ kiện và có thể bỏ qua điều kiện.
\item (Sai) Điện trường và từ trường biến thiên tồn tại hoàn toàn độc lập.
\item (Đúng) Khi giải cần xác định đúng dữ kiện, phương chiều, điều kiện áp dụng và đơn vị.
\end{itemchoice}
}
\end{ex}

% ===== Câu 159 =====
% ID: L12C3B19-35-TL
% Mức: VD
\dangbt{Tính bước sóng, tần số và tốc độ truyền sóng điện từ}
\begin{bt}
% ID: L12C3B19-35-TL
% Mức: VD
Nêu một tình huống thực tiễn liên quan đến Tính bước sóng, tần số và tốc độ truyền sóng điện từ và giải thích bằng kiến thức Vật lí.
\loigiai{
Cần nêu được: Trong chân không, bước sóng lambda = c/f với c xấp xỉ 3.10^ $8\,\text{m}$ /s. Khi vận dụng phải xác định đúng dữ kiện, phương chiều nếu có, điều kiện áp dụng, hệ đơn vị và kiểm tra tính hợp lí. Nhận định “Bước sóng trong chân không bằng f/c.” sai.
}
\end{bt}

% ===== Câu 160 =====
% ID: L12C3B19-35-TLN
% Mức: VD
\begin{ex}
% ID: L12C3B19-35-TLN
% Mức: VD
Trong bốn phát biểu về Tính bước sóng, tần số và tốc độ truyền sóng điện từ, có bao nhiêu phát biểu đúng? (Chỉ ghi một số nguyên.) (1) Trong chân không, bước sóng lambda = c/f với c xấp xỉ 3.10^ $8\,\text{m}$ /s. (2) Bước sóng trong chân không bằng f/c. (3) Cần kiểm tra điều kiện áp dụng. (4) Có thể bỏ qua đơn vị.
\shortans{2}
\loigiai{
Đối chiếu với kiến thức: Trong chân không, bước sóng lambda = c/f với c xấp xỉ 3.10^ $8\,\text{m}$ /s. Có 2 phát biểu đúng.
}
\end{ex}

% ===== Câu 161 =====
% ID: L12C3B19-35-TN
% Mức: TH
\dangbt{Nhận biết điện trường biến thiên, từ trường biến thiên và điện từ trường}
\begin{ex}
% ID: L12C3B19-35-TN
% Mức: TH
Khi phân tích Nhận biết điện trường biến thiên, từ trường biến thiên và điện từ trường?
\choice
{Điện trường và từ trường biến thiên tồn tại hoàn toàn độc lập.}
{Có thể bỏ qua phương, chiều và đơn vị trong mọi trường hợp.}
{Kết quả không phụ thuộc dữ kiện và điều kiện áp dụng.}
{\True Điện trường biến thiên sinh ra từ trường biến thiên và từ trường biến thiên sinh ra điện trường xoáy.}
\loigiai{
Điện trường biến thiên sinh ra từ trường biến thiên và từ trường biến thiên sinh ra điện trường xoáy. Vì vậy phương án $D$ đúng; các phương án còn lại sai.
}
\end{ex}

% ===== Câu 162 =====
% ID: L12C3B19-36-DS
% Mức: TH
\dangbt{Mô tả sự hình thành và lan truyền của sóng điện từ}
\begin{ex}
% ID: L12C3B19-36-DS
% Mức: TH
Xét Mô tả sự hình thành và lan truyền của sóng điện từ (Tình huống 1). Hãy xác định tính đúng, sai của bốn phát biểu sau.
\choiceTF
{\True Sóng điện từ là điện từ trường lan truyền trong không gian và truyền được trong chân không.}
{Sóng điện từ bắt buộc cần môi trường vật chất để truyền.}
{\True Khi giải cần xác định đúng dữ kiện, phương chiều, điều kiện áp dụng và đơn vị.}
{Kết quả không phụ thuộc dữ kiện và có thể bỏ qua điều kiện.}
\loigiai{
\begin{itemchoice}
\item (Đúng) Sóng điện từ là điện từ trường lan truyền trong không gian và truyền được trong chân không. (Vì): Sóng điện từ là điện từ trường lan truyền trong không gian và truyền được trong chân không. b) Sai: Sóng điện từ bắt buộc cần môi trường vật chất để truyền. c) Đúng: Khi giải cần xác định đúng dữ kiện, phương chiều, điều kiện áp dụng và đơn vị. d) Sai: Kết quả không phụ thuộc dữ kiện và có thể bỏ qua điều kiện. Kiến thức cốt lõi: Sóng điện từ là điện từ trường lan truyền trong không gian và truyền được trong chân không
\item (Sai) Sóng điện từ bắt buộc cần môi trường vật chất để truyền.
\item (Đúng) Khi giải cần xác định đúng dữ kiện, phương chiều, điều kiện áp dụng và đơn vị.
\item (Sai) Kết quả không phụ thuộc dữ kiện và có thể bỏ qua điều kiện.
\end{itemchoice}
}
\end{ex}

% ===== Câu 163 =====
% ID: L12C3B19-36-TL
% Mức: VDC
\dangbt{Tính bước sóng, tần số và tốc độ truyền sóng điện từ}
\begin{bt}
% ID: L12C3B19-36-TL
% Mức: VDC
So sánh nhận định đúng “Trong chân không, bước sóng lambda = c/f với c xấp xỉ 3.10^ $8\,\text{m}$ /s.” với nhận định sai “Bước sóng trong chân không bằng f/c.”; chỉ rõ căn cứ Vật lí.
\loigiai{
Cần nêu được: Trong chân không, bước sóng lambda = c/f với c xấp xỉ 3.10^ $8\,\text{m}$ /s. Khi vận dụng phải xác định đúng dữ kiện, phương chiều nếu có, điều kiện áp dụng, hệ đơn vị và kiểm tra tính hợp lí. Nhận định “Bước sóng trong chân không bằng f/c.” sai.
}
\end{bt}

% ===== Câu 164 =====
% ID: L12C3B19-36-TLN
% Mức: VDC
\begin{ex}
% ID: L12C3B19-36-TLN
% Mức: VDC
Trong bốn phát biểu về Tính bước sóng, tần số và tốc độ truyền sóng điện từ, có bao nhiêu phát biểu đúng? (Chỉ ghi một số nguyên.) (1) Khi giải cần xác định đúng dữ kiện, phương chiều, điều kiện áp dụng và đơn vị. (2) Bước sóng trong chân không bằng f/c. (3) Phải dùng đúng định luật cho tình huống. (4) Trong chân không, bước sóng lambda = c/f với c xấp xỉ 3.10^ $8\,\text{m}$ /s.
\shortans{3}
\loigiai{
Đối chiếu với kiến thức: Trong chân không, bước sóng lambda = c/f với c xấp xỉ 3.10^ $8\,\text{m}$ /s. Có 3 phát biểu đúng.
}
\end{ex}

% ===== Câu 165 =====
% ID: L12C3B19-36-TN
% Mức: TH
\dangbt{Nhận biết điện trường biến thiên, từ trường biến thiên và điện từ trường}
\begin{ex}
% ID: L12C3B19-36-TN
% Mức: TH
Một học sinh ôn tập Nhận biết điện trường biến thiên, từ trường biến thiên và điện từ trường?
\choice
{\True Điện trường biến thiên sinh ra từ trường biến thiên và từ trường biến thiên sinh ra điện trường xoáy.}
{Điện trường và từ trường biến thiên tồn tại hoàn toàn độc lập.}
{Có thể bỏ qua phương, chiều và đơn vị trong mọi trường hợp.}
{Kết quả không phụ thuộc dữ kiện và điều kiện áp dụng.}
\loigiai{
Điện trường biến thiên sinh ra từ trường biến thiên và từ trường biến thiên sinh ra điện trường xoáy. Vì vậy phương án $A$ đúng; các phương án còn lại sai.
}
\end{ex}

% ===== Câu 166 =====
% ID: L12C3B19-37-DS
% Mức: TH
\dangbt{Mô tả sự hình thành và lan truyền của sóng điện từ}
\begin{ex}
% ID: L12C3B19-37-DS
% Mức: TH
Xét Mô tả sự hình thành và lan truyền của sóng điện từ (Tình huống 2). Hãy xác định tính đúng, sai của bốn phát biểu sau.
\choiceTF
{Sóng điện từ bắt buộc cần môi trường vật chất để truyền.}
{\True Sóng điện từ là điện từ trường lan truyền trong không gian và truyền được trong chân không.}
{Kết quả không phụ thuộc dữ kiện và có thể bỏ qua điều kiện.}
{\True Khi giải cần xác định đúng dữ kiện, phương chiều, điều kiện áp dụng và đơn vị.}
\loigiai{
\begin{itemchoice}
\item (Sai) Sóng điện từ bắt buộc cần môi trường vật chất để truyền. (Vì): óng điện từ bắt buộc cần môi trường vật chất để truyền. b) Đúng: Sóng điện từ là điện từ trường lan truyền trong không gian và truyền được trong chân không. c) Sai: Kết quả không phụ thuộc dữ kiện và có thể bỏ qua điều kiện. d) Đúng: Khi giải cần xác định đúng dữ kiện, phương chiều, điều kiện áp dụng và đơn vị. Kiến thức cốt lõi: Sóng điện từ là điện từ trường lan truyền trong không gian và truyền được trong chân không
\item (Đúng) Sóng điện từ là điện từ trường lan truyền trong không gian và truyền được trong chân không.
\item (Sai) Kết quả không phụ thuộc dữ kiện và có thể bỏ qua điều kiện.
\item (Đúng) Khi giải cần xác định đúng dữ kiện, phương chiều, điều kiện áp dụng và đơn vị.
\end{itemchoice}
}
\end{ex}

% ===== Câu 167 =====
% ID: L12C3B19-37-TL
% Mức: TH
\dangbt{Nhận biết điện trường biến thiên, từ trường biến thiên và điện từ trường}
\begin{bt}
% ID: L12C3B19-37-TL
% Mức: TH
Trình bày kiến thức cốt lõi của Nhận biết điện trường biến thiên, từ trường biến thiên và điện từ trường và nêu điều kiện áp dụng.
\loigiai{
Cần nêu được: Điện trường biến thiên sinh ra từ trường biến thiên và từ trường biến thiên sinh ra điện trường xoáy. Khi vận dụng phải xác định đúng dữ kiện, phương chiều nếu có, điều kiện áp dụng, hệ đơn vị và kiểm tra tính hợp lí. Nhận định “Điện trường và từ trường biến thiên tồn tại hoàn toàn độc lập.” sai.
}
\end{bt}

% ===== Câu 168 =====
% ID: L12C3B19-37-TLN
% Mức: TH
\begin{ex}
% ID: L12C3B19-37-TLN
% Mức: TH
Trong bốn phát biểu về Nhận biết điện trường biến thiên, từ trường biến thiên và điện từ trường, có bao nhiêu phát biểu đúng? (Chỉ ghi một số nguyên.) (1) Điện trường biến thiên sinh ra từ trường biến thiên và từ trường biến thiên sinh ra điện trường xoáy. (2) Điện trường và từ trường biến thiên tồn tại hoàn toàn độc lập. (3) Cần kiểm tra điều kiện áp dụng. (4) Có thể bỏ qua đơn vị.
\shortans{2}
\loigiai{
Đối chiếu với kiến thức: Điện trường biến thiên sinh ra từ trường biến thiên và từ trường biến thiên sinh ra điện trường xoáy. Có 2 phát biểu đúng.
}
\end{ex}

% ===== Câu 169 =====
% ID: L12C3B19-37-TN
% Mức: TH
\begin{ex}
% ID: L12C3B19-37-TN
% Mức: TH
Chọn phát biểu không trái với kiến thức về Nhận biết điện trường biến thiên, từ trường biến thiên và điện từ trường?
\choice
{Kết quả không phụ thuộc dữ kiện và điều kiện áp dụng.}
{\True Điện trường biến thiên sinh ra từ trường biến thiên và từ trường biến thiên sinh ra điện trường xoáy.}
{Điện trường và từ trường biến thiên tồn tại hoàn toàn độc lập.}
{Có thể bỏ qua phương, chiều và đơn vị trong mọi trường hợp.}
\loigiai{
Điện trường biến thiên sinh ra từ trường biến thiên và từ trường biến thiên sinh ra điện trường xoáy. Vì vậy phương án $B$ đúng; các phương án còn lại sai.
}
\end{ex}

% ===== Câu 170 =====
% ID: L12C3B19-38-DS
% Mức: VD
\dangbt{Mô tả sự hình thành và lan truyền của sóng điện từ}
\begin{ex}
% ID: L12C3B19-38-DS
% Mức: VD
Xét Mô tả sự hình thành và lan truyền của sóng điện từ (Tình huống 3). Hãy xác định tính đúng, sai của bốn phát biểu sau.
\choiceTF
{\True Khi giải cần xác định đúng dữ kiện, phương chiều, điều kiện áp dụng và đơn vị.}
{Kết quả không phụ thuộc dữ kiện và có thể bỏ qua điều kiện.}
{\True Sóng điện từ là điện từ trường lan truyền trong không gian và truyền được trong chân không.}
{Sóng điện từ bắt buộc cần môi trường vật chất để truyền.}
\loigiai{
\begin{itemchoice}
\item (Đúng) Khi giải cần xác định đúng dữ kiện, phương chiều, điều kiện áp dụng và đơn vị. (Vì): Khi giải cần xác định đúng dữ kiện, phương chiều, điều kiện áp dụng và đơn vị. b) Sai: Kết quả không phụ thuộc dữ kiện và có thể bỏ qua điều kiện. c) Đúng: Sóng điện từ là điện từ trường lan truyền trong không gian và truyền được trong chân không. d) Sai: Sóng điện từ bắt buộc cần môi trường vật chất để truyền. Kiến thức cốt lõi: Sóng điện từ là điện từ trường lan truyền trong không gian và truyền được trong chân không
\item (Sai) Kết quả không phụ thuộc dữ kiện và có thể bỏ qua điều kiện.
\item (Đúng) Sóng điện từ là điện từ trường lan truyền trong không gian và truyền được trong chân không.
\item (Sai) Sóng điện từ bắt buộc cần môi trường vật chất để truyền.
\end{itemchoice}
}
\end{ex}

% ===== Câu 171 =====
% ID: L12C3B19-38-TL
% Mức: TH
\dangbt{Nhận biết điện trường biến thiên, từ trường biến thiên và điện từ trường}
\begin{bt}
% ID: L12C3B19-38-TL
% Mức: TH
Giải thích vì sao nhận định sau đúng: “Điện trường biến thiên sinh ra từ trường biến thiên và từ trường biến thiên sinh ra điện trường xoáy.”
\loigiai{
Cần nêu được: Điện trường biến thiên sinh ra từ trường biến thiên và từ trường biến thiên sinh ra điện trường xoáy. Khi vận dụng phải xác định đúng dữ kiện, phương chiều nếu có, điều kiện áp dụng, hệ đơn vị và kiểm tra tính hợp lí. Nhận định “Điện trường và từ trường biến thiên tồn tại hoàn toàn độc lập.” sai.
}
\end{bt}

% ===== Câu 172 =====
% ID: L12C3B19-38-TLN
% Mức: TH
\begin{ex}
% ID: L12C3B19-38-TLN
% Mức: TH
Trong bốn phát biểu về Nhận biết điện trường biến thiên, từ trường biến thiên và điện từ trường, có bao nhiêu phát biểu đúng? (Chỉ ghi một số nguyên.) (1) Điện trường và từ trường biến thiên tồn tại hoàn toàn độc lập. (2) Điện trường biến thiên sinh ra từ trường biến thiên và từ trường biến thiên sinh ra điện trường xoáy. (3) Kết quả phải phù hợp ý nghĩa vật lí. (4) Mọi đại lượng đều là vô hướng.
\shortans{1}
\loigiai{
Đối chiếu với kiến thức: Điện trường biến thiên sinh ra từ trường biến thiên và từ trường biến thiên sinh ra điện trường xoáy. Có 1 phát biểu đúng.
}
\end{ex}

% ===== Câu 173 =====
% ID: L12C3B19-38-TN
% Mức: VD
\begin{ex}
% ID: L12C3B19-38-TN
% Mức: VD
Cơ sở Vật lí đúng dùng để giải Nhận biết điện trường biến thiên, từ trường biến thiên và điện từ trường?
\choice
{Có thể bỏ qua phương, chiều và đơn vị trong mọi trường hợp.}
{Kết quả không phụ thuộc dữ kiện và điều kiện áp dụng.}
{\True Điện trường biến thiên sinh ra từ trường biến thiên và từ trường biến thiên sinh ra điện trường xoáy.}
{Điện trường và từ trường biến thiên tồn tại hoàn toàn độc lập.}
\loigiai{
Điện trường biến thiên sinh ra từ trường biến thiên và từ trường biến thiên sinh ra điện trường xoáy. Vì vậy phương án $C$ đúng; các phương án còn lại sai.
}
\end{ex}

% ===== Câu 174 =====
% ID: L12C3B19-39-DS
% Mức: VD
\dangbt{Mô tả sự hình thành và lan truyền của sóng điện từ}
\begin{ex}
% ID: L12C3B19-39-DS
% Mức: VD
Xét Mô tả sự hình thành và lan truyền của sóng điện từ (Tình huống 4). Hãy xác định tính đúng, sai của bốn phát biểu sau.
\choiceTF
{Kết quả không phụ thuộc dữ kiện và có thể bỏ qua điều kiện.}
{\True Khi giải cần xác định đúng dữ kiện, phương chiều, điều kiện áp dụng và đơn vị.}
{Sóng điện từ bắt buộc cần môi trường vật chất để truyền.}
{\True Sóng điện từ là điện từ trường lan truyền trong không gian và truyền được trong chân không.}
\loigiai{
\begin{itemchoice}
\item (Sai) Kết quả không phụ thuộc dữ kiện và có thể bỏ qua điều kiện. (Vì): Kết quả không phụ thuộc dữ kiện và có thể bỏ qua điều kiện. b) Đúng: Khi giải cần xác định đúng dữ kiện, phương chiều, điều kiện áp dụng và đơn vị. c) Sai: Sóng điện từ bắt buộc cần môi trường vật chất để truyền. d) Đúng: Sóng điện từ là điện từ trường lan truyền trong không gian và truyền được trong chân không. Kiến thức cốt lõi: Sóng điện từ là điện từ trường lan truyền trong không gian và truyền được trong chân không
\item (Đúng) Khi giải cần xác định đúng dữ kiện, phương chiều, điều kiện áp dụng và đơn vị.
\item (Sai) Sóng điện từ bắt buộc cần môi trường vật chất để truyền.
\item (Đúng) Sóng điện từ là điện từ trường lan truyền trong không gian và truyền được trong chân không.
\end{itemchoice}
}
\end{ex}

% ===== Câu 175 =====
% ID: L12C3B19-39-TL
% Mức: VD
\dangbt{Nhận biết điện trường biến thiên, từ trường biến thiên và điện từ trường}
\begin{bt}
% ID: L12C3B19-39-TL
% Mức: VD
Phân tích sai lầm trong nhận định: “Điện trường và từ trường biến thiên tồn tại hoàn toàn độc lập.”
\loigiai{
Cần nêu được: Điện trường biến thiên sinh ra từ trường biến thiên và từ trường biến thiên sinh ra điện trường xoáy. Khi vận dụng phải xác định đúng dữ kiện, phương chiều nếu có, điều kiện áp dụng, hệ đơn vị và kiểm tra tính hợp lí. Nhận định “Điện trường và từ trường biến thiên tồn tại hoàn toàn độc lập.” sai.
}
\end{bt}

% ===== Câu 176 =====
% ID: L12C3B19-39-TLN
% Mức: TH
\begin{ex}
% ID: L12C3B19-39-TLN
% Mức: TH
Trong bốn phát biểu về Nhận biết điện trường biến thiên, từ trường biến thiên và điện từ trường, có bao nhiêu phát biểu đúng? (Chỉ ghi một số nguyên.) (1) Điện trường biến thiên sinh ra từ trường biến thiên và từ trường biến thiên sinh ra điện trường xoáy. (2) Điện trường và từ trường biến thiên tồn tại hoàn toàn độc lập. (3) Cần kiểm tra điều kiện áp dụng. (4) Có thể bỏ qua đơn vị.
\shortans{2}
\loigiai{
Đối chiếu với kiến thức: Điện trường biến thiên sinh ra từ trường biến thiên và từ trường biến thiên sinh ra điện trường xoáy. Có 2 phát biểu đúng.
}
\end{ex}

% ===== Câu 177 =====
% ID: L12C3B19-39-TN
% Mức: VD
\begin{ex}
% ID: L12C3B19-39-TN
% Mức: VD
Trong một bài thực hành về Nhận biết điện trường biến thiên, từ trường biến thiên và điện từ trường?
\choice
{Điện trường và từ trường biến thiên tồn tại hoàn toàn độc lập.}
{Có thể bỏ qua phương, chiều và đơn vị trong mọi trường hợp.}
{Kết quả không phụ thuộc dữ kiện và điều kiện áp dụng.}
{\True Điện trường biến thiên sinh ra từ trường biến thiên và từ trường biến thiên sinh ra điện trường xoáy.}
\loigiai{
Điện trường biến thiên sinh ra từ trường biến thiên và từ trường biến thiên sinh ra điện trường xoáy. Vì vậy phương án $D$ đúng; các phương án còn lại sai.
}
\end{ex}

% ===== Câu 178 =====
% ID: L12C3B19-40-DS
% Mức: VD
\dangbt{Mô tả sự hình thành và lan truyền của sóng điện từ}
\begin{ex}
% ID: L12C3B19-40-DS
% Mức: VD
Xét Mô tả sự hình thành và lan truyền của sóng điện từ (Tình huống 5). Hãy xác định tính đúng, sai của bốn phát biểu sau.
\choiceTF
{\True Sóng điện từ là điện từ trường lan truyền trong không gian và truyền được trong chân không.}
{Kết quả không phụ thuộc dữ kiện và có thể bỏ qua điều kiện.}
{Sóng điện từ bắt buộc cần môi trường vật chất để truyền.}
{\True Khi giải cần xác định đúng dữ kiện, phương chiều, điều kiện áp dụng và đơn vị.}
\loigiai{
\begin{itemchoice}
\item (Đúng) Sóng điện từ là điện từ trường lan truyền trong không gian và truyền được trong chân không. (Vì): Sóng điện từ là điện từ trường lan truyền trong không gian và truyền được trong chân không. b) Sai: Kết quả không phụ thuộc dữ kiện và có thể bỏ qua điều kiện. c) Sai: Sóng điện từ bắt buộc cần môi trường vật chất để truyền. d) Đúng: Khi giải cần xác định đúng dữ kiện, phương chiều, điều kiện áp dụng và đơn vị. Kiến thức cốt lõi: Sóng điện từ là điện từ trường lan truyền trong không gian và truyền được trong chân không
\item (Sai) Kết quả không phụ thuộc dữ kiện và có thể bỏ qua điều kiện.
\item (Sai) Sóng điện từ bắt buộc cần môi trường vật chất để truyền.
\item (Đúng) Khi giải cần xác định đúng dữ kiện, phương chiều, điều kiện áp dụng và đơn vị.
\end{itemchoice}
}
\end{ex}

% ===== Câu 179 =====
% ID: L12C3B19-40-TL
% Mức: VD
\dangbt{Nhận biết điện trường biến thiên, từ trường biến thiên và điện từ trường}
\begin{bt}
% ID: L12C3B19-40-TL
% Mức: VD
Đề xuất các bước giải một bài toán thuộc dạng Nhận biết điện trường biến thiên, từ trường biến thiên và điện từ trường.
\loigiai{
Cần nêu được: Điện trường biến thiên sinh ra từ trường biến thiên và từ trường biến thiên sinh ra điện trường xoáy. Khi vận dụng phải xác định đúng dữ kiện, phương chiều nếu có, điều kiện áp dụng, hệ đơn vị và kiểm tra tính hợp lí. Nhận định “Điện trường và từ trường biến thiên tồn tại hoàn toàn độc lập.” sai.
}
\end{bt}

% ===== Câu 180 =====
% ID: L12C3B19-40-TLN
% Mức: VD
\begin{ex}
% ID: L12C3B19-40-TLN
% Mức: VD
Trong bốn phát biểu về Nhận biết điện trường biến thiên, từ trường biến thiên và điện từ trường, có bao nhiêu phát biểu đúng? (Chỉ ghi một số nguyên.) (1) Điện trường và từ trường biến thiên tồn tại hoàn toàn độc lập. (2) Điện trường biến thiên sinh ra từ trường biến thiên và từ trường biến thiên sinh ra điện trường xoáy. (3) Kết quả phải phù hợp ý nghĩa vật lí. (4) Mọi đại lượng đều là vô hướng.
\shortans{2}
\loigiai{
Đối chiếu với kiến thức: Điện trường biến thiên sinh ra từ trường biến thiên và từ trường biến thiên sinh ra điện trường xoáy. Có 2 phát biểu đúng.
}
\end{ex}

% ===== Câu 181 =====
% ID: L12C3B19-40-TN
% Mức: VD
\begin{ex}
% ID: L12C3B19-40-TN
% Mức: VD
Kết luận đúng khi vận dụng Nhận biết điện trường biến thiên, từ trường biến thiên và điện từ trường?
\choice
{\True Điện trường biến thiên sinh ra từ trường biến thiên và từ trường biến thiên sinh ra điện trường xoáy.}
{Điện trường và từ trường biến thiên tồn tại hoàn toàn độc lập.}
{Có thể bỏ qua phương, chiều và đơn vị trong mọi trường hợp.}
{Kết quả không phụ thuộc dữ kiện và điều kiện áp dụng.}
\loigiai{
Điện trường biến thiên sinh ra từ trường biến thiên và từ trường biến thiên sinh ra điện trường xoáy. Vì vậy phương án $A$ đúng; các phương án còn lại sai.
}
\end{ex}

% ===== Câu 182 =====
% ID: L12C3B19-41-DS
% Mức: TH
\dangbt{Nhận biết các tính chất cơ bản của sóng điện từ}
\begin{ex}
% ID: L12C3B19-41-DS
% Mức: TH
Xét Nhận biết các tính chất cơ bản của sóng điện từ (Tình huống 1). Hãy xác định tính đúng, sai của bốn phát biểu sau.
\choiceTF
{\True Trong chân không, vectơ $E$ và $B$ vuông góc nhau và cùng vuông góc phương truyền sóng.}
{Trong sóng điện từ, $E$ và $B$ song song với phương truyền.}
{\True Khi giải cần xác định đúng dữ kiện, phương chiều, điều kiện áp dụng và đơn vị.}
{Kết quả không phụ thuộc dữ kiện và có thể bỏ qua điều kiện.}
\loigiai{
\begin{itemchoice}
\item (Đúng) Trong chân không, vectơ $E$ và $B$ vuông góc nhau và cùng vuông góc phương truyền sóng. (Vì): Trong chân không, vectơ $E$ và $B$ vuông góc nhau và cùng vuông góc phương truyền sóng. b) Sai: Trong sóng điện từ, $E$ và $B$ song song với phương truyền. c) Đúng: Khi giải cần xác định đúng dữ kiện, phương chiều, điều kiện áp dụng và đơn vị. d) Sai: Kết quả không phụ thuộc dữ kiện và có thể bỏ qua điều kiện. Kiến thức cốt lõi: Trong chân không, vectơ $E$ và $B$ vuông góc nhau và cùng vuông góc phương truyền sóng
\item (Sai) Trong sóng điện từ, $E$ và $B$ song song với phương truyền.
\item (Đúng) Khi giải cần xác định đúng dữ kiện, phương chiều, điều kiện áp dụng và đơn vị.
\item (Sai) Kết quả không phụ thuộc dữ kiện và có thể bỏ qua điều kiện.
\end{itemchoice}
}
\end{ex}

% ===== Câu 183 =====
% ID: L12C3B19-41-TL
% Mức: VD
\dangbt{Nhận biết điện trường biến thiên, từ trường biến thiên và điện từ trường}
\begin{bt}
% ID: L12C3B19-41-TL
% Mức: VD
Nêu một tình huống thực tiễn liên quan đến Nhận biết điện trường biến thiên, từ trường biến thiên và điện từ trường và giải thích bằng kiến thức Vật lí.
\loigiai{
Cần nêu được: Điện trường biến thiên sinh ra từ trường biến thiên và từ trường biến thiên sinh ra điện trường xoáy. Khi vận dụng phải xác định đúng dữ kiện, phương chiều nếu có, điều kiện áp dụng, hệ đơn vị và kiểm tra tính hợp lí. Nhận định “Điện trường và từ trường biến thiên tồn tại hoàn toàn độc lập.” sai.
}
\end{bt}

% ===== Câu 184 =====
% ID: L12C3B19-41-TLN
% Mức: VD
\begin{ex}
% ID: L12C3B19-41-TLN
% Mức: VD
Trong bốn phát biểu về Nhận biết điện trường biến thiên, từ trường biến thiên và điện từ trường, có bao nhiêu phát biểu đúng? (Chỉ ghi một số nguyên.) (1) Điện trường biến thiên sinh ra từ trường biến thiên và từ trường biến thiên sinh ra điện trường xoáy. (2) Điện trường và từ trường biến thiên tồn tại hoàn toàn độc lập. (3) Cần kiểm tra điều kiện áp dụng. (4) Có thể bỏ qua đơn vị.
\shortans{2}
\loigiai{
Đối chiếu với kiến thức: Điện trường biến thiên sinh ra từ trường biến thiên và từ trường biến thiên sinh ra điện trường xoáy. Có 2 phát biểu đúng.
}
\end{ex}

% ===== Câu 185 =====
% ID: L12C3B19-41-TN
% Mức: NB
\dangbt{Phân loại phổ sóng điện từ và nêu ứng dụng}
\begin{ex}
% ID: L12C3B19-41-TN
% Mức: NB
Phát biểu nào sau đây đúng về Phân loại phổ sóng điện từ và nêu ứng dụng?
\choice
{Âm thanh là một miền của phổ sóng điện từ.}
{Có thể bỏ qua phương, chiều và đơn vị trong mọi trường hợp.}
{Kết quả không phụ thuộc dữ kiện và điều kiện áp dụng.}
{\True Phổ điện từ gồm sóng vô tuyến, vi ba, hồng ngoại, ánh sáng nhìn thấy, tử ngoại, tia $X$ và gamma.}
\loigiai{
Phổ điện từ gồm sóng vô tuyến, vi ba, hồng ngoại, ánh sáng nhìn thấy, tử ngoại, tia $X$ và gamma. Vì vậy phương án $D$ đúng; các phương án còn lại sai.
}
\end{ex}

% ===== Câu 186 =====
% ID: L12C3B19-42-DS
% Mức: TH
\dangbt{Nhận biết các tính chất cơ bản của sóng điện từ}
\begin{ex}
% ID: L12C3B19-42-DS
% Mức: TH
Xét Nhận biết các tính chất cơ bản của sóng điện từ (Tình huống 2). Hãy xác định tính đúng, sai của bốn phát biểu sau.
\choiceTF
{Trong sóng điện từ, $E$ và $B$ song song với phương truyền.}
{\True Trong chân không, vectơ $E$ và $B$ vuông góc nhau và cùng vuông góc phương truyền sóng.}
{Kết quả không phụ thuộc dữ kiện và có thể bỏ qua điều kiện.}
{\True Khi giải cần xác định đúng dữ kiện, phương chiều, điều kiện áp dụng và đơn vị.}
\loigiai{
\begin{itemchoice}
\item (Sai) Trong sóng điện từ, $E$ và $B$ song song với phương truyền. (Vì): Trong sóng điện từ, $E$ và $B$ song song với phương truyền. b) Đúng: Trong chân không, vectơ $E$ và $B$ vuông góc nhau và cùng vuông góc phương truyền sóng. c) Sai: Kết quả không phụ thuộc dữ kiện và có thể bỏ qua điều kiện. d) Đúng: Khi giải cần xác định đúng dữ kiện, phương chiều, điều kiện áp dụng và đơn vị. Kiến thức cốt lõi: Trong chân không, vectơ $E$ và $B$ vuông góc nhau và cùng vuông góc phương truyền sóng
\item (Đúng) Trong chân không, vectơ $E$ và $B$ vuông góc nhau và cùng vuông góc phương truyền sóng.
\item (Sai) Kết quả không phụ thuộc dữ kiện và có thể bỏ qua điều kiện.
\item (Đúng) Khi giải cần xác định đúng dữ kiện, phương chiều, điều kiện áp dụng và đơn vị.
\end{itemchoice}
}
\end{ex}

% ===== Câu 187 =====
% ID: L12C3B19-42-TL
% Mức: VDC
\dangbt{Nhận biết điện trường biến thiên, từ trường biến thiên và điện từ trường}
\begin{bt}
% ID: L12C3B19-42-TL
% Mức: VDC
So sánh nhận định đúng “Điện trường biến thiên sinh ra từ trường biến thiên và từ trường biến thiên sinh ra điện trường xoáy.” với nhận định sai “Điện trường và từ trường biến thiên tồn tại hoàn toàn độc lập.”; chỉ rõ căn cứ Vật lí.
\loigiai{
Cần nêu được: Điện trường biến thiên sinh ra từ trường biến thiên và từ trường biến thiên sinh ra điện trường xoáy. Khi vận dụng phải xác định đúng dữ kiện, phương chiều nếu có, điều kiện áp dụng, hệ đơn vị và kiểm tra tính hợp lí. Nhận định “Điện trường và từ trường biến thiên tồn tại hoàn toàn độc lập.” sai.
}
\end{bt}

% ===== Câu 188 =====
% ID: L12C3B19-42-TLN
% Mức: VDC
\begin{ex}
% ID: L12C3B19-42-TLN
% Mức: VDC
Trong bốn phát biểu về Nhận biết điện trường biến thiên, từ trường biến thiên và điện từ trường, có bao nhiêu phát biểu đúng? (Chỉ ghi một số nguyên.) (1) Khi giải cần xác định đúng dữ kiện, phương chiều, điều kiện áp dụng và đơn vị. (2) Điện trường và từ trường biến thiên tồn tại hoàn toàn độc lập. (3) Phải dùng đúng định luật cho tình huống. (4) Điện trường biến thiên sinh ra từ trường biến thiên và từ trường biến thiên sinh ra điện trường xoáy.
\shortans{3}
\loigiai{
Đối chiếu với kiến thức: Điện trường biến thiên sinh ra từ trường biến thiên và từ trường biến thiên sinh ra điện trường xoáy. Có 3 phát biểu đúng.
}
\end{ex}

% ===== Câu 189 =====
% ID: L12C3B19-42-TN
% Mức: NB
\dangbt{Phân loại phổ sóng điện từ và nêu ứng dụng}
\begin{ex}
% ID: L12C3B19-42-TN
% Mức: NB
Trong nội dung Phân loại phổ sóng điện từ và nêu ứng dụng?
\choice
{\True Phổ điện từ gồm sóng vô tuyến, vi ba, hồng ngoại, ánh sáng nhìn thấy, tử ngoại, tia $X$ và gamma.}
{Âm thanh là một miền của phổ sóng điện từ.}
{Có thể bỏ qua phương, chiều và đơn vị trong mọi trường hợp.}
{Kết quả không phụ thuộc dữ kiện và điều kiện áp dụng.}
\loigiai{
Phổ điện từ gồm sóng vô tuyến, vi ba, hồng ngoại, ánh sáng nhìn thấy, tử ngoại, tia $X$ và gamma. Vì vậy phương án $A$ đúng; các phương án còn lại sai.
}
\end{ex}

% ===== Câu 190 =====
% ID: L12C3B19-43-DS
% Mức: VD
\dangbt{Nhận biết các tính chất cơ bản của sóng điện từ}
\begin{ex}
% ID: L12C3B19-43-DS
% Mức: VD
Xét Nhận biết các tính chất cơ bản của sóng điện từ (Tình huống 3). Hãy xác định tính đúng, sai của bốn phát biểu sau.
\choiceTF
{\True Khi giải cần xác định đúng dữ kiện, phương chiều, điều kiện áp dụng và đơn vị.}
{Kết quả không phụ thuộc dữ kiện và có thể bỏ qua điều kiện.}
{\True Trong chân không, vectơ $E$ và $B$ vuông góc nhau và cùng vuông góc phương truyền sóng.}
{Trong sóng điện từ, $E$ và $B$ song song với phương truyền.}
\loigiai{
\begin{itemchoice}
\item (Đúng) Khi giải cần xác định đúng dữ kiện, phương chiều, điều kiện áp dụng và đơn vị. (Vì): Khi giải cần xác định đúng dữ kiện, phương chiều, điều kiện áp dụng và đơn vị. b) Sai: Kết quả không phụ thuộc dữ kiện và có thể bỏ qua điều kiện. c) Đúng: Trong chân không, vectơ $E$ và $B$ vuông góc nhau và cùng vuông góc phương truyền sóng. d) Sai: Trong sóng điện từ, $E$ và $B$ song song với phương truyền. Kiến thức cốt lõi: Trong chân không, vectơ $E$ và $B$ vuông góc nhau và cùng vuông góc phương truyền sóng
\item (Sai) Kết quả không phụ thuộc dữ kiện và có thể bỏ qua điều kiện.
\item (Đúng) Trong chân không, vectơ $E$ và $B$ vuông góc nhau và cùng vuông góc phương truyền sóng.
\item (Sai) Trong sóng điện từ, $E$ và $B$ song song với phương truyền.
\end{itemchoice}
}
\end{ex}

% ===== Câu 191 =====
% ID: L12C3B19-43-TL
% Mức: TH
\dangbt{Mô tả sự hình thành và lan truyền của sóng điện từ}
\begin{bt}
% ID: L12C3B19-43-TL
% Mức: TH
Trình bày kiến thức cốt lõi của Mô tả sự hình thành và lan truyền của sóng điện từ và nêu điều kiện áp dụng.
\loigiai{
Cần nêu được: Sóng điện từ là điện từ trường lan truyền trong không gian và truyền được trong chân không. Khi vận dụng phải xác định đúng dữ kiện, phương chiều nếu có, điều kiện áp dụng, hệ đơn vị và kiểm tra tính hợp lí. Nhận định “Sóng điện từ bắt buộc cần môi trường vật chất để truyền.” sai.
}
\end{bt}

% ===== Câu 192 =====
% ID: L12C3B19-43-TLN
% Mức: TH
\begin{ex}
% ID: L12C3B19-43-TLN
% Mức: TH
Trong bốn phát biểu về Mô tả sự hình thành và lan truyền của sóng điện từ, có bao nhiêu phát biểu đúng? (Chỉ ghi một số nguyên.) (1) Sóng điện từ là điện từ trường lan truyền trong không gian và truyền được trong chân không. (2) Sóng điện từ bắt buộc cần môi trường vật chất để truyền. (3) Cần kiểm tra điều kiện áp dụng. (4) Có thể bỏ qua đơn vị.
\shortans{2}
\loigiai{
Đối chiếu với kiến thức: Sóng điện từ là điện từ trường lan truyền trong không gian và truyền được trong chân không. Có 2 phát biểu đúng.
}
\end{ex}

% ===== Câu 193 =====
% ID: L12C3B19-43-TN
% Mức: NB
\dangbt{Phân loại phổ sóng điện từ và nêu ứng dụng}
\begin{ex}
% ID: L12C3B19-43-TN
% Mức: NB
Tình huống 3: chọn kết luận chính xác về Phân loại phổ sóng điện từ và nêu ứng dụng?
\choice
{Kết quả không phụ thuộc dữ kiện và điều kiện áp dụng.}
{\True Phổ điện từ gồm sóng vô tuyến, vi ba, hồng ngoại, ánh sáng nhìn thấy, tử ngoại, tia $X$ và gamma.}
{Âm thanh là một miền của phổ sóng điện từ.}
{Có thể bỏ qua phương, chiều và đơn vị trong mọi trường hợp.}
\loigiai{
Phổ điện từ gồm sóng vô tuyến, vi ba, hồng ngoại, ánh sáng nhìn thấy, tử ngoại, tia $X$ và gamma. Vì vậy phương án $B$ đúng; các phương án còn lại sai.
}
\end{ex}

% ===== Câu 194 =====
% ID: L12C3B19-44-DS
% Mức: VD
\dangbt{Nhận biết các tính chất cơ bản của sóng điện từ}
\begin{ex}
% ID: L12C3B19-44-DS
% Mức: VD
Xét Nhận biết các tính chất cơ bản của sóng điện từ (Tình huống 4). Hãy xác định tính đúng, sai của bốn phát biểu sau.
\choiceTF
{Kết quả không phụ thuộc dữ kiện và có thể bỏ qua điều kiện.}
{\True Khi giải cần xác định đúng dữ kiện, phương chiều, điều kiện áp dụng và đơn vị.}
{Trong sóng điện từ, $E$ và $B$ song song với phương truyền.}
{\True Trong chân không, vectơ $E$ và $B$ vuông góc nhau và cùng vuông góc phương truyền sóng.}
\loigiai{
\begin{itemchoice}
\item (Sai) Kết quả không phụ thuộc dữ kiện và có thể bỏ qua điều kiện. (Vì): Kết quả không phụ thuộc dữ kiện và có thể bỏ qua điều kiện. b) Đúng: Khi giải cần xác định đúng dữ kiện, phương chiều, điều kiện áp dụng và đơn vị. c) Sai: Trong sóng điện từ, $E$ và $B$ song song với phương truyền. d) Đúng: Trong chân không, vectơ $E$ và $B$ vuông góc nhau và cùng vuông góc phương truyền sóng. Kiến thức cốt lõi: Trong chân không, vectơ $E$ và $B$ vuông góc nhau và cùng vuông góc phương truyền sóng
\item (Đúng) Khi giải cần xác định đúng dữ kiện, phương chiều, điều kiện áp dụng và đơn vị.
\item (Sai) Trong sóng điện từ, $E$ và $B$ song song với phương truyền.
\item (Đúng) Trong chân không, vectơ $E$ và $B$ vuông góc nhau và cùng vuông góc phương truyền sóng.
\end{itemchoice}
}
\end{ex}

% ===== Câu 195 =====
% ID: L12C3B19-44-TL
% Mức: TH
\dangbt{Mô tả sự hình thành và lan truyền của sóng điện từ}
\begin{bt}
% ID: L12C3B19-44-TL
% Mức: TH
Giải thích vì sao nhận định sau đúng: “Sóng điện từ là điện từ trường lan truyền trong không gian và truyền được trong chân không.”
\loigiai{
Cần nêu được: Sóng điện từ là điện từ trường lan truyền trong không gian và truyền được trong chân không. Khi vận dụng phải xác định đúng dữ kiện, phương chiều nếu có, điều kiện áp dụng, hệ đơn vị và kiểm tra tính hợp lí. Nhận định “Sóng điện từ bắt buộc cần môi trường vật chất để truyền.” sai.
}
\end{bt}

% ===== Câu 196 =====
% ID: L12C3B19-44-TLN
% Mức: TH
\begin{ex}
% ID: L12C3B19-44-TLN
% Mức: TH
Trong bốn phát biểu về Mô tả sự hình thành và lan truyền của sóng điện từ, có bao nhiêu phát biểu đúng? (Chỉ ghi một số nguyên.) (1) Sóng điện từ bắt buộc cần môi trường vật chất để truyền. (2) Sóng điện từ là điện từ trường lan truyền trong không gian và truyền được trong chân không. (3) Kết quả phải phù hợp ý nghĩa vật lí. (4) Mọi đại lượng đều là vô hướng.
\shortans{1}
\loigiai{
Đối chiếu với kiến thức: Sóng điện từ là điện từ trường lan truyền trong không gian và truyền được trong chân không. Có 1 phát biểu đúng.
}
\end{ex}

% ===== Câu 197 =====
% ID: L12C3B19-44-TN
% Mức: TH
\dangbt{Phân loại phổ sóng điện từ và nêu ứng dụng}
\begin{ex}
% ID: L12C3B19-44-TN
% Mức: TH
Nội dung nào mô tả đúng nhất Phân loại phổ sóng điện từ và nêu ứng dụng?
\choice
{Có thể bỏ qua phương, chiều và đơn vị trong mọi trường hợp.}
{Kết quả không phụ thuộc dữ kiện và điều kiện áp dụng.}
{\True Phổ điện từ gồm sóng vô tuyến, vi ba, hồng ngoại, ánh sáng nhìn thấy, tử ngoại, tia $X$ và gamma.}
{Âm thanh là một miền của phổ sóng điện từ.}
\loigiai{
Phổ điện từ gồm sóng vô tuyến, vi ba, hồng ngoại, ánh sáng nhìn thấy, tử ngoại, tia $X$ và gamma. Vì vậy phương án $C$ đúng; các phương án còn lại sai.
}
\end{ex}

% ===== Câu 198 =====
% ID: L12C3B19-45-DS
% Mức: VD
\dangbt{Nhận biết các tính chất cơ bản của sóng điện từ}
\begin{ex}
% ID: L12C3B19-45-DS
% Mức: VD
Xét Nhận biết các tính chất cơ bản của sóng điện từ (Tình huống 5). Hãy xác định tính đúng, sai của bốn phát biểu sau.
\choiceTF
{\True Trong chân không, vectơ $E$ và $B$ vuông góc nhau và cùng vuông góc phương truyền sóng.}
{Kết quả không phụ thuộc dữ kiện và có thể bỏ qua điều kiện.}
{Trong sóng điện từ, $E$ và $B$ song song với phương truyền.}
{\True Khi giải cần xác định đúng dữ kiện, phương chiều, điều kiện áp dụng và đơn vị.}
\loigiai{
\begin{itemchoice}
\item (Đúng) Trong chân không, vectơ $E$ và $B$ vuông góc nhau và cùng vuông góc phương truyền sóng. (Vì): Trong chân không, vectơ $E$ và $B$ vuông góc nhau và cùng vuông góc phương truyền sóng. b) Sai: Kết quả không phụ thuộc dữ kiện và có thể bỏ qua điều kiện. c) Sai: Trong sóng điện từ, $E$ và $B$ song song với phương truyền. d) Đúng: Khi giải cần xác định đúng dữ kiện, phương chiều, điều kiện áp dụng và đơn vị. Kiến thức cốt lõi: Trong chân không, vectơ $E$ và $B$ vuông góc nhau và cùng vuông góc phương truyền sóng
\item (Sai) Kết quả không phụ thuộc dữ kiện và có thể bỏ qua điều kiện.
\item (Sai) Trong sóng điện từ, $E$ và $B$ song song với phương truyền.
\item (Đúng) Khi giải cần xác định đúng dữ kiện, phương chiều, điều kiện áp dụng và đơn vị.
\end{itemchoice}
}
\end{ex}

% ===== Câu 199 =====
% ID: L12C3B19-45-TL
% Mức: VD
\dangbt{Mô tả sự hình thành và lan truyền của sóng điện từ}
\begin{bt}
% ID: L12C3B19-45-TL
% Mức: VD
Phân tích sai lầm trong nhận định: “Sóng điện từ bắt buộc cần môi trường vật chất để truyền.”
\loigiai{
Cần nêu được: Sóng điện từ là điện từ trường lan truyền trong không gian và truyền được trong chân không. Khi vận dụng phải xác định đúng dữ kiện, phương chiều nếu có, điều kiện áp dụng, hệ đơn vị và kiểm tra tính hợp lí. Nhận định “Sóng điện từ bắt buộc cần môi trường vật chất để truyền.” sai.
}
\end{bt}

% ===== Câu 200 =====
% ID: L12C3B19-45-TLN
% Mức: TH
\begin{ex}
% ID: L12C3B19-45-TLN
% Mức: TH
Trong bốn phát biểu về Mô tả sự hình thành và lan truyền của sóng điện từ, có bao nhiêu phát biểu đúng? (Chỉ ghi một số nguyên.) (1) Sóng điện từ là điện từ trường lan truyền trong không gian và truyền được trong chân không. (2) Sóng điện từ bắt buộc cần môi trường vật chất để truyền. (3) Cần kiểm tra điều kiện áp dụng. (4) Có thể bỏ qua đơn vị.
\shortans{2}
\loigiai{
Đối chiếu với kiến thức: Sóng điện từ là điện từ trường lan truyền trong không gian và truyền được trong chân không. Có 2 phát biểu đúng.
}
\end{ex}

% ===== Câu 201 =====
% ID: L12C3B19-45-TN
% Mức: TH
\dangbt{Phân loại phổ sóng điện từ và nêu ứng dụng}
\begin{ex}
% ID: L12C3B19-45-TN
% Mức: TH
Khi phân tích Phân loại phổ sóng điện từ và nêu ứng dụng?
\choice
{Âm thanh là một miền của phổ sóng điện từ.}
{Có thể bỏ qua phương, chiều và đơn vị trong mọi trường hợp.}
{Kết quả không phụ thuộc dữ kiện và điều kiện áp dụng.}
{\True Phổ điện từ gồm sóng vô tuyến, vi ba, hồng ngoại, ánh sáng nhìn thấy, tử ngoại, tia $X$ và gamma.}
\loigiai{
Phổ điện từ gồm sóng vô tuyến, vi ba, hồng ngoại, ánh sáng nhìn thấy, tử ngoại, tia $X$ và gamma. Vì vậy phương án $D$ đúng; các phương án còn lại sai.
}
\end{ex}

% ===== Câu 202 =====
% ID: L12C3B19-46-DS
% Mức: TH
\dangbt{Tính bước sóng, tần số và tốc độ truyền sóng điện từ}
\begin{ex}
% ID: L12C3B19-46-DS
% Mức: TH
Xét Tính bước sóng, tần số và tốc độ truyền sóng điện từ (Tình huống 1). Hãy xác định tính đúng, sai của bốn phát biểu sau.
\choiceTF
{\True Trong chân không, bước sóng lambda = c/f với c xấp xỉ 3.10^ $8\,\text{m}$ /s.}
{Bước sóng trong chân không bằng f/c.}
{\True Khi giải cần xác định đúng dữ kiện, phương chiều, điều kiện áp dụng và đơn vị.}
{Kết quả không phụ thuộc dữ kiện và có thể bỏ qua điều kiện.}
\loigiai{
\begin{itemchoice}
\item (Đúng) Trong chân không, bước sóng lambda = c/f với c xấp xỉ 3.10^ $8\,\text{m}$ /s. (Vì): Trong chân không, bước sóng lambda = c/f với c xấp xỉ 3.10^ $8\,\text{m}$ /s. b) Sai: Bước sóng trong chân không bằng f/c. c) Đúng: Khi giải cần xác định đúng dữ kiện, phương chiều, điều kiện áp dụng và đơn vị. d) Sai: Kết quả không phụ thuộc dữ kiện và có thể bỏ qua điều kiện. Kiến thức cốt lõi: Trong chân không, bước sóng lambda = c/f với c xấp xỉ 3.10^ $8\,\text{m}$ /s
\item (Sai) Bước sóng trong chân không bằng f/c.
\item (Đúng) Khi giải cần xác định đúng dữ kiện, phương chiều, điều kiện áp dụng và đơn vị.
\item (Sai) Kết quả không phụ thuộc dữ kiện và có thể bỏ qua điều kiện.
\end{itemchoice}
}
\end{ex}

% ===== Câu 203 =====
% ID: L12C3B19-46-TL
% Mức: VD
\dangbt{Mô tả sự hình thành và lan truyền của sóng điện từ}
\begin{bt}
% ID: L12C3B19-46-TL
% Mức: VD
Đề xuất các bước giải một bài toán thuộc dạng Mô tả sự hình thành và lan truyền của sóng điện từ.
\loigiai{
Cần nêu được: Sóng điện từ là điện từ trường lan truyền trong không gian và truyền được trong chân không. Khi vận dụng phải xác định đúng dữ kiện, phương chiều nếu có, điều kiện áp dụng, hệ đơn vị và kiểm tra tính hợp lí. Nhận định “Sóng điện từ bắt buộc cần môi trường vật chất để truyền.” sai.
}
\end{bt}

% ===== Câu 204 =====
% ID: L12C3B19-46-TLN
% Mức: VD
\begin{ex}
% ID: L12C3B19-46-TLN
% Mức: VD
Trong bốn phát biểu về Mô tả sự hình thành và lan truyền của sóng điện từ, có bao nhiêu phát biểu đúng? (Chỉ ghi một số nguyên.) (1) Sóng điện từ bắt buộc cần môi trường vật chất để truyền. (2) Sóng điện từ là điện từ trường lan truyền trong không gian và truyền được trong chân không. (3) Kết quả phải phù hợp ý nghĩa vật lí. (4) Mọi đại lượng đều là vô hướng.
\shortans{2}
\loigiai{
Đối chiếu với kiến thức: Sóng điện từ là điện từ trường lan truyền trong không gian và truyền được trong chân không. Có 2 phát biểu đúng.
}
\end{ex}

% ===== Câu 205 =====
% ID: L12C3B19-46-TN
% Mức: TH
\dangbt{Phân loại phổ sóng điện từ và nêu ứng dụng}
\begin{ex}
% ID: L12C3B19-46-TN
% Mức: TH
Một học sinh ôn tập Phân loại phổ sóng điện từ và nêu ứng dụng?
\choice
{\True Phổ điện từ gồm sóng vô tuyến, vi ba, hồng ngoại, ánh sáng nhìn thấy, tử ngoại, tia $X$ và gamma.}
{Âm thanh là một miền của phổ sóng điện từ.}
{Có thể bỏ qua phương, chiều và đơn vị trong mọi trường hợp.}
{Kết quả không phụ thuộc dữ kiện và điều kiện áp dụng.}
\loigiai{
Phổ điện từ gồm sóng vô tuyến, vi ba, hồng ngoại, ánh sáng nhìn thấy, tử ngoại, tia $X$ và gamma. Vì vậy phương án $A$ đúng; các phương án còn lại sai.
}
\end{ex}

% ===== Câu 206 =====
% ID: L12C3B19-47-DS
% Mức: TH
\dangbt{Tính bước sóng, tần số và tốc độ truyền sóng điện từ}
\begin{ex}
% ID: L12C3B19-47-DS
% Mức: TH
Xét Tính bước sóng, tần số và tốc độ truyền sóng điện từ (Tình huống 2). Hãy xác định tính đúng, sai của bốn phát biểu sau.
\choiceTF
{Bước sóng trong chân không bằng f/c.}
{\True Trong chân không, bước sóng lambda = c/f với c xấp xỉ 3.10^ $8\,\text{m}$ /s.}
{Kết quả không phụ thuộc dữ kiện và có thể bỏ qua điều kiện.}
{\True Khi giải cần xác định đúng dữ kiện, phương chiều, điều kiện áp dụng và đơn vị.}
\loigiai{
\begin{itemchoice}
\item (Sai) Bước sóng trong chân không bằng f/c. (Vì): Bước sóng trong chân không bằng f/c. b) Đúng: Trong chân không, bước sóng lambda = c/f với c xấp xỉ 3.10^ $8\,\text{m}$ /s. c) Sai: Kết quả không phụ thuộc dữ kiện và có thể bỏ qua điều kiện. d) Đúng: Khi giải cần xác định đúng dữ kiện, phương chiều, điều kiện áp dụng và đơn vị. Kiến thức cốt lõi: Trong chân không, bước sóng lambda = c/f với c xấp xỉ 3.10^ $8\,\text{m}$ /s
\item (Đúng) Trong chân không, bước sóng lambda = c/f với c xấp xỉ 3.10^ $8\,\text{m}$ /s.
\item (Sai) Kết quả không phụ thuộc dữ kiện và có thể bỏ qua điều kiện.
\item (Đúng) Khi giải cần xác định đúng dữ kiện, phương chiều, điều kiện áp dụng và đơn vị.
\end{itemchoice}
}
\end{ex}

% ===== Câu 207 =====
% ID: L12C3B19-47-TL
% Mức: VD
\dangbt{Mô tả sự hình thành và lan truyền của sóng điện từ}
\begin{bt}
% ID: L12C3B19-47-TL
% Mức: VD
Nêu một tình huống thực tiễn liên quan đến Mô tả sự hình thành và lan truyền của sóng điện từ và giải thích bằng kiến thức Vật lí.
\loigiai{
Cần nêu được: Sóng điện từ là điện từ trường lan truyền trong không gian và truyền được trong chân không. Khi vận dụng phải xác định đúng dữ kiện, phương chiều nếu có, điều kiện áp dụng, hệ đơn vị và kiểm tra tính hợp lí. Nhận định “Sóng điện từ bắt buộc cần môi trường vật chất để truyền.” sai.
}
\end{bt}

% ===== Câu 208 =====
% ID: L12C3B19-47-TLN
% Mức: VD
\begin{ex}
% ID: L12C3B19-47-TLN
% Mức: VD
Trong bốn phát biểu về Mô tả sự hình thành và lan truyền của sóng điện từ, có bao nhiêu phát biểu đúng? (Chỉ ghi một số nguyên.) (1) Sóng điện từ là điện từ trường lan truyền trong không gian và truyền được trong chân không. (2) Sóng điện từ bắt buộc cần môi trường vật chất để truyền. (3) Cần kiểm tra điều kiện áp dụng. (4) Có thể bỏ qua đơn vị.
\shortans{2}
\loigiai{
Đối chiếu với kiến thức: Sóng điện từ là điện từ trường lan truyền trong không gian và truyền được trong chân không. Có 2 phát biểu đúng.
}
\end{ex}

% ===== Câu 209 =====
% ID: L12C3B19-47-TN
% Mức: TH
\dangbt{Phân loại phổ sóng điện từ và nêu ứng dụng}
\begin{ex}
% ID: L12C3B19-47-TN
% Mức: TH
Chọn phát biểu không trái với kiến thức về Phân loại phổ sóng điện từ và nêu ứng dụng?
\choice
{Kết quả không phụ thuộc dữ kiện và điều kiện áp dụng.}
{\True Phổ điện từ gồm sóng vô tuyến, vi ba, hồng ngoại, ánh sáng nhìn thấy, tử ngoại, tia $X$ và gamma.}
{Âm thanh là một miền của phổ sóng điện từ.}
{Có thể bỏ qua phương, chiều và đơn vị trong mọi trường hợp.}
\loigiai{
Phổ điện từ gồm sóng vô tuyến, vi ba, hồng ngoại, ánh sáng nhìn thấy, tử ngoại, tia $X$ và gamma. Vì vậy phương án $B$ đúng; các phương án còn lại sai.
}
\end{ex}

% ===== Câu 210 =====
% ID: L12C3B19-48-DS
% Mức: VD
\dangbt{Tính bước sóng, tần số và tốc độ truyền sóng điện từ}
\begin{ex}
% ID: L12C3B19-48-DS
% Mức: VD
Xét Tính bước sóng, tần số và tốc độ truyền sóng điện từ (Tình huống 3). Hãy xác định tính đúng, sai của bốn phát biểu sau.
\choiceTF
{\True Khi giải cần xác định đúng dữ kiện, phương chiều, điều kiện áp dụng và đơn vị.}
{Kết quả không phụ thuộc dữ kiện và có thể bỏ qua điều kiện.}
{\True Trong chân không, bước sóng lambda = c/f với c xấp xỉ 3.10^ $8\,\text{m}$ /s.}
{Bước sóng trong chân không bằng f/c.}
\loigiai{
\begin{itemchoice}
\item (Đúng) Khi giải cần xác định đúng dữ kiện, phương chiều, điều kiện áp dụng và đơn vị. (Vì): Khi giải cần xác định đúng dữ kiện, phương chiều, điều kiện áp dụng và đơn vị. b) Sai: Kết quả không phụ thuộc dữ kiện và có thể bỏ qua điều kiện. c) Đúng: Trong chân không, bước sóng lambda = c/f với c xấp xỉ 3.10^ $8\,\text{m}$ /s. d) Sai: Bước sóng trong chân không bằng f/c. Kiến thức cốt lõi: Trong chân không, bước sóng lambda = c/f với c xấp xỉ 3.10^ $8\,\text{m}$ /s
\item (Sai) Kết quả không phụ thuộc dữ kiện và có thể bỏ qua điều kiện.
\item (Đúng) Trong chân không, bước sóng lambda = c/f với c xấp xỉ 3.10^ $8\,\text{m}$ /s.
\item (Sai) Bước sóng trong chân không bằng f/c.
\end{itemchoice}
}
\end{ex}

% ===== Câu 211 =====
% ID: L12C3B19-48-TL
% Mức: VDC
\dangbt{Mô tả sự hình thành và lan truyền của sóng điện từ}
\begin{bt}
% ID: L12C3B19-48-TL
% Mức: VDC
So sánh nhận định đúng “Sóng điện từ là điện từ trường lan truyền trong không gian và truyền được trong chân không.” với nhận định sai “Sóng điện từ bắt buộc cần môi trường vật chất để truyền.”; chỉ rõ căn cứ Vật lí.
\loigiai{
Cần nêu được: Sóng điện từ là điện từ trường lan truyền trong không gian và truyền được trong chân không. Khi vận dụng phải xác định đúng dữ kiện, phương chiều nếu có, điều kiện áp dụng, hệ đơn vị và kiểm tra tính hợp lí. Nhận định “Sóng điện từ bắt buộc cần môi trường vật chất để truyền.” sai.
}
\end{bt}

% ===== Câu 212 =====
% ID: L12C3B19-48-TLN
% Mức: VDC
\begin{ex}
% ID: L12C3B19-48-TLN
% Mức: VDC
Trong bốn phát biểu về Mô tả sự hình thành và lan truyền của sóng điện từ, có bao nhiêu phát biểu đúng? (Chỉ ghi một số nguyên.) (1) Khi giải cần xác định đúng dữ kiện, phương chiều, điều kiện áp dụng và đơn vị. (2) Sóng điện từ bắt buộc cần môi trường vật chất để truyền. (3) Phải dùng đúng định luật cho tình huống. (4) Sóng điện từ là điện từ trường lan truyền trong không gian và truyền được trong chân không.
\shortans{3}
\loigiai{
Đối chiếu với kiến thức: Sóng điện từ là điện từ trường lan truyền trong không gian và truyền được trong chân không. Có 3 phát biểu đúng.
}
\end{ex}

% ===== Câu 213 =====
% ID: L12C3B19-48-TN
% Mức: VD
\dangbt{Phân loại phổ sóng điện từ và nêu ứng dụng}
\begin{ex}
% ID: L12C3B19-48-TN
% Mức: VD
Cơ sở Vật lí đúng dùng để giải Phân loại phổ sóng điện từ và nêu ứng dụng?
\choice
{Có thể bỏ qua phương, chiều và đơn vị trong mọi trường hợp.}
{Kết quả không phụ thuộc dữ kiện và điều kiện áp dụng.}
{\True Phổ điện từ gồm sóng vô tuyến, vi ba, hồng ngoại, ánh sáng nhìn thấy, tử ngoại, tia $X$ và gamma.}
{Âm thanh là một miền của phổ sóng điện từ.}
\loigiai{
Phổ điện từ gồm sóng vô tuyến, vi ba, hồng ngoại, ánh sáng nhìn thấy, tử ngoại, tia $X$ và gamma. Vì vậy phương án $C$ đúng; các phương án còn lại sai.
}
\end{ex}

% ===== Câu 214 =====
% ID: L12C3B19-49-DS
% Mức: VD
\dangbt{Tính bước sóng, tần số và tốc độ truyền sóng điện từ}
\begin{ex}
% ID: L12C3B19-49-DS
% Mức: VD
Xét Tính bước sóng, tần số và tốc độ truyền sóng điện từ (Tình huống 4). Hãy xác định tính đúng, sai của bốn phát biểu sau.
\choiceTF
{Kết quả không phụ thuộc dữ kiện và có thể bỏ qua điều kiện.}
{\True Khi giải cần xác định đúng dữ kiện, phương chiều, điều kiện áp dụng và đơn vị.}
{Bước sóng trong chân không bằng f/c.}
{\True Trong chân không, bước sóng lambda = c/f với c xấp xỉ 3.10^ $8\,\text{m}$ /s.}
\loigiai{
\begin{itemchoice}
\item (Sai) Kết quả không phụ thuộc dữ kiện và có thể bỏ qua điều kiện. (Vì): Kết quả không phụ thuộc dữ kiện và có thể bỏ qua điều kiện. b) Đúng: Khi giải cần xác định đúng dữ kiện, phương chiều, điều kiện áp dụng và đơn vị. c) Sai: Bước sóng trong chân không bằng f/c. d) Đúng: Trong chân không, bước sóng lambda = c/f với c xấp xỉ 3.10^ $8\,\text{m}$ /s. Kiến thức cốt lõi: Trong chân không, bước sóng lambda = c/f với c xấp xỉ 3.10^ $8\,\text{m}$ /s
\item (Đúng) Khi giải cần xác định đúng dữ kiện, phương chiều, điều kiện áp dụng và đơn vị.
\item (Sai) Bước sóng trong chân không bằng f/c.
\item (Đúng) Trong chân không, bước sóng lambda = c/f với c xấp xỉ 3.10^ $8\,\text{m}$ /s.
\end{itemchoice}
}
\end{ex}

% ===== Câu 215 =====
% ID: L12C3B19-49-TL
% Mức: TH
\dangbt{Nhận biết các tính chất cơ bản của sóng điện từ}
\begin{bt}
% ID: L12C3B19-49-TL
% Mức: TH
Trình bày kiến thức cốt lõi của Nhận biết các tính chất cơ bản của sóng điện từ và nêu điều kiện áp dụng.
\loigiai{
Cần nêu được: Trong chân không, vectơ $E$ và $B$ vuông góc nhau và cùng vuông góc phương truyền sóng. Khi vận dụng phải xác định đúng dữ kiện, phương chiều nếu có, điều kiện áp dụng, hệ đơn vị và kiểm tra tính hợp lí. Nhận định “Trong sóng điện từ, $E$ và $B$ song song với phương truyền.” sai.
}
\end{bt}

% ===== Câu 216 =====
% ID: L12C3B19-49-TLN
% Mức: TH
\begin{ex}
% ID: L12C3B19-49-TLN
% Mức: TH
Trong bốn phát biểu về Nhận biết các tính chất cơ bản của sóng điện từ, có bao nhiêu phát biểu đúng? (Chỉ ghi một số nguyên.) (1) Trong chân không, vectơ $E$ và $B$ vuông góc nhau và cùng vuông góc phương truyền sóng. (2) Trong sóng điện từ, $E$ và $B$ song song với phương truyền. (3) Cần kiểm tra điều kiện áp dụng. (4) Có thể bỏ qua đơn vị.
\shortans{2}
\loigiai{
Đối chiếu với kiến thức: Trong chân không, vectơ $E$ và $B$ vuông góc nhau và cùng vuông góc phương truyền sóng. Có 2 phát biểu đúng.
}
\end{ex}

% ===== Câu 217 =====
% ID: L12C3B19-49-TN
% Mức: VD
\dangbt{Phân loại phổ sóng điện từ và nêu ứng dụng}
\begin{ex}
% ID: L12C3B19-49-TN
% Mức: VD
Trong một bài thực hành về Phân loại phổ sóng điện từ và nêu ứng dụng?
\choice
{Âm thanh là một miền của phổ sóng điện từ.}
{Có thể bỏ qua phương, chiều và đơn vị trong mọi trường hợp.}
{Kết quả không phụ thuộc dữ kiện và điều kiện áp dụng.}
{\True Phổ điện từ gồm sóng vô tuyến, vi ba, hồng ngoại, ánh sáng nhìn thấy, tử ngoại, tia $X$ và gamma.}
\loigiai{
Phổ điện từ gồm sóng vô tuyến, vi ba, hồng ngoại, ánh sáng nhìn thấy, tử ngoại, tia $X$ và gamma. Vì vậy phương án $D$ đúng; các phương án còn lại sai.
}
\end{ex}

% ===== Câu 218 =====
% ID: L12C3B19-50-DS
% Mức: VD
\dangbt{Tính bước sóng, tần số và tốc độ truyền sóng điện từ}
\begin{ex}
% ID: L12C3B19-50-DS
% Mức: VD
Xét Tính bước sóng, tần số và tốc độ truyền sóng điện từ (Tình huống 5). Hãy xác định tính đúng, sai của bốn phát biểu sau.
\choiceTF
{\True Trong chân không, bước sóng lambda = c/f với c xấp xỉ 3.10^ $8\,\text{m}$ /s.}
{Kết quả không phụ thuộc dữ kiện và có thể bỏ qua điều kiện.}
{Bước sóng trong chân không bằng f/c.}
{\True Khi giải cần xác định đúng dữ kiện, phương chiều, điều kiện áp dụng và đơn vị.}
\loigiai{
\begin{itemchoice}
\item (Đúng) Trong chân không, bước sóng lambda = c/f với c xấp xỉ 3.10^ $8\,\text{m}$ /s. (Vì): Trong chân không, bước sóng lambda = c/f với c xấp xỉ 3.10^ $8\,\text{m}$ /s. b) Sai: Kết quả không phụ thuộc dữ kiện và có thể bỏ qua điều kiện. c) Sai: Bước sóng trong chân không bằng f/c. d) Đúng: Khi giải cần xác định đúng dữ kiện, phương chiều, điều kiện áp dụng và đơn vị. Kiến thức cốt lõi: Trong chân không, bước sóng lambda = c/f với c xấp xỉ 3.10^ $8\,\text{m}$ /s
\item (Sai) Kết quả không phụ thuộc dữ kiện và có thể bỏ qua điều kiện.
\item (Sai) Bước sóng trong chân không bằng f/c.
\item (Đúng) Khi giải cần xác định đúng dữ kiện, phương chiều, điều kiện áp dụng và đơn vị.
\end{itemchoice}
}
\end{ex}

% ===== Câu 219 =====
% ID: L12C3B19-50-TL
% Mức: TH
\dangbt{Nhận biết các tính chất cơ bản của sóng điện từ}
\begin{bt}
% ID: L12C3B19-50-TL
% Mức: TH
Giải thích vì sao nhận định sau đúng: “Trong chân không, vectơ $E$ và $B$ vuông góc nhau và cùng vuông góc phương truyền sóng.”
\loigiai{
Cần nêu được: Trong chân không, vectơ $E$ và $B$ vuông góc nhau và cùng vuông góc phương truyền sóng. Khi vận dụng phải xác định đúng dữ kiện, phương chiều nếu có, điều kiện áp dụng, hệ đơn vị và kiểm tra tính hợp lí. Nhận định “Trong sóng điện từ, $E$ và $B$ song song với phương truyền.” sai.
}
\end{bt}

% ===== Câu 220 =====
% ID: L12C3B19-50-TLN
% Mức: TH
\begin{ex}
% ID: L12C3B19-50-TLN
% Mức: TH
Trong bốn phát biểu về Nhận biết các tính chất cơ bản của sóng điện từ, có bao nhiêu phát biểu đúng? (Chỉ ghi một số nguyên.) (1) Trong sóng điện từ, $E$ và $B$ song song với phương truyền. (2) Trong chân không, vectơ $E$ và $B$ vuông góc nhau và cùng vuông góc phương truyền sóng. (3) Kết quả phải phù hợp ý nghĩa vật lí. (4) Mọi đại lượng đều là vô hướng.
\shortans{1}
\loigiai{
Đối chiếu với kiến thức: Trong chân không, vectơ $E$ và $B$ vuông góc nhau và cùng vuông góc phương truyền sóng. Có 1 phát biểu đúng.
}
\end{ex}

% ===== Câu 221 =====
% ID: L12C3B19-50-TN
% Mức: VD
\dangbt{Phân loại phổ sóng điện từ và nêu ứng dụng}
\begin{ex}
% ID: L12C3B19-50-TN
% Mức: VD
Kết luận đúng khi vận dụng Phân loại phổ sóng điện từ và nêu ứng dụng?
\choice
{\True Phổ điện từ gồm sóng vô tuyến, vi ba, hồng ngoại, ánh sáng nhìn thấy, tử ngoại, tia $X$ và gamma.}
{Âm thanh là một miền của phổ sóng điện từ.}
{Có thể bỏ qua phương, chiều và đơn vị trong mọi trường hợp.}
{Kết quả không phụ thuộc dữ kiện và điều kiện áp dụng.}
\loigiai{
Phổ điện từ gồm sóng vô tuyến, vi ba, hồng ngoại, ánh sáng nhìn thấy, tử ngoại, tia $X$ và gamma. Vì vậy phương án $A$ đúng; các phương án còn lại sai.
}
\end{ex}

% ===== Câu 222 =====
% ID: L12C3B19-51-DS
% Mức: TH
\begin{ex}
% ID: L12C3B19-51-DS
% Mức: TH
Xét Phân loại phổ sóng điện từ và nêu ứng dụng (Tình huống 1). Hãy xác định tính đúng, sai của bốn phát biểu sau.
\choiceTF
{\True Phổ điện từ gồm sóng vô tuyến, vi ba, hồng ngoại, ánh sáng nhìn thấy, tử ngoại, tia $X$ và gamma.}
{Âm thanh là một miền của phổ sóng điện từ.}
{\True Khi giải cần xác định đúng dữ kiện, phương chiều, điều kiện áp dụng và đơn vị.}
{Kết quả không phụ thuộc dữ kiện và có thể bỏ qua điều kiện.}
\loigiai{
\begin{itemchoice}
\item (Đúng) Phổ điện từ gồm sóng vô tuyến, vi ba, hồng ngoại, ánh sáng nhìn thấy, tử ngoại, tia $X$ và gamma. (Vì): Phổ điện từ gồm sóng vô tuyến, vi ba, hồng ngoại, ánh sáng nhìn thấy, tử ngoại, tia $X$ và gamma. b) Sai: Âm thanh là một miền của phổ sóng điện từ. c) Đúng: Khi giải cần xác định đúng dữ kiện, phương chiều, điều kiện áp dụng và đơn vị. d) Sai: Kết quả không phụ thuộc dữ kiện và có thể bỏ qua điều kiện. Kiến thức cốt lõi: Phổ điện từ gồm sóng vô tuyến, vi ba, hồng ngoại, ánh sáng nhìn thấy, tử ngoại, tia $X$ và gamma
\item (Sai) Âm thanh là một miền của phổ sóng điện từ.
\item (Đúng) Khi giải cần xác định đúng dữ kiện, phương chiều, điều kiện áp dụng và đơn vị.
\item (Sai) Kết quả không phụ thuộc dữ kiện và có thể bỏ qua điều kiện.
\end{itemchoice}
}
\end{ex}

% ===== Câu 223 =====
% ID: L12C3B19-51-TL
% Mức: VD
\dangbt{Nhận biết các tính chất cơ bản của sóng điện từ}
\begin{bt}
% ID: L12C3B19-51-TL
% Mức: VD
Phân tích sai lầm trong nhận định: “Trong sóng điện từ, $E$ và $B$ song song với phương truyền.”
\loigiai{
Cần nêu được: Trong chân không, vectơ $E$ và $B$ vuông góc nhau và cùng vuông góc phương truyền sóng. Khi vận dụng phải xác định đúng dữ kiện, phương chiều nếu có, điều kiện áp dụng, hệ đơn vị và kiểm tra tính hợp lí. Nhận định “Trong sóng điện từ, $E$ và $B$ song song với phương truyền.” sai.
}
\end{bt}

% ===== Câu 224 =====
% ID: L12C3B19-51-TLN
% Mức: TH
\begin{ex}
% ID: L12C3B19-51-TLN
% Mức: TH
Trong bốn phát biểu về Nhận biết các tính chất cơ bản của sóng điện từ, có bao nhiêu phát biểu đúng? (Chỉ ghi một số nguyên.) (1) Trong chân không, vectơ $E$ và $B$ vuông góc nhau và cùng vuông góc phương truyền sóng. (2) Trong sóng điện từ, $E$ và $B$ song song với phương truyền. (3) Cần kiểm tra điều kiện áp dụng. (4) Có thể bỏ qua đơn vị.
\shortans{2}
\loigiai{
Đối chiếu với kiến thức: Trong chân không, vectơ $E$ và $B$ vuông góc nhau và cùng vuông góc phương truyền sóng. Có 2 phát biểu đúng.
}
\end{ex}

% ===== Câu 225 =====
% ID: L12C3B19-51-TN
% Mức: NB
\dangbt{Tính bước sóng, tần số và tốc độ truyền sóng điện từ}
\begin{ex}
% ID: L12C3B19-51-TN
% Mức: NB
Phát biểu nào sau đây đúng về Tính bước sóng, tần số và tốc độ truyền sóng điện từ?
\choice
{Bước sóng trong chân không bằng f/c.}
{Có thể bỏ qua phương, chiều và đơn vị trong mọi trường hợp.}
{Kết quả không phụ thuộc dữ kiện và điều kiện áp dụng.}
{\True Trong chân không, bước sóng lambda = c/f với c xấp xỉ 3.10^ $8\,\text{m}$ /s.}
\loigiai{
Trong chân không, bước sóng lambda = c/f với c xấp xỉ 3.10^ $8\,\text{m}$ /s. Vì vậy phương án $D$ đúng; các phương án còn lại sai.
}
\end{ex}

% ===== Câu 226 =====
% ID: L12C3B19-52-DS
% Mức: TH
\dangbt{Phân loại phổ sóng điện từ và nêu ứng dụng}
\begin{ex}
% ID: L12C3B19-52-DS
% Mức: TH
Xét Phân loại phổ sóng điện từ và nêu ứng dụng (Tình huống 2). Hãy xác định tính đúng, sai của bốn phát biểu sau.
\choiceTF
{Âm thanh là một miền của phổ sóng điện từ.}
{\True Phổ điện từ gồm sóng vô tuyến, vi ba, hồng ngoại, ánh sáng nhìn thấy, tử ngoại, tia $X$ và gamma.}
{Kết quả không phụ thuộc dữ kiện và có thể bỏ qua điều kiện.}
{\True Khi giải cần xác định đúng dữ kiện, phương chiều, điều kiện áp dụng và đơn vị.}
\loigiai{
\begin{itemchoice}
\item (Sai) Âm thanh là một miền của phổ sóng điện từ. (Vì): Âm thanh là một miền của phổ sóng điện từ. b) Đúng: Phổ điện từ gồm sóng vô tuyến, vi ba, hồng ngoại, ánh sáng nhìn thấy, tử ngoại, tia $X$ và gamma. c) Sai: Kết quả không phụ thuộc dữ kiện và có thể bỏ qua điều kiện. d) Đúng: Khi giải cần xác định đúng dữ kiện, phương chiều, điều kiện áp dụng và đơn vị. Kiến thức cốt lõi: Phổ điện từ gồm sóng vô tuyến, vi ba, hồng ngoại, ánh sáng nhìn thấy, tử ngoại, tia $X$ và gamma
\item (Đúng) Phổ điện từ gồm sóng vô tuyến, vi ba, hồng ngoại, ánh sáng nhìn thấy, tử ngoại, tia $X$ và gamma.
\item (Sai) Kết quả không phụ thuộc dữ kiện và có thể bỏ qua điều kiện.
\item (Đúng) Khi giải cần xác định đúng dữ kiện, phương chiều, điều kiện áp dụng và đơn vị.
\end{itemchoice}
}
\end{ex}

% ===== Câu 227 =====
% ID: L12C3B19-52-TL
% Mức: VD
\dangbt{Nhận biết các tính chất cơ bản của sóng điện từ}
\begin{bt}
% ID: L12C3B19-52-TL
% Mức: VD
Đề xuất các bước giải một bài toán thuộc dạng Nhận biết các tính chất cơ bản của sóng điện từ.
\loigiai{
Cần nêu được: Trong chân không, vectơ $E$ và $B$ vuông góc nhau và cùng vuông góc phương truyền sóng. Khi vận dụng phải xác định đúng dữ kiện, phương chiều nếu có, điều kiện áp dụng, hệ đơn vị và kiểm tra tính hợp lí. Nhận định “Trong sóng điện từ, $E$ và $B$ song song với phương truyền.” sai.
}
\end{bt}

% ===== Câu 228 =====
% ID: L12C3B19-52-TLN
% Mức: VD
\begin{ex}
% ID: L12C3B19-52-TLN
% Mức: VD
Trong bốn phát biểu về Nhận biết các tính chất cơ bản của sóng điện từ, có bao nhiêu phát biểu đúng? (Chỉ ghi một số nguyên.) (1) Trong sóng điện từ, $E$ và $B$ song song với phương truyền. (2) Trong chân không, vectơ $E$ và $B$ vuông góc nhau và cùng vuông góc phương truyền sóng. (3) Kết quả phải phù hợp ý nghĩa vật lí. (4) Mọi đại lượng đều là vô hướng.
\shortans{2}
\loigiai{
Đối chiếu với kiến thức: Trong chân không, vectơ $E$ và $B$ vuông góc nhau và cùng vuông góc phương truyền sóng. Có 2 phát biểu đúng.
}
\end{ex}

% ===== Câu 229 =====
% ID: L12C3B19-52-TN
% Mức: NB
\dangbt{Tính bước sóng, tần số và tốc độ truyền sóng điện từ}
\begin{ex}
% ID: L12C3B19-52-TN
% Mức: NB
Trong nội dung Tính bước sóng, tần số và tốc độ truyền sóng điện từ?
\choice
{\True Trong chân không, bước sóng lambda = c/f với c xấp xỉ 3.10^ $8\,\text{m}$ /s.}
{Bước sóng trong chân không bằng f/c.}
{Có thể bỏ qua phương, chiều và đơn vị trong mọi trường hợp.}
{Kết quả không phụ thuộc dữ kiện và điều kiện áp dụng.}
\loigiai{
Trong chân không, bước sóng lambda = c/f với c xấp xỉ 3.10^ $8\,\text{m}$ /s. Vì vậy phương án $A$ đúng; các phương án còn lại sai.
}
\end{ex}

% ===== Câu 230 =====
% ID: L12C3B19-53-DS
% Mức: VD
\dangbt{Phân loại phổ sóng điện từ và nêu ứng dụng}
\begin{ex}
% ID: L12C3B19-53-DS
% Mức: VD
Xét Phân loại phổ sóng điện từ và nêu ứng dụng (Tình huống 3). Hãy xác định tính đúng, sai của bốn phát biểu sau.
\choiceTF
{\True Khi giải cần xác định đúng dữ kiện, phương chiều, điều kiện áp dụng và đơn vị.}
{Kết quả không phụ thuộc dữ kiện và có thể bỏ qua điều kiện.}
{\True Phổ điện từ gồm sóng vô tuyến, vi ba, hồng ngoại, ánh sáng nhìn thấy, tử ngoại, tia $X$ và gamma.}
{Âm thanh là một miền của phổ sóng điện từ.}
\loigiai{
\begin{itemchoice}
\item (Đúng) Khi giải cần xác định đúng dữ kiện, phương chiều, điều kiện áp dụng và đơn vị. (Vì): Khi giải cần xác định đúng dữ kiện, phương chiều, điều kiện áp dụng và đơn vị. b) Sai: Kết quả không phụ thuộc dữ kiện và có thể bỏ qua điều kiện. c) Đúng: Phổ điện từ gồm sóng vô tuyến, vi ba, hồng ngoại, ánh sáng nhìn thấy, tử ngoại, tia $X$ và gamma. d) Sai: Âm thanh là một miền của phổ sóng điện từ. Kiến thức cốt lõi: Phổ điện từ gồm sóng vô tuyến, vi ba, hồng ngoại, ánh sáng nhìn thấy, tử ngoại, tia $X$ và gamma
\item (Sai) Kết quả không phụ thuộc dữ kiện và có thể bỏ qua điều kiện.
\item (Đúng) Phổ điện từ gồm sóng vô tuyến, vi ba, hồng ngoại, ánh sáng nhìn thấy, tử ngoại, tia $X$ và gamma.
\item (Sai) Âm thanh là một miền của phổ sóng điện từ.
\end{itemchoice}
}
\end{ex}

% ===== Câu 231 =====
% ID: L12C3B19-53-TL
% Mức: VD
\dangbt{Nhận biết các tính chất cơ bản của sóng điện từ}
\begin{bt}
% ID: L12C3B19-53-TL
% Mức: VD
Nêu một tình huống thực tiễn liên quan đến Nhận biết các tính chất cơ bản của sóng điện từ và giải thích bằng kiến thức Vật lí.
\loigiai{
Cần nêu được: Trong chân không, vectơ $E$ và $B$ vuông góc nhau và cùng vuông góc phương truyền sóng. Khi vận dụng phải xác định đúng dữ kiện, phương chiều nếu có, điều kiện áp dụng, hệ đơn vị và kiểm tra tính hợp lí. Nhận định “Trong sóng điện từ, $E$ và $B$ song song với phương truyền.” sai.
}
\end{bt}

% ===== Câu 232 =====
% ID: L12C3B19-53-TLN
% Mức: VD
\begin{ex}
% ID: L12C3B19-53-TLN
% Mức: VD
Trong bốn phát biểu về Nhận biết các tính chất cơ bản của sóng điện từ, có bao nhiêu phát biểu đúng? (Chỉ ghi một số nguyên.) (1) Trong chân không, vectơ $E$ và $B$ vuông góc nhau và cùng vuông góc phương truyền sóng. (2) Trong sóng điện từ, $E$ và $B$ song song với phương truyền. (3) Cần kiểm tra điều kiện áp dụng. (4) Có thể bỏ qua đơn vị.
\shortans{2}
\loigiai{
Đối chiếu với kiến thức: Trong chân không, vectơ $E$ và $B$ vuông góc nhau và cùng vuông góc phương truyền sóng. Có 2 phát biểu đúng.
}
\end{ex}

% ===== Câu 233 =====
% ID: L12C3B19-53-TN
% Mức: NB
\dangbt{Tính bước sóng, tần số và tốc độ truyền sóng điện từ}
\begin{ex}
% ID: L12C3B19-53-TN
% Mức: NB
Tình huống 3: chọn kết luận chính xác về Tính bước sóng, tần số và tốc độ truyền sóng điện từ?
\choice
{Kết quả không phụ thuộc dữ kiện và điều kiện áp dụng.}
{\True Trong chân không, bước sóng lambda = c/f với c xấp xỉ 3.10^ $8\,\text{m}$ /s.}
{Bước sóng trong chân không bằng f/c.}
{Có thể bỏ qua phương, chiều và đơn vị trong mọi trường hợp.}
\loigiai{
Trong chân không, bước sóng lambda = c/f với c xấp xỉ 3.10^ $8\,\text{m}$ /s. Vì vậy phương án $B$ đúng; các phương án còn lại sai.
}
\end{ex}

% ===== Câu 234 =====
% ID: L12C3B19-54-DS
% Mức: VD
\dangbt{Phân loại phổ sóng điện từ và nêu ứng dụng}
\begin{ex}
% ID: L12C3B19-54-DS
% Mức: VD
Xét Phân loại phổ sóng điện từ và nêu ứng dụng (Tình huống 4). Hãy xác định tính đúng, sai của bốn phát biểu sau.
\choiceTF
{Kết quả không phụ thuộc dữ kiện và có thể bỏ qua điều kiện.}
{\True Khi giải cần xác định đúng dữ kiện, phương chiều, điều kiện áp dụng và đơn vị.}
{Âm thanh là một miền của phổ sóng điện từ.}
{\True Phổ điện từ gồm sóng vô tuyến, vi ba, hồng ngoại, ánh sáng nhìn thấy, tử ngoại, tia $X$ và gamma.}
\loigiai{
\begin{itemchoice}
\item (Sai) Kết quả không phụ thuộc dữ kiện và có thể bỏ qua điều kiện. (Vì): Kết quả không phụ thuộc dữ kiện và có thể bỏ qua điều kiện. b) Đúng: Khi giải cần xác định đúng dữ kiện, phương chiều, điều kiện áp dụng và đơn vị. c) Sai: Âm thanh là một miền của phổ sóng điện từ. d) Đúng: Phổ điện từ gồm sóng vô tuyến, vi ba, hồng ngoại, ánh sáng nhìn thấy, tử ngoại, tia $X$ và gamma. Kiến thức cốt lõi: Phổ điện từ gồm sóng vô tuyến, vi ba, hồng ngoại, ánh sáng nhìn thấy, tử ngoại, tia $X$ và gamma
\item (Đúng) Khi giải cần xác định đúng dữ kiện, phương chiều, điều kiện áp dụng và đơn vị.
\item (Sai) Âm thanh là một miền của phổ sóng điện từ.
\item (Đúng) Phổ điện từ gồm sóng vô tuyến, vi ba, hồng ngoại, ánh sáng nhìn thấy, tử ngoại, tia $X$ và gamma.
\end{itemchoice}
}
\end{ex}

% ===== Câu 235 =====
% ID: L12C3B19-54-TL
% Mức: VDC
\dangbt{Nhận biết các tính chất cơ bản của sóng điện từ}
\begin{bt}
% ID: L12C3B19-54-TL
% Mức: VDC
So sánh nhận định đúng “Trong chân không, vectơ $E$ và $B$ vuông góc nhau và cùng vuông góc phương truyền sóng.” với nhận định sai “Trong sóng điện từ, $E$ và $B$ song song với phương truyền.”; chỉ rõ căn cứ Vật lí.
\loigiai{
Cần nêu được: Trong chân không, vectơ $E$ và $B$ vuông góc nhau và cùng vuông góc phương truyền sóng. Khi vận dụng phải xác định đúng dữ kiện, phương chiều nếu có, điều kiện áp dụng, hệ đơn vị và kiểm tra tính hợp lí. Nhận định “Trong sóng điện từ, $E$ và $B$ song song với phương truyền.” sai.
}
\end{bt}

% ===== Câu 236 =====
% ID: L12C3B19-54-TLN
% Mức: VDC
\begin{ex}
% ID: L12C3B19-54-TLN
% Mức: VDC
Trong bốn phát biểu về Nhận biết các tính chất cơ bản của sóng điện từ, có bao nhiêu phát biểu đúng? (Chỉ ghi một số nguyên.) (1) Khi giải cần xác định đúng dữ kiện, phương chiều, điều kiện áp dụng và đơn vị. (2) Trong sóng điện từ, $E$ và $B$ song song với phương truyền. (3) Phải dùng đúng định luật cho tình huống. (4) Trong chân không, vectơ $E$ và $B$ vuông góc nhau và cùng vuông góc phương truyền sóng.
\shortans{3}
\loigiai{
Đối chiếu với kiến thức: Trong chân không, vectơ $E$ và $B$ vuông góc nhau và cùng vuông góc phương truyền sóng. Có 3 phát biểu đúng.
}
\end{ex}

% ===== Câu 237 =====
% ID: L12C3B19-54-TN
% Mức: TH
\dangbt{Tính bước sóng, tần số và tốc độ truyền sóng điện từ}
\begin{ex}
% ID: L12C3B19-54-TN
% Mức: TH
Nội dung nào mô tả đúng nhất Tính bước sóng, tần số và tốc độ truyền sóng điện từ?
\choice
{Có thể bỏ qua phương, chiều và đơn vị trong mọi trường hợp.}
{Kết quả không phụ thuộc dữ kiện và điều kiện áp dụng.}
{\True Trong chân không, bước sóng lambda = c/f với c xấp xỉ 3.10^ $8\,\text{m}$ /s.}
{Bước sóng trong chân không bằng f/c.}
\loigiai{
Trong chân không, bước sóng lambda = c/f với c xấp xỉ 3.10^ $8\,\text{m}$ /s. Vì vậy phương án $C$ đúng; các phương án còn lại sai.
}
\end{ex}

% ===== Câu 238 =====
% ID: L12C3B19-55-DS
% Mức: VD
\dangbt{Phân loại phổ sóng điện từ và nêu ứng dụng}
\begin{ex}
% ID: L12C3B19-55-DS
% Mức: VD
Xét Phân loại phổ sóng điện từ và nêu ứng dụng (Tình huống 5). Hãy xác định tính đúng, sai của bốn phát biểu sau.
\choiceTF
{\True Phổ điện từ gồm sóng vô tuyến, vi ba, hồng ngoại, ánh sáng nhìn thấy, tử ngoại, tia $X$ và gamma.}
{Kết quả không phụ thuộc dữ kiện và có thể bỏ qua điều kiện.}
{Âm thanh là một miền của phổ sóng điện từ.}
{\True Khi giải cần xác định đúng dữ kiện, phương chiều, điều kiện áp dụng và đơn vị.}
\loigiai{
\begin{itemchoice}
\item (Đúng) Phổ điện từ gồm sóng vô tuyến, vi ba, hồng ngoại, ánh sáng nhìn thấy, tử ngoại, tia $X$ và gamma. (Vì): Phổ điện từ gồm sóng vô tuyến, vi ba, hồng ngoại, ánh sáng nhìn thấy, tử ngoại, tia $X$ và gamma. b) Sai: Kết quả không phụ thuộc dữ kiện và có thể bỏ qua điều kiện. c) Sai: Âm thanh là một miền của phổ sóng điện từ. d) Đúng: Khi giải cần xác định đúng dữ kiện, phương chiều, điều kiện áp dụng và đơn vị. Kiến thức cốt lõi: Phổ điện từ gồm sóng vô tuyến, vi ba, hồng ngoại, ánh sáng nhìn thấy, tử ngoại, tia $X$ và gamma
\item (Sai) Kết quả không phụ thuộc dữ kiện và có thể bỏ qua điều kiện.
\item (Sai) Âm thanh là một miền của phổ sóng điện từ.
\item (Đúng) Khi giải cần xác định đúng dữ kiện, phương chiều, điều kiện áp dụng và đơn vị.
\end{itemchoice}
}
\end{ex}

% ===== Câu 239 =====
% ID: L12C3B19-55-TL
% Mức: TH
\dangbt{Tính bước sóng, tần số và tốc độ truyền sóng điện từ}
\begin{bt}
% ID: L12C3B19-55-TL
% Mức: TH
Trình bày kiến thức cốt lõi của Tính bước sóng, tần số và tốc độ truyền sóng điện từ và nêu điều kiện áp dụng.
\loigiai{
Cần nêu được: Trong chân không, bước sóng lambda = c/f với c xấp xỉ 3.10^ $8\,\text{m}$ /s. Khi vận dụng phải xác định đúng dữ kiện, phương chiều nếu có, điều kiện áp dụng, hệ đơn vị và kiểm tra tính hợp lí. Nhận định “Bước sóng trong chân không bằng f/c.” sai.
}
\end{bt}

% ===== Câu 240 =====
% ID: L12C3B19-55-TLN
% Mức: TH
\begin{ex}
% ID: L12C3B19-55-TLN
% Mức: TH
Trong bốn phát biểu về Tính bước sóng, tần số và tốc độ truyền sóng điện từ, có bao nhiêu phát biểu đúng? (Chỉ ghi một số nguyên.) (1) Trong chân không, bước sóng lambda = c/f với c xấp xỉ 3.10^ $8\,\text{m}$ /s. (2) Bước sóng trong chân không bằng f/c. (3) Cần kiểm tra điều kiện áp dụng. (4) Có thể bỏ qua đơn vị.
\shortans{2}
\loigiai{
Đối chiếu với kiến thức: Trong chân không, bước sóng lambda = c/f với c xấp xỉ 3.10^ $8\,\text{m}$ /s. Có 2 phát biểu đúng.
}
\end{ex}

% ===== Câu 241 =====
% ID: L12C3B19-55-TN
% Mức: TH
\begin{ex}
% ID: L12C3B19-55-TN
% Mức: TH
Khi phân tích Tính bước sóng, tần số và tốc độ truyền sóng điện từ?
\choice
{Bước sóng trong chân không bằng f/c.}
{Có thể bỏ qua phương, chiều và đơn vị trong mọi trường hợp.}
{Kết quả không phụ thuộc dữ kiện và điều kiện áp dụng.}
{\True Trong chân không, bước sóng lambda = c/f với c xấp xỉ 3.10^ $8\,\text{m}$ /s.}
\loigiai{
Trong chân không, bước sóng lambda = c/f với c xấp xỉ 3.10^ $8\,\text{m}$ /s. Vì vậy phương án $D$ đúng; các phương án còn lại sai.
}
\end{ex}

% ===== Câu 242 =====
% ID: L12C3B19-56-DS
% Mức: TH
\dangbt{Giải thích truyền thông tin bằng sóng điện từ và các tình huống thực tiễn}
\begin{ex}
% ID: L12C3B19-56-DS
% Mức: TH
Xét Giải thích truyền thông tin bằng sóng điện từ và các tình huống thực tiễn (Tình huống 1). Hãy xác định tính đúng, sai của bốn phát biểu sau.
\choiceTF
{\True Thông tin được ghép lên sóng mang bằng biến điệu, truyền đi và được máy thu giải điều chế.}
{Máy thu nhận thông tin mà không cần chọn sóng hay giải điều chế.}
{\True Khi giải cần xác định đúng dữ kiện, phương chiều, điều kiện áp dụng và đơn vị.}
{Kết quả không phụ thuộc dữ kiện và có thể bỏ qua điều kiện.}
\loigiai{
\begin{itemchoice}
\item (Đúng) Thông tin được ghép lên sóng mang bằng biến điệu, truyền đi và được máy thu giải điều chế. (Vì): Thông tin được ghép lên sóng mang bằng biến điệu, truyền đi và được máy thu giải điều chế. b) Sai: Máy thu nhận thông tin mà không cần chọn sóng hay giải điều chế. c) Đúng: Khi giải cần xác định đúng dữ kiện, phương chiều, điều kiện áp dụng và đơn vị. d) Sai: Kết quả không phụ thuộc dữ kiện và có thể bỏ qua điều kiện. Kiến thức cốt lõi: Thông tin được ghép lên sóng mang bằng biến điệu, truyền đi và được máy thu giải điều chế
\item (Sai) Máy thu nhận thông tin mà không cần chọn sóng hay giải điều chế.
\item (Đúng) Khi giải cần xác định đúng dữ kiện, phương chiều, điều kiện áp dụng và đơn vị.
\item (Sai) Kết quả không phụ thuộc dữ kiện và có thể bỏ qua điều kiện.
\end{itemchoice}
}
\end{ex}

% ===== Câu 243 =====
% ID: L12C3B19-56-TL
% Mức: TH
\dangbt{Tính bước sóng, tần số và tốc độ truyền sóng điện từ}
\begin{bt}
% ID: L12C3B19-56-TL
% Mức: TH
Giải thích vì sao nhận định sau đúng: “Trong chân không, bước sóng lambda = c/f với c xấp xỉ 3.10^ $8\,\text{m}$ /s.”
\loigiai{
Cần nêu được: Trong chân không, bước sóng lambda = c/f với c xấp xỉ 3.10^ $8\,\text{m}$ /s. Khi vận dụng phải xác định đúng dữ kiện, phương chiều nếu có, điều kiện áp dụng, hệ đơn vị và kiểm tra tính hợp lí. Nhận định “Bước sóng trong chân không bằng f/c.” sai.
}
\end{bt}

% ===== Câu 244 =====
% ID: L12C3B19-56-TLN
% Mức: TH
\begin{ex}
% ID: L12C3B19-56-TLN
% Mức: TH
Trong bốn phát biểu về Tính bước sóng, tần số và tốc độ truyền sóng điện từ, có bao nhiêu phát biểu đúng? (Chỉ ghi một số nguyên.) (1) Bước sóng trong chân không bằng f/c. (2) Trong chân không, bước sóng lambda = c/f với c xấp xỉ 3.10^ $8\,\text{m}$ /s. (3) Kết quả phải phù hợp ý nghĩa vật lí. (4) Mọi đại lượng đều là vô hướng.
\shortans{1}
\loigiai{
Đối chiếu với kiến thức: Trong chân không, bước sóng lambda = c/f với c xấp xỉ 3.10^ $8\,\text{m}$ /s. Có 1 phát biểu đúng.
}
\end{ex}

% ===== Câu 245 =====
% ID: L12C3B19-56-TN
% Mức: TH
\begin{ex}
% ID: L12C3B19-56-TN
% Mức: TH
Một học sinh ôn tập Tính bước sóng, tần số và tốc độ truyền sóng điện từ?
\choice
{\True Trong chân không, bước sóng lambda = c/f với c xấp xỉ 3.10^ $8\,\text{m}$ /s.}
{Bước sóng trong chân không bằng f/c.}
{Có thể bỏ qua phương, chiều và đơn vị trong mọi trường hợp.}
{Kết quả không phụ thuộc dữ kiện và điều kiện áp dụng.}
\loigiai{
Trong chân không, bước sóng lambda = c/f với c xấp xỉ 3.10^ $8\,\text{m}$ /s. Vì vậy phương án $A$ đúng; các phương án còn lại sai.
}
\end{ex}

% ===== Câu 246 =====
% ID: L12C3B19-57-DS
% Mức: TH
\dangbt{Giải thích truyền thông tin bằng sóng điện từ và các tình huống thực tiễn}
\begin{ex}
% ID: L12C3B19-57-DS
% Mức: TH
Xét Giải thích truyền thông tin bằng sóng điện từ và các tình huống thực tiễn (Tình huống 2). Hãy xác định tính đúng, sai của bốn phát biểu sau.
\choiceTF
{Máy thu nhận thông tin mà không cần chọn sóng hay giải điều chế.}
{\True Thông tin được ghép lên sóng mang bằng biến điệu, truyền đi và được máy thu giải điều chế.}
{Kết quả không phụ thuộc dữ kiện và có thể bỏ qua điều kiện.}
{\True Khi giải cần xác định đúng dữ kiện, phương chiều, điều kiện áp dụng và đơn vị.}
\loigiai{
\begin{itemchoice}
\item (Sai) Máy thu nhận thông tin mà không cần chọn sóng hay giải điều chế. (Vì): Máy thu nhận thông tin mà không cần chọn sóng hay giải điều chế. b) Đúng: Thông tin được ghép lên sóng mang bằng biến điệu, truyền đi và được máy thu giải điều chế. c) Sai: Kết quả không phụ thuộc dữ kiện và có thể bỏ qua điều kiện. d) Đúng: Khi giải cần xác định đúng dữ kiện, phương chiều, điều kiện áp dụng và đơn vị. Kiến thức cốt lõi: Thông tin được ghép lên sóng mang bằng biến điệu, truyền đi và được máy thu giải điều chế
\item (Đúng) Thông tin được ghép lên sóng mang bằng biến điệu, truyền đi và được máy thu giải điều chế.
\item (Sai) Kết quả không phụ thuộc dữ kiện và có thể bỏ qua điều kiện.
\item (Đúng) Khi giải cần xác định đúng dữ kiện, phương chiều, điều kiện áp dụng và đơn vị.
\end{itemchoice}
}
\end{ex}

% ===== Câu 247 =====
% ID: L12C3B19-57-TL
% Mức: VD
\dangbt{Tính bước sóng, tần số và tốc độ truyền sóng điện từ}
\begin{bt}
% ID: L12C3B19-57-TL
% Mức: VD
Phân tích sai lầm trong nhận định: “Bước sóng trong chân không bằng f/c.”
\loigiai{
Cần nêu được: Trong chân không, bước sóng lambda = c/f với c xấp xỉ 3.10^ $8\,\text{m}$ /s. Khi vận dụng phải xác định đúng dữ kiện, phương chiều nếu có, điều kiện áp dụng, hệ đơn vị và kiểm tra tính hợp lí. Nhận định “Bước sóng trong chân không bằng f/c.” sai.
}
\end{bt}

% ===== Câu 248 =====
% ID: L12C3B19-57-TLN
% Mức: TH
\begin{ex}
% ID: L12C3B19-57-TLN
% Mức: TH
Trong bốn phát biểu về Tính bước sóng, tần số và tốc độ truyền sóng điện từ, có bao nhiêu phát biểu đúng? (Chỉ ghi một số nguyên.) (1) Trong chân không, bước sóng lambda = c/f với c xấp xỉ 3.10^ $8\,\text{m}$ /s. (2) Bước sóng trong chân không bằng f/c. (3) Cần kiểm tra điều kiện áp dụng. (4) Có thể bỏ qua đơn vị.
\shortans{2}
\loigiai{
Đối chiếu với kiến thức: Trong chân không, bước sóng lambda = c/f với c xấp xỉ 3.10^ $8\,\text{m}$ /s. Có 2 phát biểu đúng.
}
\end{ex}

% ===== Câu 249 =====
% ID: L12C3B19-57-TN
% Mức: TH
\begin{ex}
% ID: L12C3B19-57-TN
% Mức: TH
Chọn phát biểu không trái với kiến thức về Tính bước sóng, tần số và tốc độ truyền sóng điện từ?
\choice
{Kết quả không phụ thuộc dữ kiện và điều kiện áp dụng.}
{\True Trong chân không, bước sóng lambda = c/f với c xấp xỉ 3.10^ $8\,\text{m}$ /s.}
{Bước sóng trong chân không bằng f/c.}
{Có thể bỏ qua phương, chiều và đơn vị trong mọi trường hợp.}
\loigiai{
Trong chân không, bước sóng lambda = c/f với c xấp xỉ 3.10^ $8\,\text{m}$ /s. Vì vậy phương án $B$ đúng; các phương án còn lại sai.
}
\end{ex}

% ===== Câu 250 =====
% ID: L12C3B19-58-DS
% Mức: VD
\dangbt{Giải thích truyền thông tin bằng sóng điện từ và các tình huống thực tiễn}
\begin{ex}
% ID: L12C3B19-58-DS
% Mức: VD
Xét Giải thích truyền thông tin bằng sóng điện từ và các tình huống thực tiễn (Tình huống 3). Hãy xác định tính đúng, sai của bốn phát biểu sau.
\choiceTF
{\True Khi giải cần xác định đúng dữ kiện, phương chiều, điều kiện áp dụng và đơn vị.}
{Kết quả không phụ thuộc dữ kiện và có thể bỏ qua điều kiện.}
{\True Thông tin được ghép lên sóng mang bằng biến điệu, truyền đi và được máy thu giải điều chế.}
{Máy thu nhận thông tin mà không cần chọn sóng hay giải điều chế.}
\loigiai{
\begin{itemchoice}
\item (Đúng) Khi giải cần xác định đúng dữ kiện, phương chiều, điều kiện áp dụng và đơn vị. (Vì): Khi giải cần xác định đúng dữ kiện, phương chiều, điều kiện áp dụng và đơn vị. b) Sai: Kết quả không phụ thuộc dữ kiện và có thể bỏ qua điều kiện. c) Đúng: Thông tin được ghép lên sóng mang bằng biến điệu, truyền đi và được máy thu giải điều chế. d) Sai: Máy thu nhận thông tin mà không cần chọn sóng hay giải điều chế. Kiến thức cốt lõi: Thông tin được ghép lên sóng mang bằng biến điệu, truyền đi và được máy thu giải điều chế
\item (Sai) Kết quả không phụ thuộc dữ kiện và có thể bỏ qua điều kiện.
\item (Đúng) Thông tin được ghép lên sóng mang bằng biến điệu, truyền đi và được máy thu giải điều chế.
\item (Sai) Máy thu nhận thông tin mà không cần chọn sóng hay giải điều chế.
\end{itemchoice}
}
\end{ex}

% ===== Câu 251 =====
% ID: L12C3B19-58-TL
% Mức: VD
\dangbt{Tính bước sóng, tần số và tốc độ truyền sóng điện từ}
\begin{bt}
% ID: L12C3B19-58-TL
% Mức: VD
Đề xuất các bước giải một bài toán thuộc dạng Tính bước sóng, tần số và tốc độ truyền sóng điện từ.
\loigiai{
Cần nêu được: Trong chân không, bước sóng lambda = c/f với c xấp xỉ 3.10^ $8\,\text{m}$ /s. Khi vận dụng phải xác định đúng dữ kiện, phương chiều nếu có, điều kiện áp dụng, hệ đơn vị và kiểm tra tính hợp lí. Nhận định “Bước sóng trong chân không bằng f/c.” sai.
}
\end{bt}

% ===== Câu 252 =====
% ID: L12C3B19-58-TLN
% Mức: VD
\begin{ex}
% ID: L12C3B19-58-TLN
% Mức: VD
Trong bốn phát biểu về Tính bước sóng, tần số và tốc độ truyền sóng điện từ, có bao nhiêu phát biểu đúng? (Chỉ ghi một số nguyên.) (1) Bước sóng trong chân không bằng f/c. (2) Trong chân không, bước sóng lambda = c/f với c xấp xỉ 3.10^ $8\,\text{m}$ /s. (3) Kết quả phải phù hợp ý nghĩa vật lí. (4) Mọi đại lượng đều là vô hướng.
\shortans{2}
\loigiai{
Đối chiếu với kiến thức: Trong chân không, bước sóng lambda = c/f với c xấp xỉ 3.10^ $8\,\text{m}$ /s. Có 2 phát biểu đúng.
}
\end{ex}

% ===== Câu 253 =====
% ID: L12C3B19-58-TN
% Mức: VD
\begin{ex}
% ID: L12C3B19-58-TN
% Mức: VD
Cơ sở Vật lí đúng dùng để giải Tính bước sóng, tần số và tốc độ truyền sóng điện từ?
\choice
{Có thể bỏ qua phương, chiều và đơn vị trong mọi trường hợp.}
{Kết quả không phụ thuộc dữ kiện và điều kiện áp dụng.}
{\True Trong chân không, bước sóng lambda = c/f với c xấp xỉ 3.10^ $8\,\text{m}$ /s.}
{Bước sóng trong chân không bằng f/c.}
\loigiai{
Trong chân không, bước sóng lambda = c/f với c xấp xỉ 3.10^ $8\,\text{m}$ /s. Vì vậy phương án $C$ đúng; các phương án còn lại sai.
}
\end{ex}

% ===== Câu 254 =====
% ID: L12C3B19-59-DS
% Mức: VD
\dangbt{Giải thích truyền thông tin bằng sóng điện từ và các tình huống thực tiễn}
\begin{ex}
% ID: L12C3B19-59-DS
% Mức: VD
Xét Giải thích truyền thông tin bằng sóng điện từ và các tình huống thực tiễn (Tình huống 4). Hãy xác định tính đúng, sai của bốn phát biểu sau.
\choiceTF
{Kết quả không phụ thuộc dữ kiện và có thể bỏ qua điều kiện.}
{\True Khi giải cần xác định đúng dữ kiện, phương chiều, điều kiện áp dụng và đơn vị.}
{Máy thu nhận thông tin mà không cần chọn sóng hay giải điều chế.}
{\True Thông tin được ghép lên sóng mang bằng biến điệu, truyền đi và được máy thu giải điều chế.}
\loigiai{
\begin{itemchoice}
\item (Sai) Kết quả không phụ thuộc dữ kiện và có thể bỏ qua điều kiện. (Vì): Kết quả không phụ thuộc dữ kiện và có thể bỏ qua điều kiện. b) Đúng: Khi giải cần xác định đúng dữ kiện, phương chiều, điều kiện áp dụng và đơn vị. c) Sai: Máy thu nhận thông tin mà không cần chọn sóng hay giải điều chế. d) Đúng: Thông tin được ghép lên sóng mang bằng biến điệu, truyền đi và được máy thu giải điều chế. Kiến thức cốt lõi: Thông tin được ghép lên sóng mang bằng biến điệu, truyền đi và được máy thu giải điều chế
\item (Đúng) Khi giải cần xác định đúng dữ kiện, phương chiều, điều kiện áp dụng và đơn vị.
\item (Sai) Máy thu nhận thông tin mà không cần chọn sóng hay giải điều chế.
\item (Đúng) Thông tin được ghép lên sóng mang bằng biến điệu, truyền đi và được máy thu giải điều chế.
\end{itemchoice}
}
\end{ex}

% ===== Câu 255 =====
% ID: L12C3B19-59-TL
% Mức: VD
\dangbt{Tính bước sóng, tần số và tốc độ truyền sóng điện từ}
\begin{bt}
% ID: L12C3B19-59-TL
% Mức: VD
Nêu một tình huống thực tiễn liên quan đến Tính bước sóng, tần số và tốc độ truyền sóng điện từ và giải thích bằng kiến thức Vật lí.
\loigiai{
Cần nêu được: Trong chân không, bước sóng lambda = c/f với c xấp xỉ 3.10^ $8\,\text{m}$ /s. Khi vận dụng phải xác định đúng dữ kiện, phương chiều nếu có, điều kiện áp dụng, hệ đơn vị và kiểm tra tính hợp lí. Nhận định “Bước sóng trong chân không bằng f/c.” sai.
}
\end{bt}

% ===== Câu 256 =====
% ID: L12C3B19-59-TLN
% Mức: VD
\begin{ex}
% ID: L12C3B19-59-TLN
% Mức: VD
Trong bốn phát biểu về Tính bước sóng, tần số và tốc độ truyền sóng điện từ, có bao nhiêu phát biểu đúng? (Chỉ ghi một số nguyên.) (1) Trong chân không, bước sóng lambda = c/f với c xấp xỉ 3.10^ $8\,\text{m}$ /s. (2) Bước sóng trong chân không bằng f/c. (3) Cần kiểm tra điều kiện áp dụng. (4) Có thể bỏ qua đơn vị.
\shortans{2}
\loigiai{
Đối chiếu với kiến thức: Trong chân không, bước sóng lambda = c/f với c xấp xỉ 3.10^ $8\,\text{m}$ /s. Có 2 phát biểu đúng.
}
\end{ex}

% ===== Câu 257 =====
% ID: L12C3B19-59-TN
% Mức: VD
\begin{ex}
% ID: L12C3B19-59-TN
% Mức: VD
Trong một bài thực hành về Tính bước sóng, tần số và tốc độ truyền sóng điện từ?
\choice
{Bước sóng trong chân không bằng f/c.}
{Có thể bỏ qua phương, chiều và đơn vị trong mọi trường hợp.}
{Kết quả không phụ thuộc dữ kiện và điều kiện áp dụng.}
{\True Trong chân không, bước sóng lambda = c/f với c xấp xỉ 3.10^ $8\,\text{m}$ /s.}
\loigiai{
Trong chân không, bước sóng lambda = c/f với c xấp xỉ 3.10^ $8\,\text{m}$ /s. Vì vậy phương án $D$ đúng; các phương án còn lại sai.
}
\end{ex}

% ===== Câu 258 =====
% ID: L12C3B19-60-DS
% Mức: VD
\dangbt{Giải thích truyền thông tin bằng sóng điện từ và các tình huống thực tiễn}
\begin{ex}
% ID: L12C3B19-60-DS
% Mức: VD
Xét Giải thích truyền thông tin bằng sóng điện từ và các tình huống thực tiễn (Tình huống 5). Hãy xác định tính đúng, sai của bốn phát biểu sau.
\choiceTF
{\True Thông tin được ghép lên sóng mang bằng biến điệu, truyền đi và được máy thu giải điều chế.}
{Kết quả không phụ thuộc dữ kiện và có thể bỏ qua điều kiện.}
{Máy thu nhận thông tin mà không cần chọn sóng hay giải điều chế.}
{\True Khi giải cần xác định đúng dữ kiện, phương chiều, điều kiện áp dụng và đơn vị.}
\loigiai{
\begin{itemchoice}
\item (Đúng) Thông tin được ghép lên sóng mang bằng biến điệu, truyền đi và được máy thu giải điều chế. (Vì): Thông tin được ghép lên sóng mang bằng biến điệu, truyền đi và được máy thu giải điều chế. b) Sai: Kết quả không phụ thuộc dữ kiện và có thể bỏ qua điều kiện. c) Sai: Máy thu nhận thông tin mà không cần chọn sóng hay giải điều chế. d) Đúng: Khi giải cần xác định đúng dữ kiện, phương chiều, điều kiện áp dụng và đơn vị. Kiến thức cốt lõi: Thông tin được ghép lên sóng mang bằng biến điệu, truyền đi và được máy thu giải điều chế
\item (Sai) Kết quả không phụ thuộc dữ kiện và có thể bỏ qua điều kiện.
\item (Sai) Máy thu nhận thông tin mà không cần chọn sóng hay giải điều chế.
\item (Đúng) Khi giải cần xác định đúng dữ kiện, phương chiều, điều kiện áp dụng và đơn vị.
\end{itemchoice}
}
\end{ex}

% ===== Câu 259 =====
% ID: L12C3B19-60-TL
% Mức: VDC
\dangbt{Tính bước sóng, tần số và tốc độ truyền sóng điện từ}
\begin{bt}
% ID: L12C3B19-60-TL
% Mức: VDC
So sánh nhận định đúng “Trong chân không, bước sóng lambda = c/f với c xấp xỉ 3.10^ $8\,\text{m}$ /s.” với nhận định sai “Bước sóng trong chân không bằng f/c.”; chỉ rõ căn cứ Vật lí.
\loigiai{
Cần nêu được: Trong chân không, bước sóng lambda = c/f với c xấp xỉ 3.10^ $8\,\text{m}$ /s. Khi vận dụng phải xác định đúng dữ kiện, phương chiều nếu có, điều kiện áp dụng, hệ đơn vị và kiểm tra tính hợp lí. Nhận định “Bước sóng trong chân không bằng f/c.” sai.
}
\end{bt}

% ===== Câu 260 =====
% ID: L12C3B19-60-TLN
% Mức: VDC
\begin{ex}
% ID: L12C3B19-60-TLN
% Mức: VDC
Trong bốn phát biểu về Tính bước sóng, tần số và tốc độ truyền sóng điện từ, có bao nhiêu phát biểu đúng? (Chỉ ghi một số nguyên.) (1) Khi giải cần xác định đúng dữ kiện, phương chiều, điều kiện áp dụng và đơn vị. (2) Bước sóng trong chân không bằng f/c. (3) Phải dùng đúng định luật cho tình huống. (4) Trong chân không, bước sóng lambda = c/f với c xấp xỉ 3.10^ $8\,\text{m}$ /s.
\shortans{3}
\loigiai{
Đối chiếu với kiến thức: Trong chân không, bước sóng lambda = c/f với c xấp xỉ 3.10^ $8\,\text{m}$ /s. Có 3 phát biểu đúng.
}
\end{ex}

% ===== Câu 261 =====
% ID: L12C3B19-60-TN
% Mức: VD
\begin{ex}
% ID: L12C3B19-60-TN
% Mức: VD
Kết luận đúng khi vận dụng Tính bước sóng, tần số và tốc độ truyền sóng điện từ?
\choice
{\True Trong chân không, bước sóng lambda = c/f với c xấp xỉ 3.10^ $8\,\text{m}$ /s.}
{Bước sóng trong chân không bằng f/c.}
{Có thể bỏ qua phương, chiều và đơn vị trong mọi trường hợp.}
{Kết quả không phụ thuộc dữ kiện và điều kiện áp dụng.}
\loigiai{
Trong chân không, bước sóng lambda = c/f với c xấp xỉ 3.10^ $8\,\text{m}$ /s. Vì vậy phương án $A$ đúng; các phương án còn lại sai.
}
\end{ex}

% ===== Câu 262 =====
% ID: L12C3B19-61-TL
% Mức: TH
\dangbt{Phân loại phổ sóng điện từ và nêu ứng dụng}
\begin{bt}
% ID: L12C3B19-61-TL
% Mức: TH
Trình bày kiến thức cốt lõi của Phân loại phổ sóng điện từ và nêu ứng dụng và nêu điều kiện áp dụng.
\loigiai{
Cần nêu được: Phổ điện từ gồm sóng vô tuyến, vi ba, hồng ngoại, ánh sáng nhìn thấy, tử ngoại, tia $X$ và gamma. Khi vận dụng phải xác định đúng dữ kiện, phương chiều nếu có, điều kiện áp dụng, hệ đơn vị và kiểm tra tính hợp lí. Nhận định “Âm thanh là một miền của phổ sóng điện từ.” sai.
}
\end{bt}

% ===== Câu 263 =====
% ID: L12C3B19-61-TLN
% Mức: TH
\begin{ex}
% ID: L12C3B19-61-TLN
% Mức: TH
Trong bốn phát biểu về Phân loại phổ sóng điện từ và nêu ứng dụng, có bao nhiêu phát biểu đúng? (Chỉ ghi một số nguyên.) (1) Phổ điện từ gồm sóng vô tuyến, vi ba, hồng ngoại, ánh sáng nhìn thấy, tử ngoại, tia $X$ và gamma. (2) Âm thanh là một miền của phổ sóng điện từ. (3) Cần kiểm tra điều kiện áp dụng. (4) Có thể bỏ qua đơn vị.
\shortans{2}
\loigiai{
Đối chiếu với kiến thức: Phổ điện từ gồm sóng vô tuyến, vi ba, hồng ngoại, ánh sáng nhìn thấy, tử ngoại, tia $X$ và gamma. Có 2 phát biểu đúng.
}
\end{ex}

% ===== Câu 264 =====
% ID: L12C3B19-61-TN
% Mức: NB
\dangbt{Nhận biết điện trường biến thiên, từ trường biến thiên và điện từ trường}
\begin{ex}
% ID: L12C3B19-61-TN
% Mức: NB
Phát biểu nào sau đây đúng về Nhận biết điện trường biến thiên, từ trường biến thiên và điện từ trường?
\choice
{Điện trường và từ trường biến thiên tồn tại hoàn toàn độc lập.}
{Có thể bỏ qua phương, chiều và đơn vị trong mọi trường hợp.}
{Kết quả không phụ thuộc dữ kiện và điều kiện áp dụng.}
{\True Điện trường biến thiên sinh ra từ trường biến thiên và từ trường biến thiên sinh ra điện trường xoáy.}
\loigiai{
Điện trường biến thiên sinh ra từ trường biến thiên và từ trường biến thiên sinh ra điện trường xoáy. Vì vậy phương án $D$ đúng; các phương án còn lại sai.
}
\end{ex}

% ===== Câu 265 =====
% ID: L12C3B19-62-TL
% Mức: TH
\dangbt{Phân loại phổ sóng điện từ và nêu ứng dụng}
\begin{bt}
% ID: L12C3B19-62-TL
% Mức: TH
Giải thích vì sao nhận định sau đúng: “Phổ điện từ gồm sóng vô tuyến, vi ba, hồng ngoại, ánh sáng nhìn thấy, tử ngoại, tia $X$ và gamma.”
\loigiai{
Cần nêu được: Phổ điện từ gồm sóng vô tuyến, vi ba, hồng ngoại, ánh sáng nhìn thấy, tử ngoại, tia $X$ và gamma. Khi vận dụng phải xác định đúng dữ kiện, phương chiều nếu có, điều kiện áp dụng, hệ đơn vị và kiểm tra tính hợp lí. Nhận định “Âm thanh là một miền của phổ sóng điện từ.” sai.
}
\end{bt}

% ===== Câu 266 =====
% ID: L12C3B19-62-TLN
% Mức: TH
\begin{ex}
% ID: L12C3B19-62-TLN
% Mức: TH
Trong bốn phát biểu về Phân loại phổ sóng điện từ và nêu ứng dụng, có bao nhiêu phát biểu đúng? (Chỉ ghi một số nguyên.) (1) Âm thanh là một miền của phổ sóng điện từ. (2) Phổ điện từ gồm sóng vô tuyến, vi ba, hồng ngoại, ánh sáng nhìn thấy, tử ngoại, tia $X$ và gamma. (3) Kết quả phải phù hợp ý nghĩa vật lí. (4) Mọi đại lượng đều là vô hướng.
\shortans{1}
\loigiai{
Đối chiếu với kiến thức: Phổ điện từ gồm sóng vô tuyến, vi ba, hồng ngoại, ánh sáng nhìn thấy, tử ngoại, tia $X$ và gamma. Có 1 phát biểu đúng.
}
\end{ex}

% ===== Câu 267 =====
% ID: L12C3B19-62-TN
% Mức: NB
\dangbt{Nhận biết điện trường biến thiên, từ trường biến thiên và điện từ trường}
\begin{ex}
% ID: L12C3B19-62-TN
% Mức: NB
Trong nội dung Nhận biết điện trường biến thiên, từ trường biến thiên và điện từ trường?
\choice
{\True Điện trường biến thiên sinh ra từ trường biến thiên và từ trường biến thiên sinh ra điện trường xoáy.}
{Điện trường và từ trường biến thiên tồn tại hoàn toàn độc lập.}
{Có thể bỏ qua phương, chiều và đơn vị trong mọi trường hợp.}
{Kết quả không phụ thuộc dữ kiện và điều kiện áp dụng.}
\loigiai{
Điện trường biến thiên sinh ra từ trường biến thiên và từ trường biến thiên sinh ra điện trường xoáy. Vì vậy phương án $A$ đúng; các phương án còn lại sai.
}
\end{ex}

% ===== Câu 268 =====
% ID: L12C3B19-63-TL
% Mức: VD
\dangbt{Phân loại phổ sóng điện từ và nêu ứng dụng}
\begin{bt}
% ID: L12C3B19-63-TL
% Mức: VD
Phân tích sai lầm trong nhận định: “Âm thanh là một miền của phổ sóng điện từ.”
\loigiai{
Cần nêu được: Phổ điện từ gồm sóng vô tuyến, vi ba, hồng ngoại, ánh sáng nhìn thấy, tử ngoại, tia $X$ và gamma. Khi vận dụng phải xác định đúng dữ kiện, phương chiều nếu có, điều kiện áp dụng, hệ đơn vị và kiểm tra tính hợp lí. Nhận định “Âm thanh là một miền của phổ sóng điện từ.” sai.
}
\end{bt}

% ===== Câu 269 =====
% ID: L12C3B19-63-TLN
% Mức: TH
\begin{ex}
% ID: L12C3B19-63-TLN
% Mức: TH
Trong bốn phát biểu về Phân loại phổ sóng điện từ và nêu ứng dụng, có bao nhiêu phát biểu đúng? (Chỉ ghi một số nguyên.) (1) Phổ điện từ gồm sóng vô tuyến, vi ba, hồng ngoại, ánh sáng nhìn thấy, tử ngoại, tia $X$ và gamma. (2) Âm thanh là một miền của phổ sóng điện từ. (3) Cần kiểm tra điều kiện áp dụng. (4) Có thể bỏ qua đơn vị.
\shortans{2}
\loigiai{
Đối chiếu với kiến thức: Phổ điện từ gồm sóng vô tuyến, vi ba, hồng ngoại, ánh sáng nhìn thấy, tử ngoại, tia $X$ và gamma. Có 2 phát biểu đúng.
}
\end{ex}

% ===== Câu 270 =====
% ID: L12C3B19-63-TN
% Mức: NB
\dangbt{Nhận biết điện trường biến thiên, từ trường biến thiên và điện từ trường}
\begin{ex}
% ID: L12C3B19-63-TN
% Mức: NB
Tình huống 3: chọn kết luận chính xác về Nhận biết điện trường biến thiên, từ trường biến thiên và điện từ trường?
\choice
{Kết quả không phụ thuộc dữ kiện và điều kiện áp dụng.}
{\True Điện trường biến thiên sinh ra từ trường biến thiên và từ trường biến thiên sinh ra điện trường xoáy.}
{Điện trường và từ trường biến thiên tồn tại hoàn toàn độc lập.}
{Có thể bỏ qua phương, chiều và đơn vị trong mọi trường hợp.}
\loigiai{
Điện trường biến thiên sinh ra từ trường biến thiên và từ trường biến thiên sinh ra điện trường xoáy. Vì vậy phương án $B$ đúng; các phương án còn lại sai.
}
\end{ex}

% ===== Câu 271 =====
% ID: L12C3B19-64-TL
% Mức: VD
\dangbt{Phân loại phổ sóng điện từ và nêu ứng dụng}
\begin{bt}
% ID: L12C3B19-64-TL
% Mức: VD
Đề xuất các bước giải một bài toán thuộc dạng Phân loại phổ sóng điện từ và nêu ứng dụng.
\loigiai{
Cần nêu được: Phổ điện từ gồm sóng vô tuyến, vi ba, hồng ngoại, ánh sáng nhìn thấy, tử ngoại, tia $X$ và gamma. Khi vận dụng phải xác định đúng dữ kiện, phương chiều nếu có, điều kiện áp dụng, hệ đơn vị và kiểm tra tính hợp lí. Nhận định “Âm thanh là một miền của phổ sóng điện từ.” sai.
}
\end{bt}

% ===== Câu 272 =====
% ID: L12C3B19-64-TLN
% Mức: VD
\begin{ex}
% ID: L12C3B19-64-TLN
% Mức: VD
Trong bốn phát biểu về Phân loại phổ sóng điện từ và nêu ứng dụng, có bao nhiêu phát biểu đúng? (Chỉ ghi một số nguyên.) (1) Âm thanh là một miền của phổ sóng điện từ. (2) Phổ điện từ gồm sóng vô tuyến, vi ba, hồng ngoại, ánh sáng nhìn thấy, tử ngoại, tia $X$ và gamma. (3) Kết quả phải phù hợp ý nghĩa vật lí. (4) Mọi đại lượng đều là vô hướng.
\shortans{2}
\loigiai{
Đối chiếu với kiến thức: Phổ điện từ gồm sóng vô tuyến, vi ba, hồng ngoại, ánh sáng nhìn thấy, tử ngoại, tia $X$ và gamma. Có 2 phát biểu đúng.
}
\end{ex}

% ===== Câu 273 =====
% ID: L12C3B19-64-TN
% Mức: TH
\dangbt{Nhận biết điện trường biến thiên, từ trường biến thiên và điện từ trường}
\begin{ex}
% ID: L12C3B19-64-TN
% Mức: TH
Nội dung nào mô tả đúng nhất Nhận biết điện trường biến thiên, từ trường biến thiên và điện từ trường?
\choice
{Có thể bỏ qua phương, chiều và đơn vị trong mọi trường hợp.}
{Kết quả không phụ thuộc dữ kiện và điều kiện áp dụng.}
{\True Điện trường biến thiên sinh ra từ trường biến thiên và từ trường biến thiên sinh ra điện trường xoáy.}
{Điện trường và từ trường biến thiên tồn tại hoàn toàn độc lập.}
\loigiai{
Điện trường biến thiên sinh ra từ trường biến thiên và từ trường biến thiên sinh ra điện trường xoáy. Vì vậy phương án $C$ đúng; các phương án còn lại sai.
}
\end{ex}

% ===== Câu 274 =====
% ID: L12C3B19-65-TL
% Mức: VD
\dangbt{Phân loại phổ sóng điện từ và nêu ứng dụng}
\begin{bt}
% ID: L12C3B19-65-TL
% Mức: VD
Nêu một tình huống thực tiễn liên quan đến Phân loại phổ sóng điện từ và nêu ứng dụng và giải thích bằng kiến thức Vật lí.
\loigiai{
Cần nêu được: Phổ điện từ gồm sóng vô tuyến, vi ba, hồng ngoại, ánh sáng nhìn thấy, tử ngoại, tia $X$ và gamma. Khi vận dụng phải xác định đúng dữ kiện, phương chiều nếu có, điều kiện áp dụng, hệ đơn vị và kiểm tra tính hợp lí. Nhận định “Âm thanh là một miền của phổ sóng điện từ.” sai.
}
\end{bt}

% ===== Câu 275 =====
% ID: L12C3B19-65-TLN
% Mức: VD
\begin{ex}
% ID: L12C3B19-65-TLN
% Mức: VD
Trong bốn phát biểu về Phân loại phổ sóng điện từ và nêu ứng dụng, có bao nhiêu phát biểu đúng? (Chỉ ghi một số nguyên.) (1) Phổ điện từ gồm sóng vô tuyến, vi ba, hồng ngoại, ánh sáng nhìn thấy, tử ngoại, tia $X$ và gamma. (2) Âm thanh là một miền của phổ sóng điện từ. (3) Cần kiểm tra điều kiện áp dụng. (4) Có thể bỏ qua đơn vị.
\shortans{2}
\loigiai{
Đối chiếu với kiến thức: Phổ điện từ gồm sóng vô tuyến, vi ba, hồng ngoại, ánh sáng nhìn thấy, tử ngoại, tia $X$ và gamma. Có 2 phát biểu đúng.
}
\end{ex}

% ===== Câu 276 =====
% ID: L12C3B19-65-TN
% Mức: TH
\dangbt{Nhận biết điện trường biến thiên, từ trường biến thiên và điện từ trường}
\begin{ex}
% ID: L12C3B19-65-TN
% Mức: TH
Khi phân tích Nhận biết điện trường biến thiên, từ trường biến thiên và điện từ trường?
\choice
{Điện trường và từ trường biến thiên tồn tại hoàn toàn độc lập.}
{Có thể bỏ qua phương, chiều và đơn vị trong mọi trường hợp.}
{Kết quả không phụ thuộc dữ kiện và điều kiện áp dụng.}
{\True Điện trường biến thiên sinh ra từ trường biến thiên và từ trường biến thiên sinh ra điện trường xoáy.}
\loigiai{
Điện trường biến thiên sinh ra từ trường biến thiên và từ trường biến thiên sinh ra điện trường xoáy. Vì vậy phương án $D$ đúng; các phương án còn lại sai.
}
\end{ex}

% ===== Câu 277 =====
% ID: L12C3B19-66-TL
% Mức: VDC
\dangbt{Phân loại phổ sóng điện từ và nêu ứng dụng}
\begin{bt}
% ID: L12C3B19-66-TL
% Mức: VDC
So sánh nhận định đúng “Phổ điện từ gồm sóng vô tuyến, vi ba, hồng ngoại, ánh sáng nhìn thấy, tử ngoại, tia $X$ và gamma.” với nhận định sai “Âm thanh là một miền của phổ sóng điện từ.”; chỉ rõ căn cứ Vật lí.
\loigiai{
Cần nêu được: Phổ điện từ gồm sóng vô tuyến, vi ba, hồng ngoại, ánh sáng nhìn thấy, tử ngoại, tia $X$ và gamma. Khi vận dụng phải xác định đúng dữ kiện, phương chiều nếu có, điều kiện áp dụng, hệ đơn vị và kiểm tra tính hợp lí. Nhận định “Âm thanh là một miền của phổ sóng điện từ.” sai.
}
\end{bt}

% ===== Câu 278 =====
% ID: L12C3B19-66-TLN
% Mức: VDC
\begin{ex}
% ID: L12C3B19-66-TLN
% Mức: VDC
Trong bốn phát biểu về Phân loại phổ sóng điện từ và nêu ứng dụng, có bao nhiêu phát biểu đúng? (Chỉ ghi một số nguyên.) (1) Khi giải cần xác định đúng dữ kiện, phương chiều, điều kiện áp dụng và đơn vị. (2) Âm thanh là một miền của phổ sóng điện từ. (3) Phải dùng đúng định luật cho tình huống. (4) Phổ điện từ gồm sóng vô tuyến, vi ba, hồng ngoại, ánh sáng nhìn thấy, tử ngoại, tia $X$ và gamma.
\shortans{3}
\loigiai{
Đối chiếu với kiến thức: Phổ điện từ gồm sóng vô tuyến, vi ba, hồng ngoại, ánh sáng nhìn thấy, tử ngoại, tia $X$ và gamma. Có 3 phát biểu đúng.
}
\end{ex}

% ===== Câu 279 =====
% ID: L12C3B19-66-TN
% Mức: TH
\dangbt{Nhận biết điện trường biến thiên, từ trường biến thiên và điện từ trường}
\begin{ex}
% ID: L12C3B19-66-TN
% Mức: TH
Một học sinh ôn tập Nhận biết điện trường biến thiên, từ trường biến thiên và điện từ trường?
\choice
{\True Điện trường biến thiên sinh ra từ trường biến thiên và từ trường biến thiên sinh ra điện trường xoáy.}
{Điện trường và từ trường biến thiên tồn tại hoàn toàn độc lập.}
{Có thể bỏ qua phương, chiều và đơn vị trong mọi trường hợp.}
{Kết quả không phụ thuộc dữ kiện và điều kiện áp dụng.}
\loigiai{
Điện trường biến thiên sinh ra từ trường biến thiên và từ trường biến thiên sinh ra điện trường xoáy. Vì vậy phương án $A$ đúng; các phương án còn lại sai.
}
\end{ex}

% ===== Câu 280 =====
% ID: L12C3B19-67-TL
% Mức: TH
\dangbt{Giải thích truyền thông tin bằng sóng điện từ và các tình huống thực tiễn}
\begin{bt}
% ID: L12C3B19-67-TL
% Mức: TH
Trình bày kiến thức cốt lõi của Giải thích truyền thông tin bằng sóng điện từ và các tình huống thực tiễn và nêu điều kiện áp dụng.
\loigiai{
Cần nêu được: Thông tin được ghép lên sóng mang bằng biến điệu, truyền đi và được máy thu giải điều chế. Khi vận dụng phải xác định đúng dữ kiện, phương chiều nếu có, điều kiện áp dụng, hệ đơn vị và kiểm tra tính hợp lí. Nhận định “Máy thu nhận thông tin mà không cần chọn sóng hay giải điều chế.” sai.
}
\end{bt}

% ===== Câu 281 =====
% ID: L12C3B19-67-TLN
% Mức: TH
\begin{ex}
% ID: L12C3B19-67-TLN
% Mức: TH
Trong bốn phát biểu về Giải thích truyền thông tin bằng sóng điện từ và các tình huống thực tiễn, có bao nhiêu phát biểu đúng? (Chỉ ghi một số nguyên.) (1) Thông tin được ghép lên sóng mang bằng biến điệu, truyền đi và được máy thu giải điều chế. (2) Máy thu nhận thông tin mà không cần chọn sóng hay giải điều chế. (3) Cần kiểm tra điều kiện áp dụng. (4) Có thể bỏ qua đơn vị.
\shortans{2}
\loigiai{
Đối chiếu với kiến thức: Thông tin được ghép lên sóng mang bằng biến điệu, truyền đi và được máy thu giải điều chế. Có 2 phát biểu đúng.
}
\end{ex}

% ===== Câu 282 =====
% ID: L12C3B19-67-TN
% Mức: TH
\dangbt{Nhận biết điện trường biến thiên, từ trường biến thiên và điện từ trường}
\begin{ex}
% ID: L12C3B19-67-TN
% Mức: TH
Chọn phát biểu không trái với kiến thức về Nhận biết điện trường biến thiên, từ trường biến thiên và điện từ trường?
\choice
{Kết quả không phụ thuộc dữ kiện và điều kiện áp dụng.}
{\True Điện trường biến thiên sinh ra từ trường biến thiên và từ trường biến thiên sinh ra điện trường xoáy.}
{Điện trường và từ trường biến thiên tồn tại hoàn toàn độc lập.}
{Có thể bỏ qua phương, chiều và đơn vị trong mọi trường hợp.}
\loigiai{
Điện trường biến thiên sinh ra từ trường biến thiên và từ trường biến thiên sinh ra điện trường xoáy. Vì vậy phương án $B$ đúng; các phương án còn lại sai.
}
\end{ex}

% ===== Câu 283 =====
% ID: L12C3B19-68-TL
% Mức: TH
\dangbt{Giải thích truyền thông tin bằng sóng điện từ và các tình huống thực tiễn}
\begin{bt}
% ID: L12C3B19-68-TL
% Mức: TH
Giải thích vì sao nhận định sau đúng: “Thông tin được ghép lên sóng mang bằng biến điệu, truyền đi và được máy thu giải điều chế.”
\loigiai{
Cần nêu được: Thông tin được ghép lên sóng mang bằng biến điệu, truyền đi và được máy thu giải điều chế. Khi vận dụng phải xác định đúng dữ kiện, phương chiều nếu có, điều kiện áp dụng, hệ đơn vị và kiểm tra tính hợp lí. Nhận định “Máy thu nhận thông tin mà không cần chọn sóng hay giải điều chế.” sai.
}
\end{bt}

% ===== Câu 284 =====
% ID: L12C3B19-68-TLN
% Mức: TH
\begin{ex}
% ID: L12C3B19-68-TLN
% Mức: TH
Trong bốn phát biểu về Giải thích truyền thông tin bằng sóng điện từ và các tình huống thực tiễn, có bao nhiêu phát biểu đúng? (Chỉ ghi một số nguyên.) (1) Máy thu nhận thông tin mà không cần chọn sóng hay giải điều chế. (2) Thông tin được ghép lên sóng mang bằng biến điệu, truyền đi và được máy thu giải điều chế. (3) Kết quả phải phù hợp ý nghĩa vật lí. (4) Mọi đại lượng đều là vô hướng.
\shortans{1}
\loigiai{
Đối chiếu với kiến thức: Thông tin được ghép lên sóng mang bằng biến điệu, truyền đi và được máy thu giải điều chế. Có 1 phát biểu đúng.
}
\end{ex}

% ===== Câu 285 =====
% ID: L12C3B19-68-TN
% Mức: VD
\dangbt{Nhận biết điện trường biến thiên, từ trường biến thiên và điện từ trường}
\begin{ex}
% ID: L12C3B19-68-TN
% Mức: VD
Cơ sở Vật lí đúng dùng để giải Nhận biết điện trường biến thiên, từ trường biến thiên và điện từ trường?
\choice
{Có thể bỏ qua phương, chiều và đơn vị trong mọi trường hợp.}
{Kết quả không phụ thuộc dữ kiện và điều kiện áp dụng.}
{\True Điện trường biến thiên sinh ra từ trường biến thiên và từ trường biến thiên sinh ra điện trường xoáy.}
{Điện trường và từ trường biến thiên tồn tại hoàn toàn độc lập.}
\loigiai{
Điện trường biến thiên sinh ra từ trường biến thiên và từ trường biến thiên sinh ra điện trường xoáy. Vì vậy phương án $C$ đúng; các phương án còn lại sai.
}
\end{ex}

% ===== Câu 286 =====
% ID: L12C3B19-69-TL
% Mức: VD
\dangbt{Giải thích truyền thông tin bằng sóng điện từ và các tình huống thực tiễn}
\begin{bt}
% ID: L12C3B19-69-TL
% Mức: VD
Phân tích sai lầm trong nhận định: “Máy thu nhận thông tin mà không cần chọn sóng hay giải điều chế.”
\loigiai{
Cần nêu được: Thông tin được ghép lên sóng mang bằng biến điệu, truyền đi và được máy thu giải điều chế. Khi vận dụng phải xác định đúng dữ kiện, phương chiều nếu có, điều kiện áp dụng, hệ đơn vị và kiểm tra tính hợp lí. Nhận định “Máy thu nhận thông tin mà không cần chọn sóng hay giải điều chế.” sai.
}
\end{bt}

% ===== Câu 287 =====
% ID: L12C3B19-69-TLN
% Mức: TH
\begin{ex}
% ID: L12C3B19-69-TLN
% Mức: TH
Trong bốn phát biểu về Giải thích truyền thông tin bằng sóng điện từ và các tình huống thực tiễn, có bao nhiêu phát biểu đúng? (Chỉ ghi một số nguyên.) (1) Thông tin được ghép lên sóng mang bằng biến điệu, truyền đi và được máy thu giải điều chế. (2) Máy thu nhận thông tin mà không cần chọn sóng hay giải điều chế. (3) Cần kiểm tra điều kiện áp dụng. (4) Có thể bỏ qua đơn vị.
\shortans{2}
\loigiai{
Đối chiếu với kiến thức: Thông tin được ghép lên sóng mang bằng biến điệu, truyền đi và được máy thu giải điều chế. Có 2 phát biểu đúng.
}
\end{ex}

% ===== Câu 288 =====
% ID: L12C3B19-69-TN
% Mức: VD
\dangbt{Nhận biết điện trường biến thiên, từ trường biến thiên và điện từ trường}
\begin{ex}
% ID: L12C3B19-69-TN
% Mức: VD
Trong một bài thực hành về Nhận biết điện trường biến thiên, từ trường biến thiên và điện từ trường?
\choice
{Điện trường và từ trường biến thiên tồn tại hoàn toàn độc lập.}
{Có thể bỏ qua phương, chiều và đơn vị trong mọi trường hợp.}
{Kết quả không phụ thuộc dữ kiện và điều kiện áp dụng.}
{\True Điện trường biến thiên sinh ra từ trường biến thiên và từ trường biến thiên sinh ra điện trường xoáy.}
\loigiai{
Điện trường biến thiên sinh ra từ trường biến thiên và từ trường biến thiên sinh ra điện trường xoáy. Vì vậy phương án $D$ đúng; các phương án còn lại sai.
}
\end{ex}

% ===== Câu 289 =====
% ID: L12C3B19-70-TL
% Mức: VD
\dangbt{Giải thích truyền thông tin bằng sóng điện từ và các tình huống thực tiễn}
\begin{bt}
% ID: L12C3B19-70-TL
% Mức: VD
Đề xuất các bước giải một bài toán thuộc dạng Giải thích truyền thông tin bằng sóng điện từ và các tình huống thực tiễn.
\loigiai{
Cần nêu được: Thông tin được ghép lên sóng mang bằng biến điệu, truyền đi và được máy thu giải điều chế. Khi vận dụng phải xác định đúng dữ kiện, phương chiều nếu có, điều kiện áp dụng, hệ đơn vị và kiểm tra tính hợp lí. Nhận định “Máy thu nhận thông tin mà không cần chọn sóng hay giải điều chế.” sai.
}
\end{bt}

% ===== Câu 290 =====
% ID: L12C3B19-70-TLN
% Mức: VD
\begin{ex}
% ID: L12C3B19-70-TLN
% Mức: VD
Trong bốn phát biểu về Giải thích truyền thông tin bằng sóng điện từ và các tình huống thực tiễn, có bao nhiêu phát biểu đúng? (Chỉ ghi một số nguyên.) (1) Máy thu nhận thông tin mà không cần chọn sóng hay giải điều chế. (2) Thông tin được ghép lên sóng mang bằng biến điệu, truyền đi và được máy thu giải điều chế. (3) Kết quả phải phù hợp ý nghĩa vật lí. (4) Mọi đại lượng đều là vô hướng.
\shortans{2}
\loigiai{
Đối chiếu với kiến thức: Thông tin được ghép lên sóng mang bằng biến điệu, truyền đi và được máy thu giải điều chế. Có 2 phát biểu đúng.
}
\end{ex}

% ===== Câu 291 =====
% ID: L12C3B19-70-TN
% Mức: VD
\dangbt{Nhận biết điện trường biến thiên, từ trường biến thiên và điện từ trường}
\begin{ex}
% ID: L12C3B19-70-TN
% Mức: VD
Kết luận đúng khi vận dụng Nhận biết điện trường biến thiên, từ trường biến thiên và điện từ trường?
\choice
{\True Điện trường biến thiên sinh ra từ trường biến thiên và từ trường biến thiên sinh ra điện trường xoáy.}
{Điện trường và từ trường biến thiên tồn tại hoàn toàn độc lập.}
{Có thể bỏ qua phương, chiều và đơn vị trong mọi trường hợp.}
{Kết quả không phụ thuộc dữ kiện và điều kiện áp dụng.}
\loigiai{
Điện trường biến thiên sinh ra từ trường biến thiên và từ trường biến thiên sinh ra điện trường xoáy. Vì vậy phương án $A$ đúng; các phương án còn lại sai.
}
\end{ex}

% ===== Câu 292 =====
% ID: L12C3B19-71-TL
% Mức: VD
\dangbt{Giải thích truyền thông tin bằng sóng điện từ và các tình huống thực tiễn}
\begin{bt}
% ID: L12C3B19-71-TL
% Mức: VD
Nêu một tình huống thực tiễn liên quan đến Giải thích truyền thông tin bằng sóng điện từ và các tình huống thực tiễn và giải thích bằng kiến thức Vật lí.
\loigiai{
Cần nêu được: Thông tin được ghép lên sóng mang bằng biến điệu, truyền đi và được máy thu giải điều chế. Khi vận dụng phải xác định đúng dữ kiện, phương chiều nếu có, điều kiện áp dụng, hệ đơn vị và kiểm tra tính hợp lí. Nhận định “Máy thu nhận thông tin mà không cần chọn sóng hay giải điều chế.” sai.
}
\end{bt}

% ===== Câu 293 =====
% ID: L12C3B19-71-TLN
% Mức: VD
\begin{ex}
% ID: L12C3B19-71-TLN
% Mức: VD
Trong bốn phát biểu về Giải thích truyền thông tin bằng sóng điện từ và các tình huống thực tiễn, có bao nhiêu phát biểu đúng? (Chỉ ghi một số nguyên.) (1) Thông tin được ghép lên sóng mang bằng biến điệu, truyền đi và được máy thu giải điều chế. (2) Máy thu nhận thông tin mà không cần chọn sóng hay giải điều chế. (3) Cần kiểm tra điều kiện áp dụng. (4) Có thể bỏ qua đơn vị.
\shortans{2}
\loigiai{
Đối chiếu với kiến thức: Thông tin được ghép lên sóng mang bằng biến điệu, truyền đi và được máy thu giải điều chế. Có 2 phát biểu đúng.
}
\end{ex}

% ===== Câu 294 =====
% ID: L12C3B19-71-TN
% Mức: NB
\dangbt{Mô tả sự hình thành và lan truyền của sóng điện từ}
\begin{ex}
% ID: L12C3B19-71-TN
% Mức: NB
Phát biểu nào sau đây đúng về Mô tả sự hình thành và lan truyền của sóng điện từ?
\choice
{Sóng điện từ bắt buộc cần môi trường vật chất để truyền.}
{Có thể bỏ qua phương, chiều và đơn vị trong mọi trường hợp.}
{Kết quả không phụ thuộc dữ kiện và điều kiện áp dụng.}
{\True Sóng điện từ là điện từ trường lan truyền trong không gian và truyền được trong chân không.}
\loigiai{
Sóng điện từ là điện từ trường lan truyền trong không gian và truyền được trong chân không. Vì vậy phương án $D$ đúng; các phương án còn lại sai.
}
\end{ex}

% ===== Câu 295 =====
% ID: L12C3B19-72-TL
% Mức: VDC
\dangbt{Giải thích truyền thông tin bằng sóng điện từ và các tình huống thực tiễn}
\begin{bt}
% ID: L12C3B19-72-TL
% Mức: VDC
So sánh nhận định đúng “Thông tin được ghép lên sóng mang bằng biến điệu, truyền đi và được máy thu giải điều chế.” với nhận định sai “Máy thu nhận thông tin mà không cần chọn sóng hay giải điều chế.”; chỉ rõ căn cứ Vật lí.
\loigiai{
Cần nêu được: Thông tin được ghép lên sóng mang bằng biến điệu, truyền đi và được máy thu giải điều chế. Khi vận dụng phải xác định đúng dữ kiện, phương chiều nếu có, điều kiện áp dụng, hệ đơn vị và kiểm tra tính hợp lí. Nhận định “Máy thu nhận thông tin mà không cần chọn sóng hay giải điều chế.” sai.
}
\end{bt}

% ===== Câu 296 =====
% ID: L12C3B19-72-TLN
% Mức: VDC
\begin{ex}
% ID: L12C3B19-72-TLN
% Mức: VDC
Trong bốn phát biểu về Giải thích truyền thông tin bằng sóng điện từ và các tình huống thực tiễn, có bao nhiêu phát biểu đúng? (Chỉ ghi một số nguyên.) (1) Khi giải cần xác định đúng dữ kiện, phương chiều, điều kiện áp dụng và đơn vị. (2) Máy thu nhận thông tin mà không cần chọn sóng hay giải điều chế. (3) Phải dùng đúng định luật cho tình huống. (4) Thông tin được ghép lên sóng mang bằng biến điệu, truyền đi và được máy thu giải điều chế.
\shortans{3}
\loigiai{
Đối chiếu với kiến thức: Thông tin được ghép lên sóng mang bằng biến điệu, truyền đi và được máy thu giải điều chế. Có 3 phát biểu đúng.
}
\end{ex}

% ===== Câu 297 =====
% ID: L12C3B19-72-TN
% Mức: NB
\dangbt{Mô tả sự hình thành và lan truyền của sóng điện từ}
\begin{ex}
% ID: L12C3B19-72-TN
% Mức: NB
Trong nội dung Mô tả sự hình thành và lan truyền của sóng điện từ?
\choice
{\True Sóng điện từ là điện từ trường lan truyền trong không gian và truyền được trong chân không.}
{Sóng điện từ bắt buộc cần môi trường vật chất để truyền.}
{Có thể bỏ qua phương, chiều và đơn vị trong mọi trường hợp.}
{Kết quả không phụ thuộc dữ kiện và điều kiện áp dụng.}
\loigiai{
Sóng điện từ là điện từ trường lan truyền trong không gian và truyền được trong chân không. Vì vậy phương án $A$ đúng; các phương án còn lại sai.
}
\end{ex}

% ===== Câu 298 =====
% ID: L12C3B19-73-TN
% Mức: NB
\begin{ex}
% ID: L12C3B19-73-TN
% Mức: NB
Tình huống 3: chọn kết luận chính xác về Mô tả sự hình thành và lan truyền của sóng điện từ?
\choice
{Kết quả không phụ thuộc dữ kiện và điều kiện áp dụng.}
{\True Sóng điện từ là điện từ trường lan truyền trong không gian và truyền được trong chân không.}
{Sóng điện từ bắt buộc cần môi trường vật chất để truyền.}
{Có thể bỏ qua phương, chiều và đơn vị trong mọi trường hợp.}
\loigiai{
Sóng điện từ là điện từ trường lan truyền trong không gian và truyền được trong chân không. Vì vậy phương án $B$ đúng; các phương án còn lại sai.
}
\end{ex}

% ===== Câu 299 =====
% ID: L12C3B19-74-TN
% Mức: TH
\begin{ex}
% ID: L12C3B19-74-TN
% Mức: TH
Nội dung nào mô tả đúng nhất Mô tả sự hình thành và lan truyền của sóng điện từ?
\choice
{Có thể bỏ qua phương, chiều và đơn vị trong mọi trường hợp.}
{Kết quả không phụ thuộc dữ kiện và điều kiện áp dụng.}
{\True Sóng điện từ là điện từ trường lan truyền trong không gian và truyền được trong chân không.}
{Sóng điện từ bắt buộc cần môi trường vật chất để truyền.}
\loigiai{
Sóng điện từ là điện từ trường lan truyền trong không gian và truyền được trong chân không. Vì vậy phương án $C$ đúng; các phương án còn lại sai.
}
\end{ex}

% ===== Câu 300 =====
% ID: L12C3B19-75-TN
% Mức: TH
\begin{ex}
% ID: L12C3B19-75-TN
% Mức: TH
Khi phân tích Mô tả sự hình thành và lan truyền của sóng điện từ?
\choice
{Sóng điện từ bắt buộc cần môi trường vật chất để truyền.}
{Có thể bỏ qua phương, chiều và đơn vị trong mọi trường hợp.}
{Kết quả không phụ thuộc dữ kiện và điều kiện áp dụng.}
{\True Sóng điện từ là điện từ trường lan truyền trong không gian và truyền được trong chân không.}
\loigiai{
Sóng điện từ là điện từ trường lan truyền trong không gian và truyền được trong chân không. Vì vậy phương án $D$ đúng; các phương án còn lại sai.
}
\end{ex}

% ===== Câu 301 =====
% ID: L12C3B19-76-TN
% Mức: TH
\begin{ex}
% ID: L12C3B19-76-TN
% Mức: TH
Một học sinh ôn tập Mô tả sự hình thành và lan truyền của sóng điện từ?
\choice
{\True Sóng điện từ là điện từ trường lan truyền trong không gian và truyền được trong chân không.}
{Sóng điện từ bắt buộc cần môi trường vật chất để truyền.}
{Có thể bỏ qua phương, chiều và đơn vị trong mọi trường hợp.}
{Kết quả không phụ thuộc dữ kiện và điều kiện áp dụng.}
\loigiai{
Sóng điện từ là điện từ trường lan truyền trong không gian và truyền được trong chân không. Vì vậy phương án $A$ đúng; các phương án còn lại sai.
}
\end{ex}

% ===== Câu 302 =====
% ID: L12C3B19-77-TN
% Mức: TH
\begin{ex}
% ID: L12C3B19-77-TN
% Mức: TH
Chọn phát biểu không trái với kiến thức về Mô tả sự hình thành và lan truyền của sóng điện từ?
\choice
{Kết quả không phụ thuộc dữ kiện và điều kiện áp dụng.}
{\True Sóng điện từ là điện từ trường lan truyền trong không gian và truyền được trong chân không.}
{Sóng điện từ bắt buộc cần môi trường vật chất để truyền.}
{Có thể bỏ qua phương, chiều và đơn vị trong mọi trường hợp.}
\loigiai{
Sóng điện từ là điện từ trường lan truyền trong không gian và truyền được trong chân không. Vì vậy phương án $B$ đúng; các phương án còn lại sai.
}
\end{ex}

% ===== Câu 303 =====
% ID: L12C3B19-78-TN
% Mức: VD
\begin{ex}
% ID: L12C3B19-78-TN
% Mức: VD
Cơ sở Vật lí đúng dùng để giải Mô tả sự hình thành và lan truyền của sóng điện từ?
\choice
{Có thể bỏ qua phương, chiều và đơn vị trong mọi trường hợp.}
{Kết quả không phụ thuộc dữ kiện và điều kiện áp dụng.}
{\True Sóng điện từ là điện từ trường lan truyền trong không gian và truyền được trong chân không.}
{Sóng điện từ bắt buộc cần môi trường vật chất để truyền.}
\loigiai{
Sóng điện từ là điện từ trường lan truyền trong không gian và truyền được trong chân không. Vì vậy phương án $C$ đúng; các phương án còn lại sai.
}
\end{ex}

% ===== Câu 304 =====
% ID: L12C3B19-79-TN
% Mức: VD
\begin{ex}
% ID: L12C3B19-79-TN
% Mức: VD
Trong một bài thực hành về Mô tả sự hình thành và lan truyền của sóng điện từ?
\choice
{Sóng điện từ bắt buộc cần môi trường vật chất để truyền.}
{Có thể bỏ qua phương, chiều và đơn vị trong mọi trường hợp.}
{Kết quả không phụ thuộc dữ kiện và điều kiện áp dụng.}
{\True Sóng điện từ là điện từ trường lan truyền trong không gian và truyền được trong chân không.}
\loigiai{
Sóng điện từ là điện từ trường lan truyền trong không gian và truyền được trong chân không. Vì vậy phương án $D$ đúng; các phương án còn lại sai.
}
\end{ex}

% ===== Câu 305 =====
% ID: L12C3B19-80-TN
% Mức: VD
\begin{ex}
% ID: L12C3B19-80-TN
% Mức: VD
Kết luận đúng khi vận dụng Mô tả sự hình thành và lan truyền của sóng điện từ?
\choice
{\True Sóng điện từ là điện từ trường lan truyền trong không gian và truyền được trong chân không.}
{Sóng điện từ bắt buộc cần môi trường vật chất để truyền.}
{Có thể bỏ qua phương, chiều và đơn vị trong mọi trường hợp.}
{Kết quả không phụ thuộc dữ kiện và điều kiện áp dụng.}
\loigiai{
Sóng điện từ là điện từ trường lan truyền trong không gian và truyền được trong chân không. Vì vậy phương án $A$ đúng; các phương án còn lại sai.
}
\end{ex}

% ===== Câu 306 =====
% ID: L12C3B19-81-TN
% Mức: NB
\dangbt{Nhận biết các tính chất cơ bản của sóng điện từ}
\begin{ex}
% ID: L12C3B19-81-TN
% Mức: NB
Phát biểu nào sau đây đúng về Nhận biết các tính chất cơ bản của sóng điện từ?
\choice
{Trong sóng điện từ, $E$ và $B$ song song với phương truyền.}
{Có thể bỏ qua phương, chiều và đơn vị trong mọi trường hợp.}
{Kết quả không phụ thuộc dữ kiện và điều kiện áp dụng.}
{\True Trong chân không, vectơ $E$ và $B$ vuông góc nhau và cùng vuông góc phương truyền sóng.}
\loigiai{
Trong chân không, vectơ $E$ và $B$ vuông góc nhau và cùng vuông góc phương truyền sóng. Vì vậy phương án $D$ đúng; các phương án còn lại sai.
}
\end{ex}

% ===== Câu 307 =====
% ID: L12C3B19-82-TN
% Mức: NB
\begin{ex}
% ID: L12C3B19-82-TN
% Mức: NB
Trong nội dung Nhận biết các tính chất cơ bản của sóng điện từ?
\choice
{\True Trong chân không, vectơ $E$ và $B$ vuông góc nhau và cùng vuông góc phương truyền sóng.}
{Trong sóng điện từ, $E$ và $B$ song song với phương truyền.}
{Có thể bỏ qua phương, chiều và đơn vị trong mọi trường hợp.}
{Kết quả không phụ thuộc dữ kiện và điều kiện áp dụng.}
\loigiai{
Trong chân không, vectơ $E$ và $B$ vuông góc nhau và cùng vuông góc phương truyền sóng. Vì vậy phương án $A$ đúng; các phương án còn lại sai.
}
\end{ex}

% ===== Câu 308 =====
% ID: L12C3B19-83-TN
% Mức: NB
\begin{ex}
% ID: L12C3B19-83-TN
% Mức: NB
Tình huống 3: chọn kết luận chính xác về Nhận biết các tính chất cơ bản của sóng điện từ?
\choice
{Kết quả không phụ thuộc dữ kiện và điều kiện áp dụng.}
{\True Trong chân không, vectơ $E$ và $B$ vuông góc nhau và cùng vuông góc phương truyền sóng.}
{Trong sóng điện từ, $E$ và $B$ song song với phương truyền.}
{Có thể bỏ qua phương, chiều và đơn vị trong mọi trường hợp.}
\loigiai{
Trong chân không, vectơ $E$ và $B$ vuông góc nhau và cùng vuông góc phương truyền sóng. Vì vậy phương án $B$ đúng; các phương án còn lại sai.
}
\end{ex}

% ===== Câu 309 =====
% ID: L12C3B19-84-TN
% Mức: TH
\begin{ex}
% ID: L12C3B19-84-TN
% Mức: TH
Nội dung nào mô tả đúng nhất Nhận biết các tính chất cơ bản của sóng điện từ?
\choice
{Có thể bỏ qua phương, chiều và đơn vị trong mọi trường hợp.}
{Kết quả không phụ thuộc dữ kiện và điều kiện áp dụng.}
{\True Trong chân không, vectơ $E$ và $B$ vuông góc nhau và cùng vuông góc phương truyền sóng.}
{Trong sóng điện từ, $E$ và $B$ song song với phương truyền.}
\loigiai{
Trong chân không, vectơ $E$ và $B$ vuông góc nhau và cùng vuông góc phương truyền sóng. Vì vậy phương án $C$ đúng; các phương án còn lại sai.
}
\end{ex}

% ===== Câu 310 =====
% ID: L12C3B19-85-TN
% Mức: TH
\begin{ex}
% ID: L12C3B19-85-TN
% Mức: TH
Khi phân tích Nhận biết các tính chất cơ bản của sóng điện từ?
\choice
{Trong sóng điện từ, $E$ và $B$ song song với phương truyền.}
{Có thể bỏ qua phương, chiều và đơn vị trong mọi trường hợp.}
{Kết quả không phụ thuộc dữ kiện và điều kiện áp dụng.}
{\True Trong chân không, vectơ $E$ và $B$ vuông góc nhau và cùng vuông góc phương truyền sóng.}
\loigiai{
Trong chân không, vectơ $E$ và $B$ vuông góc nhau và cùng vuông góc phương truyền sóng. Vì vậy phương án $D$ đúng; các phương án còn lại sai.
}
\end{ex}

% ===== Câu 311 =====
% ID: L12C3B19-86-TN
% Mức: TH
\begin{ex}
% ID: L12C3B19-86-TN
% Mức: TH
Một học sinh ôn tập Nhận biết các tính chất cơ bản của sóng điện từ?
\choice
{\True Trong chân không, vectơ $E$ và $B$ vuông góc nhau và cùng vuông góc phương truyền sóng.}
{Trong sóng điện từ, $E$ và $B$ song song với phương truyền.}
{Có thể bỏ qua phương, chiều và đơn vị trong mọi trường hợp.}
{Kết quả không phụ thuộc dữ kiện và điều kiện áp dụng.}
\loigiai{
Trong chân không, vectơ $E$ và $B$ vuông góc nhau và cùng vuông góc phương truyền sóng. Vì vậy phương án $A$ đúng; các phương án còn lại sai.
}
\end{ex}

% ===== Câu 312 =====
% ID: L12C3B19-87-TN
% Mức: TH
\begin{ex}
% ID: L12C3B19-87-TN
% Mức: TH
Chọn phát biểu không trái với kiến thức về Nhận biết các tính chất cơ bản của sóng điện từ?
\choice
{Kết quả không phụ thuộc dữ kiện và điều kiện áp dụng.}
{\True Trong chân không, vectơ $E$ và $B$ vuông góc nhau và cùng vuông góc phương truyền sóng.}
{Trong sóng điện từ, $E$ và $B$ song song với phương truyền.}
{Có thể bỏ qua phương, chiều và đơn vị trong mọi trường hợp.}
\loigiai{
Trong chân không, vectơ $E$ và $B$ vuông góc nhau và cùng vuông góc phương truyền sóng. Vì vậy phương án $B$ đúng; các phương án còn lại sai.
}
\end{ex}

% ===== Câu 313 =====
% ID: L12C3B19-88-TN
% Mức: VD
\begin{ex}
% ID: L12C3B19-88-TN
% Mức: VD
Cơ sở Vật lí đúng dùng để giải Nhận biết các tính chất cơ bản của sóng điện từ?
\choice
{Có thể bỏ qua phương, chiều và đơn vị trong mọi trường hợp.}
{Kết quả không phụ thuộc dữ kiện và điều kiện áp dụng.}
{\True Trong chân không, vectơ $E$ và $B$ vuông góc nhau và cùng vuông góc phương truyền sóng.}
{Trong sóng điện từ, $E$ và $B$ song song với phương truyền.}
\loigiai{
Trong chân không, vectơ $E$ và $B$ vuông góc nhau và cùng vuông góc phương truyền sóng. Vì vậy phương án $C$ đúng; các phương án còn lại sai.
}
\end{ex}

% ===== Câu 314 =====
% ID: L12C3B19-89-TN
% Mức: VD
\begin{ex}
% ID: L12C3B19-89-TN
% Mức: VD
Trong một bài thực hành về Nhận biết các tính chất cơ bản của sóng điện từ?
\choice
{Trong sóng điện từ, $E$ và $B$ song song với phương truyền.}
{Có thể bỏ qua phương, chiều và đơn vị trong mọi trường hợp.}
{Kết quả không phụ thuộc dữ kiện và điều kiện áp dụng.}
{\True Trong chân không, vectơ $E$ và $B$ vuông góc nhau và cùng vuông góc phương truyền sóng.}
\loigiai{
Trong chân không, vectơ $E$ và $B$ vuông góc nhau và cùng vuông góc phương truyền sóng. Vì vậy phương án $D$ đúng; các phương án còn lại sai.
}
\end{ex}

% ===== Câu 315 =====
% ID: L12C3B19-90-TN
% Mức: VD
\begin{ex}
% ID: L12C3B19-90-TN
% Mức: VD
Kết luận đúng khi vận dụng Nhận biết các tính chất cơ bản của sóng điện từ?
\choice
{\True Trong chân không, vectơ $E$ và $B$ vuông góc nhau và cùng vuông góc phương truyền sóng.}
{Trong sóng điện từ, $E$ và $B$ song song với phương truyền.}
{Có thể bỏ qua phương, chiều và đơn vị trong mọi trường hợp.}
{Kết quả không phụ thuộc dữ kiện và điều kiện áp dụng.}
\loigiai{
Trong chân không, vectơ $E$ và $B$ vuông góc nhau và cùng vuông góc phương truyền sóng. Vì vậy phương án $A$ đúng; các phương án còn lại sai.
}
\end{ex}

% ===== Câu 316 =====
% ID: L12C3B19-91-TN
% Mức: NB
\dangbt{Tính bước sóng, tần số và tốc độ truyền sóng điện từ}
\begin{ex}
% ID: L12C3B19-91-TN
% Mức: NB
Phát biểu nào sau đây đúng về Tính bước sóng, tần số và tốc độ truyền sóng điện từ?
\choice
{Bước sóng trong chân không bằng f/c.}
{Có thể bỏ qua phương, chiều và đơn vị trong mọi trường hợp.}
{Kết quả không phụ thuộc dữ kiện và điều kiện áp dụng.}
{\True Trong chân không, bước sóng lambda = c/f với c xấp xỉ 3.10^ $8\,\text{m}$ /s.}
\loigiai{
Trong chân không, bước sóng lambda = c/f với c xấp xỉ 3.10^ $8\,\text{m}$ /s. Vì vậy phương án $D$ đúng; các phương án còn lại sai.
}
\end{ex}

% ===== Câu 317 =====
% ID: L12C3B19-92-TN
% Mức: NB
\begin{ex}
% ID: L12C3B19-92-TN
% Mức: NB
Trong nội dung Tính bước sóng, tần số và tốc độ truyền sóng điện từ?
\choice
{\True Trong chân không, bước sóng lambda = c/f với c xấp xỉ 3.10^ $8\,\text{m}$ /s.}
{Bước sóng trong chân không bằng f/c.}
{Có thể bỏ qua phương, chiều và đơn vị trong mọi trường hợp.}
{Kết quả không phụ thuộc dữ kiện và điều kiện áp dụng.}
\loigiai{
Trong chân không, bước sóng lambda = c/f với c xấp xỉ 3.10^ $8\,\text{m}$ /s. Vì vậy phương án $A$ đúng; các phương án còn lại sai.
}
\end{ex}

% ===== Câu 318 =====
% ID: L12C3B19-93-TN
% Mức: NB
\begin{ex}
% ID: L12C3B19-93-TN
% Mức: NB
Tình huống 3: chọn kết luận chính xác về Tính bước sóng, tần số và tốc độ truyền sóng điện từ?
\choice
{Kết quả không phụ thuộc dữ kiện và điều kiện áp dụng.}
{\True Trong chân không, bước sóng lambda = c/f với c xấp xỉ 3.10^ $8\,\text{m}$ /s.}
{Bước sóng trong chân không bằng f/c.}
{Có thể bỏ qua phương, chiều và đơn vị trong mọi trường hợp.}
\loigiai{
Trong chân không, bước sóng lambda = c/f với c xấp xỉ 3.10^ $8\,\text{m}$ /s. Vì vậy phương án $B$ đúng; các phương án còn lại sai.
}
\end{ex}

% ===== Câu 319 =====
% ID: L12C3B19-94-TN
% Mức: TH
\begin{ex}
% ID: L12C3B19-94-TN
% Mức: TH
Nội dung nào mô tả đúng nhất Tính bước sóng, tần số và tốc độ truyền sóng điện từ?
\choice
{Có thể bỏ qua phương, chiều và đơn vị trong mọi trường hợp.}
{Kết quả không phụ thuộc dữ kiện và điều kiện áp dụng.}
{\True Trong chân không, bước sóng lambda = c/f với c xấp xỉ 3.10^ $8\,\text{m}$ /s.}
{Bước sóng trong chân không bằng f/c.}
\loigiai{
Trong chân không, bước sóng lambda = c/f với c xấp xỉ 3.10^ $8\,\text{m}$ /s. Vì vậy phương án $C$ đúng; các phương án còn lại sai.
}
\end{ex}

% ===== Câu 320 =====
% ID: L12C3B19-95-TN
% Mức: TH
\begin{ex}
% ID: L12C3B19-95-TN
% Mức: TH
Khi phân tích Tính bước sóng, tần số và tốc độ truyền sóng điện từ?
\choice
{Bước sóng trong chân không bằng f/c.}
{Có thể bỏ qua phương, chiều và đơn vị trong mọi trường hợp.}
{Kết quả không phụ thuộc dữ kiện và điều kiện áp dụng.}
{\True Trong chân không, bước sóng lambda = c/f với c xấp xỉ 3.10^ $8\,\text{m}$ /s.}
\loigiai{
Trong chân không, bước sóng lambda = c/f với c xấp xỉ 3.10^ $8\,\text{m}$ /s. Vì vậy phương án $D$ đúng; các phương án còn lại sai.
}
\end{ex}

% ===== Câu 321 =====
% ID: L12C3B19-96-TN
% Mức: TH
\begin{ex}
% ID: L12C3B19-96-TN
% Mức: TH
Một học sinh ôn tập Tính bước sóng, tần số và tốc độ truyền sóng điện từ?
\choice
{\True Trong chân không, bước sóng lambda = c/f với c xấp xỉ 3.10^ $8\,\text{m}$ /s.}
{Bước sóng trong chân không bằng f/c.}
{Có thể bỏ qua phương, chiều và đơn vị trong mọi trường hợp.}
{Kết quả không phụ thuộc dữ kiện và điều kiện áp dụng.}
\loigiai{
Trong chân không, bước sóng lambda = c/f với c xấp xỉ 3.10^ $8\,\text{m}$ /s. Vì vậy phương án $A$ đúng; các phương án còn lại sai.
}
\end{ex}

% ===== Câu 322 =====
% ID: L12C3B19-97-TN
% Mức: TH
\begin{ex}
% ID: L12C3B19-97-TN
% Mức: TH
Chọn phát biểu không trái với kiến thức về Tính bước sóng, tần số và tốc độ truyền sóng điện từ?
\choice
{Kết quả không phụ thuộc dữ kiện và điều kiện áp dụng.}
{\True Trong chân không, bước sóng lambda = c/f với c xấp xỉ 3.10^ $8\,\text{m}$ /s.}
{Bước sóng trong chân không bằng f/c.}
{Có thể bỏ qua phương, chiều và đơn vị trong mọi trường hợp.}
\loigiai{
Trong chân không, bước sóng lambda = c/f với c xấp xỉ 3.10^ $8\,\text{m}$ /s. Vì vậy phương án $B$ đúng; các phương án còn lại sai.
}
\end{ex}

% ===== Câu 323 =====
% ID: L12C3B19-98-TN
% Mức: VD
\begin{ex}
% ID: L12C3B19-98-TN
% Mức: VD
Cơ sở Vật lí đúng dùng để giải Tính bước sóng, tần số và tốc độ truyền sóng điện từ?
\choice
{Có thể bỏ qua phương, chiều và đơn vị trong mọi trường hợp.}
{Kết quả không phụ thuộc dữ kiện và điều kiện áp dụng.}
{\True Trong chân không, bước sóng lambda = c/f với c xấp xỉ 3.10^ $8\,\text{m}$ /s.}
{Bước sóng trong chân không bằng f/c.}
\loigiai{
Trong chân không, bước sóng lambda = c/f với c xấp xỉ 3.10^ $8\,\text{m}$ /s. Vì vậy phương án $C$ đúng; các phương án còn lại sai.
}
\end{ex}

% ===== Câu 324 =====
% ID: L12C3B19-99-TN
% Mức: VD
\begin{ex}
% ID: L12C3B19-99-TN
% Mức: VD
Trong một bài thực hành về Tính bước sóng, tần số và tốc độ truyền sóng điện từ?
\choice
{Bước sóng trong chân không bằng f/c.}
{Có thể bỏ qua phương, chiều và đơn vị trong mọi trường hợp.}
{Kết quả không phụ thuộc dữ kiện và điều kiện áp dụng.}
{\True Trong chân không, bước sóng lambda = c/f với c xấp xỉ 3.10^ $8\,\text{m}$ /s.}
\loigiai{
Trong chân không, bước sóng lambda = c/f với c xấp xỉ 3.10^ $8\,\text{m}$ /s. Vì vậy phương án $D$ đúng; các phương án còn lại sai.
}
\end{ex}
